\documentclass[11pt,letterpaper]{article}
\usepackage[T1]{fontenc}
\usepackage[utf8]{inputenc}
\usepackage{lmodern}
\usepackage[margin=1in,headheight=15pt]{geometry}
\usepackage{amsmath,amssymb,amsthm,mathtools,mathrsfs}
\usepackage{microtype,booktabs,longtable,array,enumitem}
\usepackage[dvipsnames]{xcolor}
\usepackage{fancyhdr}
\usepackage[colorlinks=true,linkcolor=MidnightBlue,citecolor=MidnightBlue,urlcolor=MidnightBlue,breaklinks=true]{hyperref}
\usepackage[nameinlink,noabbrev]{cleveref}

\hypersetup{%
  pdftitle={Finite Hahn Deformations: Division, Conservation of Multiplicity, and Residue Duality},
  pdfauthor={Merged research report},
  pdfsubject={Support-controlled deformations of isolated complete intersections over Hahn coefficient rings}}

\pagestyle{fancy}\fancyhf{}
\fancyhead[L]{\small Finite Hahn deformations}
\fancyhead[R]{\small Division, multiplicity, residue duality}
\fancyfoot[C]{\thepage}

\setlength{\parskip}{0.3em}
\setlength{\parindent}{1.25em}
\setlist{nosep,leftmargin=2em}

\newtheorem{theorem}{Theorem}[section]
\newtheorem{lemma}[theorem]{Lemma}
\newtheorem{proposition}[theorem]{Proposition}
\newtheorem{corollary}[theorem]{Corollary}
\theoremstyle{definition}
\newtheorem{definition}[theorem]{Definition}
\newtheorem{hypothesis}[theorem]{Hypothesis}
\newtheorem{example}[theorem]{Example}
\newtheorem{construction}[theorem]{Construction}
\theoremstyle{remark}
\newtheorem{remark}[theorem]{Remark}

\newcommand{\No}{\mathbf{No}}
\newcommand{\SC}{\mathbf{SC}}
\newcommand{\C}{\mathbb C}
\newcommand{\R}{\mathbb R}
\newcommand{\N}{\mathbb N}
\newcommand{\Q}{\mathbb Q}
\newcommand{\Z}{\mathbb Z}
\newcommand{\OO}{\mathcal O}
\newcommand{\HH}{\mathscr H}
\newcommand{\Hp}{\mathscr H^{+}}
\newcommand{\Ip}{\mathscr I}
\newcommand{\Af}{\mathscr F}
\newcommand{\Alg}{\mathcal A}
\newcommand{\mm}{\mathfrak m}
\newcommand{\Kc}{K^{\circ}}
\newcommand{\sh}{^{\sharp}}
\newcommand{\dd}{\mathrm d}
\newcommand{\bx}{\mathcal B}
\newcommand{\Gram}{\mathfrak G}
\newcommand{\TGram}{\mathfrak T}
\newcommand{\Eul}{\mathfrak e}
\newcommand{\Hoint}{\mathop{\oint^{\,H}}\nolimits}
\newcommand{\HInt}{\mathop{\int^{\,H}}\nolimits}
\DeclareMathOperator{\supp}{supp}
\DeclareMathOperator{\st}{st}
\DeclareMathOperator{\NF}{NF}
\DeclareMathOperator{\Res}{Res}
\DeclareMathOperator{\Tr}{Tr}
\DeclareMathOperator{\im}{im}
\DeclareMathOperator{\Mat}{Mat}
\DeclareMathOperator{\Hom}{Hom}
\DeclareMathOperator{\Spec}{Spec}
\DeclareMathOperator{\ord}{ord}
\DeclareMathOperator{\disc}{disc}
\newcolumntype{L}[1]{>{\raggedright\arraybackslash}p{#1}}
\numberwithin{equation}{section}

\title{\bfseries Finite Hahn Deformations\\[0.35em]
\Large Division, Conservation of Multiplicity, and Residue Duality\\[0.55em]
\large Support-controlled deformations of isolated complete intersections\\
over Hahn coefficient rings}
\author{Merged research report}
\date{September 21, 2026}

\begin{document}
\maketitle

\begin{abstract}
Let $f=(f_1,\ldots,f_n)$ be an ordinary holomorphic complete intersection on a fixed polydisk $D\subset\C^n$ whose only common zero in $D$ is the origin, and let $\mu$ be its ordinary colength. We prove that every Hahn perturbation $F=f+E$ whose error support is well ordered and \emph{strictly positive} carries the entire finite geometry of $f$ intact, on the unshrunk domain $D$ and with an explicit support certificate.

Over the nonnegative-support fixed-polydisk ring $\Hp_\Gamma(D)$ the quotient by $(F_1,\ldots,F_n)$ is finite free of rank $\mu$ over the Hahn valuation ring $\Kc$; over $\HH_\Gamma(D)=\OO(D)((t^\Gamma))$ it is finite dimensional of dimension $\mu$ over $K$, on the unchanged ordinary basis. Every $G$ admits a division identity $G=\sum_a c_a e_a+\sum_j F_jQ_j$ with $\supp c_a,\supp Q_j\subseteq\supp G+M$, where $M$ is the additive monoid generated by the error support; the remainder is unique and the quotients need not be. The characters of the quotient are exactly evaluation at the actual displaced surcomplex zeros, and the local Artinian factors are exactly the formal local intersection algebras \emph{at those zeros}, not at their standard parts; hence $\sum_\zeta\mu_\zeta=\mu$, and the same count holds over every infinitesimal target. A coefficientwise multidimensional residue $\Res_F$ gives a perfect pairing over $\Kc$ which survives collisions, and satisfies the trace--Jacobian identity $\Res_F(GJ_F)=\Tr(M_G)=\sum_\zeta\mu_\zeta G\sh(\zeta)$. When all zeros are simple, $\Res_F(G)=\sum_\zeta G\sh(\zeta)/J_F\sh(\zeta)$; for a nonnegative-support numerator individual summands may be infinite while the total is finite.

Four further blocks are proved here. An automatic Hahn--Hartogs theorem: separate coherence in each variable, with the other variables fixed at arbitrary finite surcomplex values, implies joint coherence, the global well-ordered support being a conclusion rather than a hypothesis. A formal-coefficient variant over $\Af_{n,\Gamma}=\C[[z]]((t^\Gamma))$, with finite fibres over every infinitesimal target and, at $\mu=1$, a fine-analytic bijection of monads with fine-analytic inverse. A precision theorem and a discriminant threshold that certify stability without labelling roots, together with a sharp conditioned root-stability theorem whose threshold $\sigma>2\kappa$ and whose two valuation bounds are all proved attained. And a support-preserving lift of the ordinary root cover with unchanged monodromy.

Every theorem names its coefficient ring. The engine throughout is the positive-support Neumann inversion of the Hahn--Neumann support lemma; no Hahn construction uses topological convergence of its partial sums. Ordinary analytic convergence is used only for the classical finite-parameter comparisons before Hahn substitution.
\end{abstract}

\noindent\textbf{Keywords.} Surcomplex numbers; Hahn series; strong summability; isolated complete intersections; finite free deformation; Weierstrass and matrix division; homological perturbation; Grothendieck residues; B\'ezoutian; Hartogs coherence.

\clearpage
{\small\setlength{\parskip}{0pt}\tableofcontents}
\clearpage

%======================================================================
\section{Introduction}
\label{finite:sec:intro}

\subsection{The question}
\label{finite:sec:question}

An isolated zero of an ordinary holomorphic system of $n$ equations in $n$ variables carries a finite local intersection multiplicity $\mu$. Suppose each equation is perturbed by an infinitesimal of the surreal --- more precisely surcomplex --- kind: a Hahn family
\[
 E_j=\sum_{\gamma\in\Gamma_{>0}}e_{j,\gamma}(z)\,t^\gamma,
 \qquad t^\gamma=\omega^{-\gamma},
\]
whose exponent set is well ordered and strictly positive and whose coefficients $e_{j,\gamma}$ are ordinary holomorphic functions. The perturbation need not be polynomial, need not be governed by a single infinitesimal parameter, need not have finitely many scales, and its value group need not have rank one: exponents may include $1$, $\omega$ and $\omega^2$ simultaneously, so that one correction is smaller than every power of another.

The question answered here is whether the whole finite local geometry survives such a perturbation: not merely the existence of nearby roots, but the finite algebra that carries them, its multiplication, its multiplicities, its residue duality, and its trace formula --- and whether all of this can be constructed on \emph{one fixed ordinary coefficient domain}, with an explicitly stated support certificate rather than a sequence of concealed local choices.

The answer is affirmative, and the article proves it in the following delimited form. The leading system is an arbitrary isolated ordinary complete intersection; it need not be monomial, triangular, or have a simple zero. The perturbation is arbitrary of strictly positive well-ordered support. The output lives on the same polydisk that the input did.

\subsection{The central package}
\label{finite:sec:package}

Fix a set-sized divisible ordered subgroup $\Gamma\subseteq(\No,+)$, put
\[
 K=\C((t^\Gamma)),\qquad
 \Kc=\{a\in K:v(a)\geq0\},\qquad
 \mm_K=\{a\in K:v(a)>0\},
\]
and let $\HH_\Gamma(D)=\OO(D)((t^\Gamma))$ be the ring of Hahn families whose coefficient functions are all holomorphic on one fixed polydisk $D$. Let $f_1,\ldots,f_n\in\OO(D)$ have $\{0\}$ as their only common zero in $D$, let $\mu$ be the colength of the resulting ordinary local algebra, and let $e_1=1,e_2,\ldots,e_\mu$ be polynomial representatives of a basis of it. Let $F_j=f_j+E_j$ with $\supp E_j\subseteq\Gamma_{>0}$, let $E=\bigcup_j\supp E_j$ and let $M=E^*$ be the additive monoid it generates, including $0$.

The package proved in \cref{finite:sec:division,finite:sec:points,finite:sec:residues,finite:sec:trace} is:

\begin{enumerate}
\item \textbf{Division with the $S+M$ certificate.} Every $G\in\HH_\Gamma(D)$ satisfies
\begin{equation}\label{finite:eq:introdivision}
 G=\sum_{a=1}^{\mu}c_a(G)\,e_a+\sum_{j=1}^{n}F_jQ_j,
 \qquad
 \supp c_a(G),\ \supp Q_j\subseteq\supp G+M,
\end{equation}
with the scalar vector $c(G)\in K^\mu$ \emph{unique} and the witnesses $Q_j$ explicitly \emph{not} unique. All $Q_j$ have coefficient functions holomorphic on the original $D$.
\item \textbf{Integral finite freeness.} With $\Alg^+=\Hp_\Gamma(D)/(F)$ and $\Alg=\HH_\Gamma(D)/(F)$,
\begin{equation}\label{finite:eq:introfree}
 \Alg^+\simeq(\Kc)^\mu,\qquad \Alg\simeq K^\mu,
\end{equation}
on the unchanged classes of $e_1,\ldots,e_\mu$, and $\Alg^+/\mm_K\Alg^+$ is the ordinary local algebra. This is proved over a valuation ring that may have high-rank value group and need not be Noetherian, and it is not obtained by completing anything.
\item \textbf{Characters are actual evaluations; conservation of multiplicity.} The characters $\Alg\to K$ are exactly the evaluations at the actual zeros $\zeta$ of $F\sh$ in $\mm_K^n$, all zeros in the full surcomplex monad already lie in $\mm_K^n$, and
\begin{equation}\label{finite:eq:introlength}
 \sum_{F\sh(\zeta)=0}\mu_\zeta=\mu,
 \qquad
 \mu_\zeta=\dim_K K[[u_1,\ldots,u_n]]\big/
 \bigl(\widehat F_{1,\zeta},\ldots,\widehat F_{n,\zeta}\bigr),
\end{equation}
where $\widehat F_{j,\zeta}(u)$ is the Taylor expansion of $F_j$ at the \emph{actual displaced zero} $\zeta$ and emphatically not at its standard part. The same count holds for $F-\eta$ at every infinitesimal target $\eta$.
\item \textbf{Perfect residue duality and the trace--Jacobian identity.} A coefficientwise residue $\Res_F$ makes $\Alg^+$ a finite free Frobenius $\Kc$-algebra, with Gram matrix and inverse supported in $M$, and
\begin{equation}\label{finite:eq:introtrace}
 \Res_F(G\,J_F)=\Tr(M_G)=\sum_{F\sh(\zeta)=0}\mu_\zeta\,G\sh(\zeta),
 \qquad
 J_F=\det\Bigl(\frac{\partial F_i}{\partial z_j}\Bigr),
\end{equation}
so that $\Res_F(J_F)=\mu$ exactly, with no positive-support correction; and at simple zeros $\Res_F(G)=\sum_\zeta G\sh(\zeta)/J_F\sh(\zeta)$.
\end{enumerate}

Two features recur and should be stated at once. First, individual local residues $G\sh(\zeta)/J_F\sh(\zeta)$ are routinely \emph{infinite} surcomplex numbers while their total is finite; the cancellation is an exact identity in $K$, not a regularization. Second, the residue Gram determinant is a unit even when the discriminant vanishes: a collision of roots destroys the decomposition into point evaluations but cannot degenerate the pairing.

\subsection{The five archives merged here, and what each contributed}
\label{finite:sec:provenance}

This report merges five manuscripts which prove the package of \cref{finite:sec:package} with one proof strategy, differing in coefficient ring, in the generality of the leading system, and in the single extra block each appends. The shared core is between sixty and seventy per cent of each manuscript, and is printed here once.

\begin{center}\small
\begin{longtable}{@{}L{0.09\textwidth}L{0.30\textwidth}L{0.55\textwidth}@{}}
\toprule
Source & Coefficient ring and leading system & Contribution retained here beyond the shared core\\
\midrule
\endhead
\textbf{08} &
$\HH_\Gamma(D)$, arbitrary isolated complete intersection &
\emph{Base.} The integral structure \eqref{finite:eq:introfree} over $\Kc$ and perfect duality over $\Kc$ (\cref{finite:thm:finitefree,finite:thm:duality}); the Hahn contraction formula as a standalone theorem (\cref{finite:thm:hpl}); parameter preparation and division on a fixed domain with exactly one shrink (\cref{finite:thm:parameterprep}); the finite-witness transfer principle (\cref{finite:lem:witness}); the parameter-linear tensor contraction for perturbed monomial systems (\cref{finite:thm:family}); relative residue duality over a parameter polydisk (\cref{finite:cor:relative}); change of generators (\cref{finite:cor:transformation}).\\[0.35em]
\textbf{05} &
$\HH_\Gamma(D)$, arbitrary isolated complete intersection &
Koszul/homological-perturbation transport of \emph{every} relation with explicit preimage (\cref{finite:cor:koszul}); the precision theorem (\cref{finite:thm:precision}); the discriminant--Jacobian identity (\cref{finite:thm:discriminant}); the discriminant threshold for preserving simple zeros (\cref{finite:cor:discstable}); the finite-parameter comparison principle (\cref{finite:lem:finitedependence}); the coefficient-extraction residue formula (\cref{finite:cor:coefficientresidue}); the rank-three quartic over $1,\omega,\omega^2$ (\cref{finite:sec:quartic}).\\[0.35em]
\textbf{06} &
$\Af_{n,\Gamma}=\C[[z]]((t^\Gamma))$, arbitrary isolated formal complete intersection &
The formal-coefficient variant and the warning example that fixes the correct algebra (\cref{finite:ex:unit}); the fibre theorems --- statements about the \emph{map}, not the zero set (\cref{finite:thm:fibres}) --- and the monad bijection with fine-analytic inverse at $\mu=1$ (\cref{finite:cor:inverse}); the monad-by-monad conservation theorem on an ordinary base (\cref{finite:thm:monadcount}); exact cluster moments and separating linear forms (\cref{finite:prop:separation}); the valuation localization bound (\cref{finite:thm:localization}); the triple-collision example at two separated scales (\cref{finite:sec:triple}).\\[0.35em]
\textbf{11} &
$\HH_\Gamma(V\times D)$, coordinate-power leading system &
Sharp conditioned root stability with \emph{both} bounds proved attained and the threshold proved strict (\cref{finite:thm:rootstability,finite:prop:sharpness}); the matching corollary (\cref{finite:cor:matching}); the ordinary holomorphic parameter domain $V$; the summable matrix functional calculus (\cref{finite:lem:matrix}); the B\'ezoutian reproducing kernel and the resulting identification of the trace element with the Jacobian (\cref{finite:thm:bezout}); localized residues at actual surcomplex points (\cref{finite:thm:localres}); coordinate-change invariance (\cref{finite:thm:transform}).\\[0.35em]
\textbf{15} &
$\HH_\Gamma(U\times D)$, coordinate-power leading system &
The automatic Hahn--Hartogs coherence theorem, with the global support a conclusion (\cref{finite:lem:autosupport,finite:thm:hartogs,finite:cor:hole}); support-controlled \emph{matrix} division, which is what controls the coupling in an induction over the equations (\cref{finite:thm:matrixdivision}); the support-preserving lift of the ordinary root cover with unchanged monodromy (\cref{finite:thm:monodromy}); the one-variable collision-stable trace and residue formulas (\cref{finite:prop:univtrace}); the cubic example separating residue duality from trace degeneration (\cref{finite:sec:cubic}).\\
\bottomrule
\end{longtable}
\end{center}

Two pairs among the five are near-duplicates rather than variants. Manuscripts 05 and 08 prove the same theorem over the same ring for the same class of leading systems, differing in proof route and in the coda; 11 and 15 both treat the parametrized coordinate-power specialization with identical division, conservation and duality statements, their genuinely separate contributions being one section each. Manuscript 06 proves the core a third time over formal coefficients; its independent value is the fibre and monad material, not the core.

\subsection{What is classical, and what is proposed}
\label{finite:sec:classical}

Hahn support algebra and the algebraic closedness of Hahn fields with divisible value group and algebraically closed residue field are established \cite{finite:Neumann,finite:Poonen}. Analysis over surreal fields, restricted analytic structures and surreal exponentiation precede this work \cite{finite:Alling,finite:vdDE,finite:BM}. The contraction formulas are a specialization of the homological perturbation lemma and are not a new algebraic identity \cite{finite:Crainic}. Ordinary local residues, finite holomorphic maps, complete-intersection duality, and the normal-form/residue/trace correspondence are classical \cite{finite:Cartan,finite:GunningRossi,finite:GriffithsHarris,finite:Hartshorne,finite:Kunz,finite:Demailly,finite:CDS,finite:CattaniDickenstein,finite:Cox,finite:NayakSastry}. Weierstrass-type machinery over valued fields is developed by Cluckers and Lipshitz \cite{finite:CluckersLipshitz}. The ordinary Hartogs theorems are classical \cite{finite:Lebl,finite:BK}.

What is offered as new is the precise combination: arbitrary well-ordered positive supports, fixed-domain division witnesses, finite freeness over a possibly non-Noetherian valuation ring, identification of all actual surcomplex roots and of their formal local lengths at the displaced points, compatibility of a coefficientwise contour functional with that finite algebra, and the automatic-support Hartogs theorem. \Cref{finite:sec:status} records the consolidated and deliberately limited status assessment of all five sources.

%======================================================================
\section{Conventions, and the four coefficient rings}
\label{finite:sec:conventions}

\subsection{One symbol per concept}
\label{finite:sec:notation}

The five sources use incompatible notation for identical objects. The following choices are fixed once and used without exception; synonyms are recorded here and nowhere else.

\begin{itemize}
\item $\No$ is the surreal class, $\SC=\No[i]$ the surcomplex class. $\Gamma$ is a nonzero \emph{set-sized divisible} ordered subgroup of $\No$, and $t^\gamma=\omega^{-\gamma}$, so that $K=\C((t^\Gamma))$ sits inside $\SC$ by Conway normal forms.
\item $v$ is the Hahn valuation, $v(a)=\min\supp a$ for $a\neq0$ and $v(0)=+\infty$. For a coherent datum $G$ with coefficient functions, $v(G)$ is the least exponent carrying a coefficient function that is not identically zero; it is not a value at a point. For a vector or matrix, $v$ is the minimum of the entry valuations.
\item $\Kc=\{v\geq0\}$ is the valuation ring, the \emph{finite} elements of $K$; $\mm_K=\{v>0\}$ its maximal ideal, the \emph{infinitesimals}. Sources write $R$, $\OO_K$, $K^+$ and $K^\circ$ for $\Kc$; only $\Kc$ is used here.
\item $\st:\Kc\to\C$ is the standard part, equivalently the coefficient at exponent zero; it is not defined on infinite elements of $K$. On a finite thickening $U\sh_K$, the \emph{standard-part topology} is the pullback of the ordinary topology of $U$ under $\st$. Every monad has the indiscrete subspace topology for this choice; it is different from the fine topology. Three sources call this the residue topology, the reduction topology and the shadow topology. Only ``standard part'' and $\st$ are used below.
\item $\sharp$ denotes coherent evaluation on an infinitesimal thickening: $G\sh(c+\varepsilon)$.
\item $E=\bigcup_j\supp E_j$ is the perturbation support and $M=E^*$ the additive monoid it generates, including $0$; so $M+M=M$. Sources write $T^*$, $M$, $\mathcal E^*$; only $M$ is used.
\item $\mu$ is the leading colength, $\mu_\zeta$ the local multiplicity at an actual zero $\zeta$. In the coordinate-power case $\mu=\prod_i m_i$; the sources' $N$ is this $\mu$.
\item $e_1=1,e_2,\ldots,e_\mu$ are the fixed polynomial representatives of the leading basis; $\varepsilon_1,\ldots,\varepsilon_n$ are the Koszul exterior generators. Sources overload a single letter for both; here they never collide.
\item $\NF_F$ is the normal form, $c_a(G)$ its coordinates, $M_G$ multiplication by $[G]$ in the fixed basis, $J_F$ the Jacobian determinant, $\Res_F$ the residue functional, $\Gram_F$ its Gram matrix and $\TGram_F$ the trace Gram matrix.
\item $(i,p,h)$ is a contraction, $\delta$ the perturbation differential, $S=(1+\delta h)^{-1}$, and $p_F=pS$, $h_F=hS$.
\end{itemize}

One notational clash deserves its own sentence, because two sources use the single symbol $\Delta_F$ for two different quantities. Here
\begin{equation}\label{finite:eq:twodiscriminants}
 \Delta_F:=\det M_{J_F}
 \quad\text{(basis independent)},
 \qquad
 \disc_F:=\det\TGram_F
 \quad\text{(basis relative)} .
\end{equation}
\Cref{finite:thm:discriminant} proves $\disc_F=(\det\Gram_F)\,\Delta_F$ and hence $v(\disc_F)=v(\Delta_F)$, so the two vanish together; they are nevertheless different elements of $K$ and are never written with the same symbol below.

\subsection{Set-sized workspaces}
\label{finite:sec:workspace}

All index sets, supports and recursions are sets. Every finite tuple of surcomplex numbers occurring in a proof lies in some set-sized divisible workspace: take the group generated by the exponents of its Conway normal forms and pass to the divisible hull. Enlarging $\Gamma$ to a larger set-sized divisible subgroup is harmless and is done freely, for instance to accommodate an infinitesimal target or the coordinates of a proposed zero. This device never authorizes a proper-class sum or a polynomial ring over a proper class, and no statement below is made about a proper-class support. $K$ is algebraically closed, by the theorem on Hahn fields with divisible value group and algebraically closed residue field \cite[Corollary 4]{finite:Poonen}.

\subsection{The four coefficient rings, named once}
\label{finite:sec:rings}

Four different coefficient rings occur in the literature surrounding this material under nearly identical notation and carrying nearly word-for-word identical theorem statements. This is the single most dangerous place for a false attribution, so the four are named here and every theorem below names the one it is stated over. \emph{No theorem is transferred from one of these rings to another anywhere in this report.}

\begin{enumerate}
\item \textbf{The fixed-polydisk ring.} For an ordinary domain $U\subseteq\C^m$,
\[
 \HH_\Gamma(U)=\OO(U)((t^\Gamma))
 =\Bigl\{\textstyle\sum_{\gamma\in S}g_\gamma t^\gamma:
 S\subseteq\Gamma\ \text{well ordered},\ g_\gamma\in\OO(U)\Bigr\},
\]
with $\Hp_\Gamma(U)$ its nonnegative-support subring and $\Ip_\Gamma(U)$ its strictly-positive-support ideal. Every coefficient function is holomorphic on the \emph{one} fixed domain $U$; no bound on their ordinary sizes is imposed. Coefficient functions that vanish identically are omitted from the support. This is the ring of \cref{finite:sec:division,finite:sec:points,finite:sec:residues,finite:sec:trace,finite:sec:stability}, and it is the ring of sources 05, 08, 11 and 15.
\item \textbf{The formal-coefficient ring.} With $A_n=\C[[z_1,\ldots,z_n]]$,
\[
 \Af_{n,\Gamma}=A_n((t^\Gamma)),
 \qquad
 \Af^+_{n,\Gamma}=\{H\in\Af_{n,\Gamma}:\supp H\subseteq\Gamma_{\geq0}\},
\]
one well-ordered Hahn support shared by all the Taylor coefficients. No ordinary convergence is required of the $g_\gamma$. This is the ring of \cref{finite:sec:formal} and of source 06.
\item \textbf{The radius-free ring} $\C\{z_1,\ldots,z_n\}((t^\Gamma))$, in which each Hahn coefficient is individually a convergent germ and the radii may tend to zero, so that there need be \emph{no} common polydisk. It is the ring of the companion analytic-geometry report and does not occur in any theorem below.
\item \textbf{The common-domain germ ring}, consisting of germs of Hahn families that admit one ordinary neighbourhood serving all coefficients before shrinking. It too belongs to the companion report and occurs in no theorem below.
\end{enumerate}

Rings (1)--(4) are pairwise distinct, and the two containments that matter are strict. The companion report records the two results that relate them: a strictness example exhibiting an element of the radius-free ring with no common ordinary polydisk, together with the warning that invoking a theorem stated only for coefficients holomorphic on one fixed polydisk would leave a gap in the radius-free setting; and a faithful-flatness theorem showing that the inclusion of the radius-free ring into the formal ring has the same maximal ideals, residue fields, completions and finite-colength quotients. Those two results are cited here, not reproved, and they are the only permitted bridge between rings; they are \emph{not} used to move any theorem of this report.

Within this report the only bridge used between rings (1) and (2) is explicit and one-directional: Taylor expansion at an ordinary point $c\in U$,
\begin{equation}\label{finite:eq:taylorbridge}
 \HH_\Gamma(U)\longrightarrow\Af_{n,\Gamma},
 \qquad
 \sum_\gamma g_\gamma t^\gamma\longmapsto
 \sum_\gamma\Bigl(\textstyle\sum_\alpha\frac{\partial^\alpha g_\gamma(c)}{\alpha!}u^\alpha\Bigr)t^\gamma,
\end{equation}
which is a ring homomorphism and does not enlarge the Hahn support. It is used exactly once, in \cref{finite:thm:monadcount}, and its direction is never reversed.

A last remark on the meaning of the phrase \emph{global support control}: here it means a uniform well-ordered support for a family of coefficient functions on one stated ordinary domain. It does not mean an all-scale coherence condition imposed on every surreal affine chart, which is a different and much more restrictive condition.

%======================================================================
\section{The support calculus and its two negative companions}
\label{finite:sec:support}

\subsection{Strong summability}
\label{finite:sec:strong}

\begin{definition}[Strong summability]
\label{finite:def:strong}
A set-indexed family of Hahn series is \emph{strongly summable} if the union of its supports is well ordered \emph{and} each exponent receives a contribution from only finitely many members of the family. Its sum is then formed by finite coefficient addition at each exponent. The same definition applies verbatim to series with coefficients in an arbitrary complex vector space $W$, written $W((t^\Gamma))$, and hence to holomorphic coefficient functions, vectors, matrices and finite exterior powers.
\end{definition}

There is no limiting process in \cref{finite:def:strong}. Every complex-linear map $L:W\to W'$ extends coefficientwise, satisfies $\supp(LG)\subseteq\supp(G)$, and preserves every strongly summable family. Neither continuity of $L$ for the Fr\'echet topology of holomorphic functions nor a bound on the sizes of complex coefficients is needed for this; the point will be used repeatedly.

\begin{lemma}[Hahn--Neumann support lemma]
\label{finite:lem:support}
Let $S,T$ be well-ordered subsets of an ordered abelian group. Then $S+T$ is well ordered and each exponent has only finitely many representations $s+t$. If $E\subseteq\Gamma_{>0}$ is well ordered, then
\[
 E^*=\{e_1+\cdots+e_k:k\geq0,\ e_j\in E\}
\]
is well ordered and a fixed exponent has only finitely many representations as a finite word in $E$, permutations counted; the empty word contributes $0$. Consequently, for well-ordered $S$, the set $S+E^*$ is well ordered and a fixed exponent has only finitely many representations $s+e_1+\cdots+e_k$ with $s\in S$ and $e_i\in E$.
\end{lemma}

\begin{proof}
For $S+T$: an infinite strictly decreasing sequence of sums, or infinitely many pairs with one sum, forces one coordinate to decrease along an infinite subsequence, contradicting well-ordering. The monoid assertion is Neumann's lemma \cite{finite:Neumann}; one proof uses the well-quasi-ordering of finite words over a well-ordered positive alphabet under coordinatewise embedding: an embedding cannot decrease the sum of letters, and equality of sums forces equal words. This excludes both a descending chain of sums and infinitely many words with a common sum. Each finite word has finitely many permutations. Now apply the first assertion to $S$ and $E^*$.
\end{proof}

\Cref{finite:lem:support} is the engine of essentially every proof in this report, and $S+M$ is the universal support certificate. Two negative companions to it are load-bearing and are stated with equal emphasis, because both are separately restated in every one of the five sources.

\begin{remark}[Strong summability is not a valuation limit]
\label{finite:rem:notlimit}
Take $\Gamma=\Q+\Q\omega$. The scalar series $(1-t)^{-1}=\sum_{m\geq0}t^m$ is strongly summable, but the valuations $0,1,2,\ldots$ of its tails never reach $\omega$, so the sequence of finite partial sums does not approach the sum to valuation accuracy $\omega$. Therefore \emph{no} argument of the form ``the valuation improves by a fixed positive amount at each iteration, hence the iteration converges'' justifies any construction in this report. Every infinite formula below is justified by finite contribution at each exponent, never by treating $k\cdot\min E$ as cofinal in $\Gamma$. In the full surreal fine topology the obstruction is stronger still: a set of nonzero positive distances has a smaller positive surreal lower bound, so a convergent set-indexed net is eventually constant.
\end{remark}

\begin{remark}[Finitely many dependencies is not finitely many exponents below]
\label{finite:rem:truncation}
$E^*\cap(-\infty,\eta)$ need not be finite. Already for $E=\{1,\omega\}$ there are infinitely many elements of $E^*$ below $\omega$. The proofs below always select the finitely many words \emph{contributing to one specified coefficient}; they never truncate to an allegedly finite initial segment of the ordered group. The distinction is exactly what makes the finite-parameter replacement arguments of \cref{finite:sec:trace} legitimate, and it is also the obstruction to reading those arguments as an algorithm: see \cref{finite:sec:computation}. Well-ordering cannot be weakened to positivity either: $\{1/r:r\geq1\}$ is positive and not well ordered.
\end{remark}

\subsection{Positive-support operator inversion}
\label{finite:sec:neumann}

\begin{lemma}[Support-controlled operator inversion]
\label{finite:lem:operator}
Let $W$ be a complex vector space, let $E\subseteq\Gamma_{>0}$ be well ordered, and for each $e\in E$ let $L_e:W\to W$ be complex linear. Extend
\[
 \mathcal L=\sum_{e\in E}t^e L_e
\]
coefficientwise to $W((t^\Gamma))$. Then $\mathcal L$ is a well-defined support-increasing operator, and for $X$ of support $S$ the equation $Y+\mathcal LY=X$ has the unique solution
\begin{equation}\label{finite:eq:neumann}
 Y=(1+\mathcal L)^{-1}X=\sum_{k\geq0}(-\mathcal L)^kX,
 \qquad \supp Y\subseteq S+E^*,
\end{equation}
the displayed family being strongly summable; $(1+\mathcal L)^{-1}$ is a two-sided inverse. The same holds for a finite matrix of such operators, and, more generally, whenever every application of an operator $L$ built from coefficientwise complex-linear maps and from multiplication by positive-support Hahn data introduces at least one letter of $E$. If all $L_e$ preserve $\OO(U)$ for a fixed $U$, every coefficient of $Y$ is holomorphic on that same $U$.
\end{lemma}

\begin{proof}
Expand $\mathcal L^kX$ into terms indexed by an exponent $s\in S$ and an ordered word $(e_1,\ldots,e_k)$; the exponent of such a term is $s+e_1+\cdots+e_k$. By \cref{finite:lem:support} the union of these exponents lies in the well-ordered set $S+E^*$, and each exponent receives only finitely many contributions even when the word length is unbounded. Hence the family is strongly summable and each output coefficient is a finite sum of finite compositions of the $L_e$ applied to input coefficients; this also gives the fixed-domain statement, since coefficientwise linear maps introduce no exponents. Multiplying the series by $1+\mathcal L$ telescopes coefficientwise. For uniqueness, if $Z+\mathcal LZ=0$ with $Z\neq0$ then $v(\mathcal LZ)>v(Z)$, so cancellation at $v(Z)$ is impossible. Note that the cancellation argument is legitimate without any statement that the finite partial sums converge in a topology, by \cref{finite:rem:notlimit}.
\end{proof}

We shall also use an elementary nonlinear variant: an equation whose coefficient at exponent $\nu$ is a fixed linear splitting applied to the input coefficient at $\nu$, plus finite products of previously determined \emph{positive} coefficients, is solved by well-founded recursion through $E^*$. Each product uses exponents strictly below $\nu$, well-ordering supplies the recursion, and \cref{finite:lem:support} makes each coefficient formula finite. \Cref{finite:thm:parameterprep} is the explicit instance.

%======================================================================
\section{Evaluation, thickenings, and finite Taylor identities}
\label{finite:sec:evaluation}

\subsection{Coherent evaluation}
\label{finite:sec:coherenteval}

Write $\mm_{\SC}$ for the surcomplex infinitesimals, those elements whose modulus is smaller than every positive ordinary real. For an ordinary domain $U\subseteq\C^m$ put
\[
 U\sh=\{c+\varepsilon:c\in U,\ \varepsilon\in\mm_{\SC}^{\,m}\},
 \qquad
 U\sh_K=\{c+\varepsilon:c\in U,\ \varepsilon\in\mm_K^m\}.
\]

\begin{definition}[Coherent evaluation]
\label{finite:def:eval}
For $G=\sum_\gamma g_\gamma t^\gamma\in\HH_\Gamma(U)$ and $c+\varepsilon\in U\sh$,
\begin{equation}\label{finite:eq:evaluation}
 G\sh(c+\varepsilon)
 =\sum_{\gamma}\sum_{\alpha\in\N^m}
 \frac{\partial^\alpha g_\gamma(c)}{\alpha!}\,t^\gamma\varepsilon^\alpha .
\end{equation}
\end{definition}

\begin{proposition}[Evaluation is a homomorphism]
\label{finite:prop:evaluation}
The family in \eqref{finite:eq:evaluation} is strongly summable, with support contained in $\supp G+A^*$ where $A=\bigcup_j\supp\varepsilon_j$. The map $G\mapsto G\sh$ into functions on $U\sh$ is an injective algebra homomorphism and commutes with every partial derivative in the ordinary variables, leaving the monomials $t^\gamma$ fixed; evaluation at one point need not be injective. If $\supp G\subseteq\Gamma_{\geq0}$ then $\st G\sh(c+\varepsilon)=g_0(c)$. If $G\in\HH_\Gamma(U\times V)$ and $a\in U\sh$, evaluation in the $U$-variables alone yields a coherent datum on the \emph{same} domain $V$, with support in $\supp G+A^*$; no shrinking of $V$ occurs.
\end{proposition}

\begin{proof}
Each term has support in $\supp G+A^*$, and by \cref{finite:lem:support} a fixed exponent admits only finitely many decompositions into an exponent of $G$ and a finite word in $A$; there are finitely many ways to assign such a word to the $m$ coordinates. Hence strong summability. The product and derivative rules reduce coefficientwise to the ordinary multivariable Taylor rules, with the same finite-contribution justification for every regrouping. At an ordinary point the value is $\sum_\gamma g_\gamma(c)t^\gamma$; vanishing for all $c$ forces every coefficient function to vanish, which gives injectivity. In the nonnegative-support case every term other than $g_0(c)$ has positive valuation. Partial evaluation leaves the derivatives in the evaluated variables as ordinary holomorphic functions on $V$, and the support computation is independent of the point of $V$.
\end{proof}

For $G\in\Hp_\Gamma(U)$, reduction means taking the coefficient at exponent zero; at $c+\varepsilon$ the standard part of $G\sh(c+\varepsilon)$ is that coefficient evaluated at $c$. This is what distinguishes zeros lying in one monad from ordinary zeros lying in a different monad. Equalities of coherent data may be checked at all \emph{ordinary} points; in one variable, equality at ordinary points of a set with an ordinary accumulation point suffices, coefficient by coefficient. This identity principle is a statement about coefficient functions and is explicitly distinct from any statement about fine-topological sequences.

\subsection{Division at an infinitesimally displaced centre}
\label{finite:sec:shifted}

The following lemma is what rules out extraneous algebraic characters in \cref{finite:sec:points}. It is stated on the \emph{original} polydisk: no shrinking is involved.

\begin{lemma}[Shifted division and finite Taylor remainder]
\label{finite:lem:shifted}
Let $D\subset\C^n$ be a polydisk centred at $0$, let $\zeta\in\mm_K^n$, put $A=\bigcup_j\supp\zeta_j$, and let $G\in\HH_\Gamma(D)$ have support $S$. There are $Q_j\in\HH_\Gamma(D)$ with
\begin{equation}\label{finite:eq:shifteddivision}
 G(z)=G\sh(\zeta)+\sum_{j=1}^n(z_j-\zeta_j)Q_j(z),
 \qquad \supp Q_j\subseteq S+A^*.
\end{equation}
More generally, for every ordinary integer $N\geq1$ there are $Q_\alpha\in\HH_\Gamma(D)$ with
\begin{equation}\label{finite:eq:finitetaylor}
 G(z)=\sum_{|\alpha|<N}
 \frac{(\partial^\alpha G)\sh(\zeta)}{\alpha!}(z-\zeta)^\alpha
 +\sum_{|\alpha|=N}(z-\zeta)^\alpha Q_\alpha(z),
\end{equation}
with the same support bound. All coefficient functions of all remainders are holomorphic on $D$. Consequently the kernel of evaluation at $\zeta$ is the ideal $(z_1-\zeta_1,\ldots,z_n-\zeta_n)$, and Taylor expansion at $\zeta$ is a ring homomorphism $\HH_\Gamma(D)\to K[[u_1,\ldots,u_n]]$.
\end{lemma}

\begin{proof}
In one coordinate set $R_jg=g|_{z_j=0}$ and $D_jg=(g-R_jg)/z_j$; the apparent singularity is removable, both operators preserve holomorphy on $D$, and $g=R_jg+z_jD_jg$. Treating the other coordinates as parameters, division by $z_j-\zeta_j$ gives
\[
 Q=\sum_{k\geq0}\zeta_j^kD_j^{k+1}G,
 \qquad
 R=R_jG+\zeta_jR_jQ=\sum_{k\geq0}\frac{\zeta_j^k}{k!}(\partial_j^kG)|_{z_j=0},
\]
which is strongly summable with the stated support by \cref{finite:lem:operator}, and whose remainder is independent of $z_j$. Dividing successively in all coordinates gives \eqref{finite:eq:shifteddivision} with constant remainder $G\sh(\zeta)$. Applying the same division to each remainder, $N$ times, produces a polynomial of total degree $<N$ in $z-\zeta$ plus a remainder in the $N$th power of the displaced coordinate ideal; differentiating and evaluating at $\zeta$ identifies the polynomial coefficients as in \eqref{finite:eq:finitetaylor}. Finitely many repetitions do not enlarge the controlling monoid $A^*$. Finally, the finite Taylor identities give the usual finite convolution for the Taylor coefficients of a product at each multi-index, so the Taylor map is a ring homomorphism; no convergence in $u$ is asserted for the resulting formal series.
\end{proof}

%======================================================================
\section{Automatic support and the Hahn--Hartogs theorem}
\label{finite:sec:hartogs}

This section is foundational: it is a statement about the function class itself, logically independent of every deformation theorem below, and it removes a hypothesis rather than adding one. It is stated over the fixed-polydisk ring $\HH_\Gamma(U)$ of \cref{finite:sec:rings}(1).

\subsection{Pointwise support cannot conceal a descending chain}
\label{finite:sec:autosupport}

\begin{lemma}[Automatic common support]
\label{finite:lem:autosupport}
Let $X$ be a connected ordinary complex domain, or more generally a connected complex manifold. Let $(f_\gamma)_{\gamma\in A}$ be a set-indexed family of ordinary holomorphic functions on $X$, indexed by a subset $A$ of an ordered abelian group. Suppose that for every $x\in X$ the set $A_x=\{\gamma\in A:f_\gamma(x)\neq0\}$ is well ordered. Then the exact global support $A'=\{\gamma\in A:f_\gamma\not\equiv0\}$ is well ordered.
\end{lemma}

\begin{proof}
If $A'$ were not well ordered, choose a strictly decreasing sequence $\gamma_1>\gamma_2>\cdots$ in $A'$. For each $k$ the zero set of the nonzero holomorphic function $f_{\gamma_k}$ is closed with empty interior, so its complement $V_k$ is open and dense. A complex manifold is locally compact Hausdorff, hence a Baire space, so $\bigcap_kV_k$ is dense and in particular nonempty. At any of its points all the exponents $\gamma_k$ occur, contradicting well-ordering of $A_x$.
\end{proof}

No bound on the number or on the sizes of the nonzero coefficient functions enters. The force of the lemma is the impossibility of making countably many nonzero holomorphic functions vanish collectively at every point.

\subsection{Separate coherence implies joint coherence}
\label{finite:sec:separate}

\begin{theorem}[Automatic Hahn--Hartogs coherence]
\label{finite:thm:hartogs}
Let $U_1,\ldots,U_d\subseteq\C$ be nonempty connected domains and $U=\prod_jU_j$. Let
\[
 F:U_1\sh\times\cdots\times U_d\sh\longrightarrow\SC
\]
be a class function, and assume that for every coordinate $j$, and every fixed choice of the other coordinates in their \emph{surcomplex} thickenings, the resulting one-variable function is Hahn-coherent on $U_j\sh$. Then there are a set-sized exponent group $\Gamma$ and a unique $H\in\HH_\Gamma(U)$ with $F=H\sh$. No common-support hypothesis on the slices, no ordinary local boundedness hypothesis on $F$, and no fine-topological joint continuity hypothesis is required: the global well-ordered support is a \emph{conclusion}.
\end{theorem}

\begin{proof}
Restrict $F$ to the ordinary set $U$. The union $A=\bigcup_{c\in U}\supp F(c)$ is a set, by replacement and union over the set $U$, each value having set-sized support; let $\Gamma$ be the divisible hull of the group it generates. For $\gamma\in A$ put $f_\gamma(c)=[t^\gamma]F(c)$. Fixing all ordinary coordinates but the $j$th, coherence of that slice makes $f_\gamma$ holomorphic in the remaining coordinate, and $f_\gamma$ is zero there if $\gamma$ is absent from the slice support. So each $f_\gamma$ is separately ordinary holomorphic, and classical Hartogs separate holomorphy \cite{finite:BK} gives $f_\gamma\in\OO(U)$; this classical theorem is an input, not a new surcomplex assertion. At each ordinary $c$ the active support is $\supp F(c)$, well ordered by hypothesis, so \cref{finite:lem:autosupport} makes $S=\{\gamma:f_\gamma\not\equiv0\}$ well ordered. Set $H=\sum_{\gamma\in S}f_\gamma t^\gamma\in\HH_\Gamma(U)$; then $H\sh=F$ at all ordinary tuples.

Extend the equality one coordinate at a time. Suppose it holds when the first $k$ coordinates are arbitrary finite surcomplex points and the rest ordinary. Fix those $k$ points and all ordinary coordinates after $k+1$. In coordinate $k+1$, $F$ is coherent by hypothesis and $H\sh$ is coherent by \cref{finite:prop:evaluation}, after enlarging the workspace if necessary; the two agree at every ordinary value of that coordinate by the induction hypothesis, so coefficientwise identity makes them equal on the full thickening. Induction reaches $k=d$. Uniqueness is coefficientwise identity at ordinary tuples.
\end{proof}

\begin{remark}[The exact role of the slice hypothesis]
\label{finite:rem:slices}
The hypothesis concerns fixed \emph{surcomplex} values of all other coordinates, not merely fixed ordinary values. Ordinary-slice data determine the ordinary restriction and its joint Hahn series, but do not by themselves constrain values assigned arbitrarily at tuples with several nonordinary coordinates. The coordinate-by-coordinate induction uses precisely the stronger hypothesis stated.
\end{remark}

\begin{corollary}[Coherent Hartogs removal]
\label{finite:cor:hole}
Let $d\geq2$, let $\Omega\subseteq\C^d$ be a connected domain, and let $\mathcal K\subset\Omega$ be compact with $\Omega\setminus\mathcal K$ connected. Every element of $\HH_\Gamma(\Omega\setminus\mathcal K)$ extends uniquely to an element of $\HH_\Gamma(\Omega)$, with the same exact support.
\end{corollary}

\begin{proof}
Apply the classical Hartogs extension theorem \cite[Theorem 4.3.1]{finite:Lebl} to each ordinary coefficient function. A nonzero coefficient cannot become zero after extension, since its restriction is the original coefficient, so the support is unchanged. Coefficientwise uniqueness gives uniqueness.
\end{proof}

\Cref{finite:cor:hole} is a direct transfer of a classical theorem, recorded for completeness. The non-formal content is \cref{finite:thm:hartogs}: the automatic existence of a common support \emph{before} any joint datum has been given.

%======================================================================
\section{Fixed-domain preparation and matrix division with parameters}
\label{finite:sec:preparation}

Everything in this section is stated over the fixed-polydisk ring $\HH_\Gamma$ of \cref{finite:sec:rings}(1). The point of the section is that exactly \emph{one} ordinary shrink is ever performed, before the Hahn recursion begins, and none afterwards at any Hahn exponent.

\subsection{A fixed-domain ordinary division pair}
\label{finite:sec:ordinarypair}

Let $B\subset\C^r$ be a polydisk of parameters, let $D_R=\{z\in\C:|z|<R\}$, and let
\[
 P_0(y,z)=z^m+p_{m-1}(y)z^{m-1}+\cdots+p_0(y),
 \qquad p_k\in\OO(B),
\]
have, for \emph{every} $y\in B$, all its $z$-roots in $|z|<r_1$, where $r_1<r_0<R$ are fixed ordinary radii. This uniformity over the whole of $B$ is what makes the operators below fixed-domain. For $h\in\OO(B\times D_R)$ define
\begin{align}
 R_0h(y,z)&=\frac1{2\pi i}\int_{|w|=r_0}
 h(y,w)\,\frac{P_0(y,w)-P_0(y,z)}{(w-z)P_0(y,w)}\,\dd w,
 \label{finite:eq:ordinaryremainder}\\
 D_0h(y,z)&=\frac1{2\pi i}\int_{|w|=r_0}
 \frac{h(y,w)}{(w-z)P_0(y,w)}\,\dd w
 \qquad(|z|<r_0),
 \label{finite:eq:ordinaryquotient}
\end{align}
the numerator quotient in \eqref{finite:eq:ordinaryremainder} being a polynomial in $z$ of degree $<m$, so that the apparent pole at $w=z$ is removable.

\begin{lemma}[Uniform ordinary division operators]
\label{finite:lem:ordinarydivision}
The operators \eqref{finite:eq:ordinaryremainder}--\eqref{finite:eq:ordinaryquotient} are $\OO(B)$-linear, $R_0h\in\OO(B)[z]$ has degree $<m$, $D_0h$ extends holomorphically to all of $B\times D_R$ by $(h-R_0h)/P_0$ outside the root disk, and
\begin{equation}\label{finite:eq:ordinarydivision}
 h=R_0h+P_0\,D_0h,\qquad \deg_zR_0h<m .
\end{equation}
The decomposition is unique, $D_0(P_0q)=q$, and $D_0r=0$ for polynomials $r$ of degree $<m$. Both outputs are holomorphic on the same $B\times D_R$. For $P_0=z^m$ the operators are the elementary ones
\begin{equation}\label{finite:eq:monomialoperators}
 R_mh(y,z)=\sum_{k=0}^{m-1}\frac{z^k}{k!}\,(\partial_z^kh)(y,0),
 \qquad
 D_mh=\frac{h-R_mh}{z^m}.
\end{equation}
\end{lemma}

\begin{proof}
The second integrand in \eqref{finite:eq:ordinaryremainder} is polynomial of degree $\leq m-1$ in $z$, so its integral is a polynomial for every $z$; the denominator is nonzero on the circle $|w|=r_0$, and differentiation under this compact contour gives holomorphic dependence on $y$. Adding the two formulas reproduces the Cauchy integral, which is \eqref{finite:eq:ordinarydivision} for $|z|<r_0$; outside the root disk $(h-R_0h)/P_0$ is holomorphic and agrees on an annulus, so $D_0h\in\OO(B\times D_R)$. If $P_0q+r=0$ with $\deg_zr<m$, then for each fixed $y$ the polynomial $r(y,\cdot)$ vanishes at all roots of the monic $P_0(y,\cdot)$ with multiplicity, hence vanishes; then $q=0$. This is a convenient fixed-domain form of ordinary Weierstrass division \cite{finite:GunningRossi,finite:Lebl}.
\end{proof}

\subsection{Support-controlled matrix division}
\label{finite:sec:matrixdivision}

The scalar case of the next theorem suffices for one equation. The \emph{matrix} case is what controls the coupling when several equations are eliminated in succession: multiplication by the $k$th equation on the previous quotient is represented by a matrix. Its entries lie in a commutative ring, but its ordinary coefficient matrices at different Hahn exponents need not commute, nor need matrix multiplication commute with the division operator.

\begin{theorem}[Support-controlled matrix division]
\label{finite:thm:matrixdivision}
Let $P_0$ satisfy \cref{finite:lem:ordinarydivision} on $B\times D_R$, let $q\geq1$, and let
\[
 \mathcal F=P_0I_q+\mathcal E,
 \qquad \mathcal E\in\Mat_q\bigl(\Ip_\Gamma(B\times D_R)\bigr).
\]
Let $E\subseteq\Gamma_{>0}$ be a well-ordered set containing the supports of all entries of $\mathcal E$, and $M=E^*$. For every $G\in\HH_\Gamma(B\times D_R)^q$ there are \emph{unique}
\[
 Q\in\HH_\Gamma(B\times D_R)^q,
 \qquad
 \mathcal R\in\HH_\Gamma(B)[z]^q
 \ \text{with}\ \deg_z\mathcal R_a<m,
\]
such that $G=\mathcal FQ+\mathcal R$. If $\supp G\subseteq S$ then $\supp Q\cup\supp\mathcal R\subseteq S+M$, all coefficient functions remain holomorphic on the original $B\times D_R$, and nonnegative input support gives nonnegative output support. Consequently
\begin{equation}\label{finite:eq:matrixmodule}
 \Hp_\Gamma(B\times D_R)^q\big/\mathcal F\,\Hp_\Gamma(B\times D_R)^q
 \simeq\Hp_\Gamma(B)^{qm},
\end{equation}
with basis the classes of $z^k\varepsilon_a$, $0\leq k<m$, $1\leq a\leq q$, and likewise over $\HH_\Gamma$.
\end{theorem}

\begin{proof}
Apply $D_0$ componentwise to $G=\mathcal FQ+\mathcal R$. Since $D_0\mathcal R=0$ and $D_0(P_0Q)=Q$, the equation becomes $Q+D_0(\mathcal EQ)=D_0G$. The operator $Q\mapsto D_0(\mathcal EQ)$ multiplies by an ordinary holomorphic matrix of positive Hahn exponent and then applies the fixed operator $D_0$, so every application introduces a letter of $E$ and \cref{finite:lem:operator} gives
\begin{equation}\label{finite:eq:matrixQ}
 Q=\sum_{k\geq0}\bigl(-D_0\mathcal E\bigr)^kD_0G,
 \qquad
 \mathcal R=R_0\bigl(G-\mathcal EQ\bigr),
\end{equation}
with support in $S+M$ and coefficients on the original domain; here $D_0\mathcal E$ denotes the composite operator, not a matrix of derivatives, and its finite compositions are \emph{ordered} words. Splitting $G-\mathcal EQ$ by \eqref{finite:eq:ordinarydivision} gives $G=\mathcal FQ+\mathcal R$. For uniqueness, apply $D_0$ to the difference of two divisions: the quotient difference solves $(1+D_0\mathcal E)\Delta Q=0$ and vanishes by \cref{finite:lem:operator}, whence the remainder difference vanishes. Uniqueness says the bounded-degree vector polynomials form an actual complementary module, giving \eqref{finite:eq:matrixmodule}.
\end{proof}

\begin{theorem}[Fixed-domain Hahn preparation with holomorphic parameters]
\label{finite:thm:parameterprep}
Let $P_0$ be as above and
\[
 F=u_0\,(P_0+E_0),
 \qquad u_0\in\OO(B\times D_R)^\times,
 \qquad E_0\in\Ip_\Gamma(B\times D_R),
\]
and let $T\subseteq\Gamma_{>0}$ be well ordered and contain $\supp E_0$. There is a \emph{unique} normalized factorization
\begin{equation}\label{finite:eq:parameterprep}
 F=u_0\,P\,Q,
 \qquad P=P_0+A,\ \ \deg_zA<m,
 \qquad Q=1+C,
\end{equation}
with $A\in\Ip_\Gamma(B)[z]$, $C\in\Ip_\Gamma(B\times D_R)$ and $\supp A\cup\supp C\subseteq T^*\setminus\{0\}$. The factor $Q$ is a coherent unit, nonzero throughout $(B\times D_R)\sh$. Moreover division by $P$, hence by $F$, is the case $q=1$ of \cref{finite:thm:matrixdivision}, so
\[
 \HH_\Gamma(B\times D_R)/(F)\simeq\HH_\Gamma(B)^m
\]
with basis $1,z,\ldots,z^{m-1}$, with the same $S+T^*$ certificate, and likewise over $\Hp_\Gamma$. All these constructions commute with specialization of $y$ at points of $B\sh$, after placing the specialization in a common workspace.
\end{theorem}

\begin{proof}
The preparation equation is $E_0=A+P_0C+AC$. Proceed by well-founded recursion through $(T^*)_{>0}$: at exponent $\nu$ all terms of
\[
 h_\nu=(E_0)_\nu-\sum_{\substack{\alpha+\beta=\nu\\ \alpha,\beta>0}}A_\alpha C_\beta
\]
are known from strictly smaller exponents and the sum is finite by \cref{finite:lem:support}; set $A_\nu=R_0h_\nu$ and $C_\nu=D_0h_\nu$, which solves the coefficient equation by \eqref{finite:eq:ordinarydivision} and keeps all coefficients on the original domain. If a second normalized factorization differs, at the least exponent of difference $\Delta A+P_0\Delta C=0$, and ordinary division forces both differences to vanish; taking the well-ordered union of the two candidate supports legitimizes the least-exponent argument even if the rival factorization was not initially supported in $T^*$. Division is \eqref{finite:eq:matrixQ} with $q=1$ and $\mathcal E=A$. For specialization, the contour operators, coefficient differentiation in $y$, and substitution of an infinitesimal displacement in $y$ commute --- first for one ordinary coefficient by differentiation under the fixed contour, then for the Taylor series by strong summability --- and the extra exponents introduced belong to the monoid generated by the displacement support.
\end{proof}

\begin{corollary}[One shrink, and no more]
\label{finite:cor:oneshrink}
If the exponent-zero coefficient of $F$ is $z$-regular of order $m$ at an ordinary point $(y_0,z_0)$, then ordinary Weierstrass preparation followed by \emph{one} shrinking of the parameter polydisk and the $z$-disk puts $F$ into the hypotheses of \cref{finite:thm:parameterprep}. No subsequent shrinking is made at any Hahn exponent.
\end{corollary}

\begin{proof}
Ordinary preparation gives $f_0=u_0P_0$ near the point. Shrink once so that $u_0$ is nowhere zero and all roots of the monic $P_0$ remain inside a fixed smaller $z$-circle; dividing the positive-support error by $u_0$ changes no exponents and preserves holomorphy. Apply the theorem.
\end{proof}

The absence of repeated shrinking is not a formality. A Hahn family of holomorphic germs that admits no common neighbourhood is not an element of $\HH_\Gamma(B\times D_R)$ for any fixed domain, nor of the common-domain germ ring in \cref{finite:sec:rings}(4); it belongs only to the radius-free category in \cref{finite:sec:rings}(3). An infinite sequence of choices of smaller neighbourhoods need not retain any usable ordinary neighbourhood at all. The proofs here never quietly enlarge the function class at an infinite stage.

%======================================================================
\section{The support-controlled perturbation lemma}
\label{finite:sec:perturbation}

\subsection{Contractions concentrated in degree zero}
\label{finite:sec:contraction}

Let $(C_\bullet,d)$ be a chain complex of complex vector spaces in nonnegative degrees whose only homology is a vector space $V$ in degree zero. A \emph{contraction} onto $V$ is a triple $(i,p,h)$, $h$ of degree $+1$, with
\begin{equation}\label{finite:eq:contraction}
 pi=1,\qquad dh+hd=1-ip,\qquad h^2=0,\quad hi=0,\quad ph=0,
\end{equation}
where $p$ is taken to be zero in positive degrees and $i(V)\subseteq C_0$; note also $pd=0$ and $di=0$.

\begin{construction}[Existence over $\C$]
\label{finite:constr:h}
Let $B_k=\im(d:C_{k+1}\to C_k)$. In degree $0$ choose $C_0=i(V)\oplus B_0$. In each positive degree, exactness identifies $B_k$ with $\ker d_k$; choose a complex-vector-space complement $L_k$ of $B_k$ in $C_k$, so that $d:L_k\to B_{k-1}$ is an isomorphism. Define $h$ to be its inverse on each $B_{k-1}$, and zero on each $L_k$ and on $i(V)$. All identities in \eqref{finite:eq:contraction} hold on the displayed summands.
\end{construction}

\begin{remark}[What the splitting is, and what it is not]
\label{finite:rem:splitting}
The map $h$ need not be $\OO(D)$-linear and need not be continuous in a Fr\'echet topology. It maps holomorphic functions on $D$ to holomorphic functions on the \emph{same} $D$, because it is chosen inside the complex of global sections on $D$; and its coefficientwise extension preserves support. This is exactly why no analytic norm estimate and no sequence of smaller polydisks is ever needed to iterate it along positive Hahn exponents, and it is the fixed-domain feature on which the whole construction depends. The existence proof uses ordinary choices of complex-vector-space complements; it is not, by itself, an effective algorithm for an arbitrary black-box holomorphic system, and this is precisely where a noncanonical choice enters. \Cref{finite:thm:family,finite:thm:finitefree15,finite:cor:coefficientresidue} supply explicit operators in the cases where they are available.
\end{remark}

\subsection{The Hahn contraction formula}
\label{finite:sec:hpl}

Pass coefficientwise to $C_\bullet((t^\Gamma))$ and let
\[
 \delta=\sum_{e\in E}t^e\delta_e
\]
be a degree $-1$ perturbation with $E\subseteq\Gamma_{>0}$ well ordered and $(d+\delta)^2=0$; in the Koszul application $\delta i=0$ automatically, since the complex has no negative degree.

\begin{theorem}[Hahn contraction formula]
\label{finite:thm:hpl}
The following operators are well defined and strongly summable:
\begin{align}
 S&=(1+\delta h)^{-1}=\sum_{k\geq0}(-\delta h)^k,
 &\mathcal A&=S\delta=\delta(1+h\delta)^{-1},
 \label{finite:eq:hplS}\\
 p_F&=pS=p-p\mathcal Ah,
 &h_F&=hS=h-h\mathcal Ah.
 \label{finite:eq:hplPH}
\end{align}
With the original $i$ they satisfy
\begin{equation}\label{finite:eq:hplidentity}
 p_Fi=1,\qquad p_F(d+\delta)=0,\qquad
 (d+\delta)h_F+h_F(d+\delta)=1-ip_F,
\end{equation}
together with $h_F^2=0$, $h_Fi=0$, $p_Fh_F=0$. Hence the perturbed complex has vanishing positive-degree homology and degree-zero homology $V((t^\Gamma))$. Each operator in \eqref{finite:eq:hplPH} sends support $S_0$ into $S_0+M$, $M=E^*$, and all assertions hold verbatim on the nonnegative-support submodules.
\end{theorem}

\begin{proof}
All displayed inverses exist by \cref{finite:lem:operator}, because each occurrence of $\delta$ introduces a letter of $E$ while $h,d,p,i$ introduce none; the $k$th summand of $S$ therefore has support in $S_0+kE$, and \cref{finite:lem:support} supplies well-ordering and finite contribution over all $k$ at once. Consequently the following operator algebra is a coefficientwise finite identity at each exponent.

From $(1+\delta h)\delta=\delta(1+h\delta)$ come the two expressions for $\mathcal A$ and the relations
\begin{equation}\label{finite:eq:Arelations}
 \delta h\mathcal A=\mathcal Ah\delta=\delta-\mathcal A,
 \qquad \mathcal Ai=0 .
\end{equation}
Multiplying $d\mathcal A+\mathcal Ad+\mathcal Aip\mathcal A$ on the left by $1+\delta h$ and on the right by $1+h\delta$ gives
\[
 d\delta+\delta d+\delta(hd+dh+ip)\delta=d\delta+\delta d+\delta^2=0,
\]
and both multipliers are invertible, so with $\mathcal Ai=0$,
\begin{equation}\label{finite:eq:dA}
 d\mathcal A+\mathcal Ad=0 .
\end{equation}
Now
\[
 p_F(d+\delta)=p\delta-p\mathcal Ahd-p\mathcal Ah\delta
 =p\mathcal A(1-hd)=p\mathcal A(dh+ip)=0,
\]
using $pd=0$, \eqref{finite:eq:dA} and $\mathcal Ai=0$. Expanding the homotopy identity and using \eqref{finite:eq:Arelations}, its left side is $1-ip+\mathcal Ah+h\mathcal A-dh\mathcal Ah-h\mathcal Ahd$; substituting $dh=1-ip-hd$ and using \eqref{finite:eq:dA} reduces it to $1-ip+ip\mathcal Ah=1-ip_F$. Finally $p_Fi=pi-p\mathcal Ahi=1$, and $Sh=h$ gives $h_F^2=hShS=0$ and $p_Fh_F=pShS=0$. The support bounds and their nonnegative-support versions are term-by-term consequences of \cref{finite:lem:operator}.
\end{proof}

\begin{remark}[The same engine, seen three ways; and what is not new]
\label{finite:rem:conjugation}
The same operator content admits two further presentations, both used in the sources, and it is worth recording that they agree.
\begin{enumerate}
\item \emph{Conjugation.} With $Q=1+\delta h$ and $S=Q^{-1}$ one has the \emph{exact} conjugacy
\begin{equation}\label{finite:eq:conjugation}
 (d+\delta)\,Q=Q\,d,
 \qquad\text{hence}\qquad
 d+\delta=QdS,\quad S(d+\delta)=dS .
\end{equation}
Indeed $(d+\delta)Q=d+\delta+d\delta h+\delta^2h=d+\delta-\delta dh$ by $d\delta+\delta d+\delta^2=0$, while $Qd=d+\delta hd=d+\delta(1-ip-dh)=d+\delta-\delta dh$ by \eqref{finite:eq:contraction} and $\delta i=0$. Since $Qi=i$ and $Qh=h$, conjugating $dh+hd=1-ip$ by $Q$ and $S$ gives \eqref{finite:eq:hplidentity} in one line. This is an \emph{actual} conjugation, not an asymptotic one.
\item \emph{Resolvent.} The identities \eqref{finite:eq:Arelations}--\eqref{finite:eq:dA} are the resolvent identities used to transport Koszul relations in \cref{finite:cor:koszul} and to prove the precision theorem \cref{finite:thm:precision} through the difference identity $S_F-S_{\widetilde F}=S_F(\widetilde\delta-\delta)hS_{\widetilde F}$.
\end{enumerate}
None of this is a new abstract perturbation lemma. Invertibility of $1+\delta h$ is already the appropriate algebraic smallness condition in the classical lemma \cite{finite:Crainic}, whose general versions allow homology in several degrees and an induced differential on that homology; here that induced differential vanishes for a simple reason, namely that the target is concentrated in degree zero. What has been checked is that the inverse exists on the \emph{full} Hahn modules, that it respects one ordinary coefficient domain, and that it carries the exact support estimate the surcomplex application requires.
\end{remark}

%======================================================================
\section{Division and finite freeness over the fixed-polydisk ring}
\label{finite:sec:division}

\begin{hypothesis}[The fixed ordinary system]
\label{finite:hyp:ordinary}
$D\subset\C^n$ is a polydisk centred at $0$; $f_1,\ldots,f_n\in\OO(D)$; their common zero set \emph{in $D$} is exactly $\{0\}$. Put $A_0=\OO(D)/(f_1,\ldots,f_n)$ and $\mu=\dim_\C A_0$.
\end{hypothesis}

\begin{proposition}[The ordinary quotient and its Koszul resolution]
\label{finite:prop:ordinarykoszul}
Under \cref{finite:hyp:ordinary}, restriction to the germ at $0$ induces
\begin{equation}\label{finite:eq:ordinarylocal}
 A_0\simeq\OO_{\C^n,0}/(f_1,\ldots,f_n),
 \qquad 1\leq\mu<\infty,
\end{equation}
and the Koszul complex $C_\bullet=K_\bullet(f;\OO(D))$ of \emph{global} sections, with differential
\begin{equation}\label{finite:eq:koszul}
 d\bigl(g\,\varepsilon_{j_1}\wedge\cdots\wedge\varepsilon_{j_k}\bigr)
 =\sum_{a=1}^{k}(-1)^{a-1}f_{j_a}g\,
 \varepsilon_{j_1}\wedge\cdots\widehat{\varepsilon_{j_a}}\cdots\wedge\varepsilon_{j_k},
\end{equation}
is exact in positive homological degrees with degree-zero homology $A_0$. Some power of the maximal ideal of $A_0$ vanishes, so a basis may be, and is, chosen to consist of polynomials $e_1=1,e_2,\ldots,e_\mu$.
\end{proposition}

\begin{proof}
At $0$ the local ring $\OO_{\C^n,0}$ is regular of dimension $n$ and the ideal is primary to the maximal ideal because the zero is isolated; hence the $f_i$ form a system of parameters and therefore a regular sequence, and their Koszul complex is exact in positive degrees there. At every other point of $D$ some $f_i$ is a unit, so the Koszul complex of sheaves is contractible there. Thus the complex of coherent sheaves is exact in positive degrees on all of $D$, and its degree-zero quotient sheaf is supported at $0$ with the finite stalk \eqref{finite:eq:ordinarylocal}. $D$ is Stein, so Cartan's Theorem~B applied to the coherent kernels and images of this finite complex makes the complex of global sections exact in the same degrees and the global quotient map onto the supported sheaf surjective; a sheaf supported at one point has its finite stalk as global sections. These are ordinary complex-analytic facts \cite{finite:Cartan,finite:GunningRossi,finite:Demailly}.
\end{proof}

Fix a contraction $(i,p,h)$ of $C_\bullet$ as in \cref{finite:constr:h}, with $i$ sending the coordinate basis of $\C^\mu$ to the chosen polynomials. Put $F_j=f_j+E_j$ with $\supp E_j\subseteq\Gamma_{>0}$, let $\delta$ be the Koszul differential of the $E_j$, $\partial=d+\delta$ the Koszul differential of $F$, and $E,M$ as in \cref{finite:sec:notation}.

\begin{theorem}[Uniform finite-flat zero algebra; division with the $S+M$ certificate]
\label{finite:thm:finitefree}
Assume \cref{finite:hyp:ordinary}. Over the fixed-polydisk ring $\HH_\Gamma(D)$, for every $G\in\HH_\Gamma(D)$ there are $Q_1,\ldots,Q_n\in\HH_\Gamma(D)$ and a \emph{unique} $c(G)=(c_1,\ldots,c_\mu)\in K^\mu$ with
\begin{equation}\label{finite:eq:division}
 G=\sum_{a=1}^{\mu}c_a\,e_a+\sum_{j=1}^{n}F_jQ_j .
\end{equation}
One choice is
\begin{equation}\label{finite:eq:NFformula}
 c(G)=\NF_F(G)=p\,(1+\delta h)^{-1}G,
 \qquad (Q_1,\ldots,Q_n)=h\,(1+\delta h)^{-1}G,
\end{equation}
and these satisfy
\begin{equation}\label{finite:eq:divisionsupport}
 \supp c_a,\ \supp Q_j\subseteq\supp G+M,
 \qquad
 v(c_a),\ v(Q_j)\geq v(G),
\end{equation}
with all coefficient functions holomorphic on the original $D$. If $G\in\Hp_\Gamma(D)$ then $c_a\in\Kc$ and $Q_j\in\Hp_\Gamma(D)$. Consequently, with
\[
 \Alg^+=\Hp_\Gamma(D)/(F_1,\ldots,F_n),
 \qquad
 \Alg=\HH_\Gamma(D)/(F_1,\ldots,F_n),
\]
the unchanged classes of $e_1,\ldots,e_\mu$ are a basis and
\begin{equation}\label{finite:eq:finitefree}
 \Alg^+\simeq(\Kc)^\mu,\qquad \Alg\simeq K^\mu ,
\end{equation}
with $\Alg^+/\mm_K\Alg^+\simeq A_0$ as algebras. The witnesses $Q_j$ are \emph{not} claimed unique.
\end{theorem}

\begin{proof}
Each coefficient of $\delta$ is multiplication by an ordinary holomorphic coefficient of some $E_j$ composed with an exterior contraction, so $\delta$ satisfies the hypotheses of \cref{finite:thm:hpl} on the fixed vector spaces $C_k$. For a degree-zero input $\partial G=0$, so \eqref{finite:eq:hplidentity} reads
\[
 G=i\,\NF_F(G)+\partial\,h_FG,
\]
and writing the degree-one element $h_FG$ in the exterior basis produces the witnesses $Q_j$; this is \eqref{finite:eq:division}. The support and valuation estimates are those of \cref{finite:thm:hpl}, since $M\subseteq\Gamma_{\geq0}$.

For uniqueness and the module statement: $p_F\partial=0$ makes $p_F$ vanish on the ideal $(F)$, while $p_Fi=1$ makes it the identity on the chosen representatives. Hence a remainder lying in the ideal is a degree-zero boundary $i(c)=\partial Q$, and applying $p_F$ gives $c=0$. So $\ker p_F$ is exactly the ideal, proving both spanning and linear independence. Every operator preserves nonnegative support, so the same proof gives the assertion over $\Kc$ and not merely after inverting $t$: the quotient is genuinely \emph{free}, not merely torsion-free. At exponent zero $\delta$ contributes nothing, so the normal form and the multiplication reduce to their ordinary counterparts, which is the last assertion.
\end{proof}

\begin{remark}[Flatness without Noetherian hypotheses]
\label{finite:rem:nonnoetherian}
``Finite flat'' here follows from the strictly stronger, explicitly proved statement ``finite free''. No Noetherian deformation theorem is invoked for $\Kc$, which can have high-rank value group and need not be Noetherian. Ordinary Noetherian local algebra is used only once, to construct the exponent-zero resolution in \cref{finite:prop:ordinarykoszul}. In the usual algebraic sense $\Spec\Alg^+\to\Spec\Kc$ is finite flat of rank $\mu$; that statement is \emph{not} a claim about proper maps in an undeclared topology on the class $\SC^n$. Note also that the positive ideal of $\Hp_\Gamma(\mathrm{pt})$ is not finitely generated when $\Gamma$ is nonzero and divisible --- finitely many positive-valued generators have a least valuation $\delta>0$, and $t^{\delta/2}$ cannot lie in their ideal --- so a Noetherian completion theorem could not simply be assumed for this coefficient ring. This observation does not say that the \emph{field} $K$ is non-Noetherian.
\end{remark}

\begin{corollary}[Explicit lifting of Koszul relations]
\label{finite:cor:koszul}
The Koszul complex $K_\bullet(F;\Hp_\Gamma(D))$ has zero homology in positive degrees, and likewise over $\HH_\Gamma(D)$. If a cycle $Z$ of positive degree has support $S_0$, then $Z=\partial\,h_FZ$, and the explicit preimage $h_FZ$ has support in $S_0+M$ with coefficient functions holomorphic on $D$.
\end{corollary}

\begin{proof}
$p_FZ=0$ by degree, so \eqref{finite:eq:hplidentity} reads $Z=\partial h_FZ$; the support bound is \cref{finite:thm:hpl}.
\end{proof}

This is the precise sense in which the ambiguity of the witnesses in \eqref{finite:eq:division} is governed: two choices of $(Q_1,\ldots,Q_n)$ differ by a Koszul syzygy, and every syzygy is transported by $h_F$ with the same support certificate. We call this Koszul exactness rather than invoking an unproved converse that would identify it with every notion of regular sequence in a possibly non-Noetherian ring.

\subsection{Structure constants, canonicity, and restriction}
\label{finite:sec:structure}

Transporting multiplication through \eqref{finite:eq:finitefree} gives structure constants
\begin{equation}\label{finite:eq:structure}
 e_ae_b\equiv\sum_{c=1}^{\mu}C_{ab}^{\,c}e_c\ \ \bmod (F),
 \qquad C_{ab}^{\,c}=\NF_F(e_ae_b)_c,
\end{equation}
which have support in $M$ and whose exponent-zero coefficients are the ordinary structure constants of $A_0$. Let $M_G$ denote multiplication by $[G]$ in this basis, its $a$th column being $\NF_F(Ge_a)$; the coordinate matrices $M_j=M_{z_j}$ lie in $\Mat_\mu(\Kc)$ and reduce to multiplication by $z_j$ on $A_0$. They commute, because they represent multiplication in a commutative quotient --- a conclusion of the construction, not a consistency condition to be guessed during a matrix recursion.

Although $h$ is noncanonical, the normal form relative to the fixed representative subspace $\mathrm{span}_\C\{e_a\}$ is canonical: two constructions both produce a representative of the same class inside that span, so their coordinate vectors agree. Different choices of $h$ can change the witnesses and the computational cost, but not $\Alg$, its multiplication, its traces, its zero set, its local lengths, its residue pairing or its discriminants.

\begin{corollary}[No loss of support under restriction]
\label{finite:cor:restriction}
If $D'\subseteq D$ is a smaller centred polydisk on which the same ordinary system still has $\{0\}$ as its only common zero, restriction induces an isomorphism of the two finite zero algebras, preserving the chosen polynomial basis and the multiplication matrices. The witnesses constructed on $D$ restrict to valid witnesses on $D'$, and no new support appears.
\end{corollary}

\begin{proof}
Restriction carries the ideal into the corresponding ideal and each basis vector to itself; both quotients have the displayed basis, so the map is an isomorphism. Uniqueness of the remainder then gives the last two assertions.
\end{proof}

\subsection{An exact coefficient recursion, with honest scope}
\label{finite:sec:recursion}

Write $W=(1+\delta h)^{-1}G$ and $\delta=\sum_{e\in E}t^e\delta_e$. Then
\begin{equation}\label{finite:eq:recursion}
 W_\nu=G_\nu-\sum_{\substack{e\in E,\ \eta\in\supp G+M\\ e+\eta=\nu}}\delta_e\,hW_\eta,
 \qquad
 \NF_F(G)=pW,\quad (Q_j)=hW,
\end{equation}
each displayed sum being finite by \cref{finite:lem:support}, with every $\eta$ on the right strictly smaller than $\nu$. This is an exact coefficient algorithm when the ordinary operations are effective and the complete finite list of contributing input exponents and support words can be computed for each requested $\nu$. Decidable support membership alone does not supply such a terminating enumeration. It is \emph{not} a blanket claim that every valuation truncation can be listed in finite time; \cref{finite:rem:truncation} is the obstruction, and \cref{finite:sec:computation} states the scope precisely.

\subsection{Explicit contractions: perturbed coordinate powers}
\label{finite:sec:family}

The existence proof above allows an arbitrary complex-linear contraction, which need not preserve holomorphic dependence on an additional parameter. To obtain a family theorem without hiding that issue we specify a parameter-linear contraction. For the $j$th variable put, as in \eqref{finite:eq:monomialoperators},
\begin{equation}\label{finite:eq:Rj}
 R_jg=\sum_{k=0}^{m_j-1}\frac{z_j^k}{k!}(\partial_j^kg)|_{z_j=0},
 \qquad
 D_jg=\frac{g-R_jg}{z_j^{m_j}} ,
\end{equation}
which are $\OO(B)$-linear, preserve holomorphy on $B\times D$, and satisfy $z_j^{m_j}D_j=1-R_j$, $D_jz_j^{m_j}=1$, $D_jR_j=0$. Operators of distinct coordinates commute.

\begin{theorem}[Perturbed monomial complete intersections, with holomorphic parameters]
\label{finite:thm:family}
Let $B\subset\C^r$ and $D\subset\C^n$ be polydisks, $D$ centred at $0$, and let
\begin{equation}\label{finite:eq:monomialfamily}
 F_j(y,z)=z_j^{m_j}+E_j(y,z),\qquad m_j\geq1,
\end{equation}
with $E_j\in\Ip_\Gamma(B\times D)$ of common well-ordered positive support $E$, $M=E^*$, and $\mu=\prod_jm_j$. Then, over $\HH_\Gamma(B\times D)$,
\begin{equation}\label{finite:eq:familyfree}
 \HH_\Gamma(B\times D)/(F_1,\ldots,F_n)
 \simeq\bigoplus_{\alpha\in\bx}\HH_\Gamma(B)\,z^\alpha,
 \qquad
 \bx=\{\alpha\in\N^n:0\leq\alpha_j<m_j\},
\end{equation}
as $\HH_\Gamma(B)$-modules, and likewise over the nonnegative-support rings. An input of support $S$ has normal-form coefficients and witnesses supported in $S+M$, all holomorphic on the original parameter and variable domains. The quotient and its multiplication commute with specialization at every $y\in B\sh$; each specialized zero cluster in $\mm_{\SC}^{\,n}$ has total local multiplicity $\mu$, and so does the cluster over every infinitesimal target in the $z$-equations, at every such parameter point.
\end{theorem}

\begin{proof}
In the two-term Koszul complex of $z_j^{m_j}$ let $h_j$ be $D_j$ from degree zero to degree one and zero in degree one, and let $\Pi_j=i_jp_j$ act as $R_j$ in degree zero and as zero in degree one, so $d_jh_j+h_jd_j=1-\Pi_j$. On the multivariable Koszul complex use the finite tensor-contraction formula
\begin{equation}\label{finite:eq:tensorh}
 h=h_1+\Pi_1h_2+\Pi_1\Pi_2h_3+\cdots+\Pi_1\cdots\Pi_{n-1}h_n,
\end{equation}
each degree-one operator placed in its own coordinate with the usual graded convention. One need not identify holomorphic functions with an algebraic tensor product: \eqref{finite:eq:tensorh} is a finite sum of explicitly defined coordinate operators acting on holomorphic functions and exterior components. The homotopy identity telescopes,
\[
 dh+hd=(1-\Pi_1)+\Pi_1(1-\Pi_2)+\cdots+\Pi_1\cdots\Pi_{n-1}(1-\Pi_n)=1-\Pi_1\cdots\Pi_n,
\]
and the side conditions hold because $h_j^2=0$ and $h_j\Pi_j=\Pi_jh_j=0$. The degree-zero image of $\Pi_1\cdots\Pi_n$ is exactly the free $\OO(B)$-module on the box monomials. Apply \cref{finite:thm:hpl} to the positive Koszul perturbation determined by the $E_j$: all operators are parameter-linear and remain on $B\times D$, which gives \eqref{finite:eq:familyfree} and the support estimates. Parameter differentiation commutes with these coordinate operators, so strong Taylor evaluation in $y$ commutes with every normal-form and multiplication formula; at a point of $B\sh$ the specialized data are again $z_j^{m_j}$ plus positive-support terms on the ordinary $z$-polydisk. The final assertions are then \cref{finite:thm:length,finite:cor:monadmap} in the enlarged workspace.
\end{proof}

\begin{remark}[An equivalent exterior formula, and the two-variable case]
\label{finite:rem:exteriorh}
The contraction \eqref{finite:eq:tensorh} may be written in closed exterior form as
\begin{equation}\label{finite:eq:exteriorh}
 h=\sum_{i=1}^{n}\varepsilon_i\wedge D_i\,P_1\cdots P_{i-1},
\end{equation}
where $P_i$ applies $R_i$ to the coefficient of an exterior monomial not containing $\varepsilon_i$ and returns zero otherwise; equivalently $h(g\,\varepsilon_J)=\sum_{i<\min J}P_1\cdots P_{i-1}D_i(g)\,\varepsilon_i\wedge\varepsilon_J$, with $\min\varnothing=n+1$. That $h^2=0$ is visible from the formula: after an $\varepsilon_i$ has been inserted, a later summand of larger index is killed by its preceding $P_i$, a summand of the same index wedges $\varepsilon_i$ twice, and a summand of smaller index $j$ is killed by $D_jR_j=0$. In two variables, on degree zero
\[
 hG=(D_1G)\,\varepsilon_1+(R_1D_2G)\,\varepsilon_2,
 \quad pG=R_1R_2G,
 \quad
 \delta hG=E_1D_1G+E_2R_1D_2G,
\]
so that $\NF_F(G)=R_1R_2\sum_{k\geq0}(-\delta h)^kG$. The asymmetry records an order of division; it does \emph{not} assert that the equations are triangular. This matters: formula \eqref{finite:eq:NFformula} is not merely a reduction rule that might depend on the order of reductions. The full exterior calculation is what proves that the remainder annihilates the \emph{whole ideal}, not merely the particular expression produced by one division pass. Ignoring syzygies would leave a genuine gap in a several-variable argument.
\end{remark}

\begin{remark}[A more general family criterion]
\label{finite:rem:familycriterion}
The same proof works for any ordinary parameter-dependent complete intersection whose fixed-domain Koszul complex is supplied with an $\OO(B)$-linear contraction onto a finite free $\OO(B)$-module. This is a precise additional hypothesis, not an assertion that every family automatically has such a splitting on any prescribed domain. More abstractly: if an ordinary Koszul complex on a fixed domain admits complex-linear $(i,p,h)$ satisfying \eqref{finite:eq:contraction} with $p$ onto a finite-dimensional degree-zero space, all maps preserving the chosen ordinary function spaces, then every positive-support Hahn perturbation has a finite-dimensional quotient of the same dimension, with a canonical remainder and the same $S+M$ bound; the proof of \cref{finite:thm:finitefree} applies verbatim, since \eqref{finite:eq:conjugation} uses only \eqref{finite:eq:contraction}, anticommutation of the Koszul differentials, and $\delta i=0$. Recovering the actual-point and residue theorems for a new leading singularity additionally requires a finite-jet or equivalent property sufficient for \cref{finite:lem:matrix,finite:lem:witness}, and a compatible residue normalization.
\end{remark}

\subsection{The same theorem by matrix elimination}
\label{finite:sec:matrixroute}

For coordinate-power leading systems there is a second route to \eqref{finite:eq:familyfree} which does not construct a contraction at all, and which exhibits explicitly where the coupling is controlled. We record it because its intermediate statement --- that multiplication by the $k$th equation on the previous quotient is represented by a matrix of the form $w_k^{m_k}I_q+\mathcal E_k$ with $\mathcal E_k$ of positive support --- is what a one-variable root count cannot supply.

\begin{theorem}[Finite-free complete intersections by matrix induction]
\label{finite:thm:finitefree15}
Let $U$ be a connected ordinary parameter domain, $V=U\times D_{r_1}\times\cdots\times D_{r_n}$, and $F_i=w_i^{m_i}+E_i$ with $E_i\in\Ip_\Gamma(V)$, $E=\bigcup_i\supp E_i$, $M=E^*$, $\mu=\prod_im_i$. Then over $\HH_\Gamma(V)$:
\begin{enumerate}
\item $\Alg^+=\Hp_\Gamma(V)/(F_1,\ldots,F_n)$ is free of rank $\mu$ over $\Hp_\Gamma(U)$ on the classes of $w^\alpha$, $\alpha\in\bx$, and likewise with $\HH_\Gamma$ in place of $\Hp_\Gamma$;
\item every $G$ of support $S$ has $G=\sum_iF_iQ_i+\sum_{\alpha\in\bx}r_\alpha(x)w^\alpha$ with all outputs supported in $S+M$ and all coefficients holomorphic on the original domains, the $r_\alpha$ unique and the $Q_i$ not;
\item $F_1,\ldots,F_n$ is a regular sequence \emph{in any order}, and the ideal is proper;
\item all structure constants have support in $M$, with exponent-zero reduction the multiplication table of $\OO(U)[w]/(w_1^{m_1},\ldots,w_n^{m_n})$;
\item the normal form commutes with coherent finite parameter substitution and is independent of the elimination order.
\end{enumerate}
The conclusions persist after dividing equations whose leading terms are $u_i(x,w)w_i^{m_i}$ with $u_i$ ordinary holomorphic and zero-free on $V$, and after replacing $F_i$ by $F_i-b_i(x)$ for any $b_i\in\Ip_\Gamma(U)$.
\end{theorem}

\begin{proof}
Induct on the number of eliminated equations, the inductive statement being that after $k-1$ equations the module is free over $\Hp_\Gamma(U\times D_{r_k}\times\cdots\times D_{r_n})$ on $w_1^{\alpha_1}\cdots w_{k-1}^{\alpha_{k-1}}$, with the $S+M$ bound, and that each remaining equation acts by a matrix whose exponent-zero reduction is scalar multiplication by its leading coordinate power. For $k=1$ this is \cref{finite:thm:matrixdivision} with $q=1$ and $P_0=w_1^{m_1}$, all other variables being parameters. For the step, put $q=m_1\cdots m_{k-1}$ and represent multiplication by $F_k$ on the preceding quotient in the established basis; since the normal form reduces at exponent zero to ordinary division by the preceding coordinate powers, that matrix is
\begin{equation}\label{finite:eq:inducedmatrix}
 \mathcal F_k=w_k^{m_k}I_q+\mathcal E_k,
\end{equation}
every entry of $\mathcal E_k$ having positive support inside $M$ and being holomorphic on $U\times D_{r_k}\times\cdots\times D_{r_n}$. Apply \cref{finite:thm:matrixdivision} in the variable $w_k$. Since $M+M=M$ the support estimate is still $S+M$, and no domain is shrunk. Lifting a vector quotient back to the original ring uses the finitely many identities expressing each matrix column as the normal form of a product, whose differences from the actual products are combinations of the preceding $F_i$; products and sums of these preserve $S+M$. For regularity, a nonzero vector has a least exponent $\nu$, at which its image under \eqref{finite:eq:inducedmatrix} is $w_k^{m_k}$ times its nonzero ordinary coefficient vector, and the positive matrix error cannot contribute at $\nu$; so multiplication by $F_k$ is injective, in any order. The remaining assertions follow from uniqueness of the rectangular remainder, and the corollary from the fact that division by an ordinary zero-free unit changes no exponents and no ideal, while subtracting a positive-support target leaves the same leading system.
\end{proof}

The coefficient matrices in \eqref{finite:eq:inducedmatrix} need not commute with one another; their finite compositions are ordered words in \eqref{finite:eq:matrixQ}. This is exactly the point at which the coupling between the equations is controlled, and it is why the hypotheses may allow errors such as Hahn sums of $\exp(w_1w_2+xw_3)$, with no polynomial bound in the fibre variables and no independence from the coordinate being eliminated.

%======================================================================
\section{From the finite algebra to actual surcomplex zeros}
\label{finite:sec:points}

A character of an algebra of functions may not be identified with point evaluation merely by analogy; a finite-dimensional algebra is not by itself a theorem about actual analytic zeros, and an arbitrary algebra homomorphism need not commute with an unspecified infinite summation. This section closes that gap explicitly, over the fixed-polydisk ring $\HH_\Gamma(D)$ of \cref{finite:sec:rings}(1), by two independent devices: shifted division (\cref{finite:lem:shifted}) and a summable matrix functional calculus (\cref{finite:lem:matrix}). Either suffices; both are recorded because they have different ranges of applicability.

\subsection{Coordinates are infinitesimal}
\label{finite:sec:infinitesimal}

\begin{lemma}[Coordinate characteristic polynomials]
\label{finite:lem:coordpoly}
Each $P_j(T)=\det(TI-M_j)$ lies in $\Kc[T]$, is monic of degree $\mu$, and satisfies $P_j(T)\equiv T^\mu \bmod\mm_K$. Every root of $P_j$ in $K$, or in any valued field extension with the same ordered comparison, is infinitesimal.
\end{lemma}

\begin{proof}
The maximal ideal of $A_0$ is nilpotent, so multiplication by $z_j$ is nilpotent there; this gives the reduction. Write $P_j(T)=T^\mu+\sum_{k<\mu}a_kT^k$ with $a_k\in\mm_K$. If a nonzero root $\lambda$ had $v(\lambda)\leq0$, dividing its equation by $\lambda^\mu$ would give $1+\sum_{k<\mu}a_k\lambda^{k-\mu}=0$ with every summand after $1$ of positive valuation, which is impossible. The root $0$ is already infinitesimal.
\end{proof}

\begin{lemma}[Summable matrix functional calculus]
\label{finite:lem:matrix}
Let the leading system be the coordinate powers $z_i^{m_i}$ and let $X_i=M_{z_i}$ in the box basis, so that $X_i=N_i+Y_i$ with $N_i$ the ordinary nilpotent box-multiplication matrices and $v(Y_i)>0$. For $g\in\OO(D)$ with Taylor coefficients $g_\alpha$ at $0$, the matrix family $(g_\alpha X^\alpha)_{\alpha\in\N^n}$ is strongly summable and
\begin{equation}\label{finite:eq:matrixcalculus}
 M_g=\sum_{\alpha\in\N^n}g_\alpha X^\alpha ;
\end{equation}
for $G=\sum_\gamma g_\gamma t^\gamma\in\HH_\Gamma(D)$ the corresponding doubly indexed family is strongly summable with sum $M_G$, supported in $\supp G+M$.
\end{lemma}

\begin{proof}
Put $b=\sum_i(m_i-1)$, so every product of more than $b$ of the $N_i$ vanishes. Expand a matrix word in $X$'s into $N$ and $Y$ factors: a nonzero term with $k$ occurrences of $Y$ has at most $k+1$ intervening blocks of $N$ factors, each of length at most $b$, hence total length at most $(k+1)b+k$. At a fixed exponent \cref{finite:lem:support} allows only finitely many $Y$-words and lengths, and for each there are finitely many matrix words of that bounded length; this is strong summability, with support in $M$, or $\supp G+M$ with the Hahn index included. For a Taylor polynomial \eqref{finite:eq:matrixcalculus} is just multiplicativity of the quotient map. To pass to a general $g$, fix a Hahn exponent: on the normal-form side, a jet estimate shows that sufficiently high Taylor terms contribute nothing at that exponent --- an $R_i$ cannot lower total $z$-adic order, a $D_i$ lowers it by at most $m_i$, multiplication by a holomorphic function cannot lower it, and $p$ kills every monomial of total degree $>b$, so a term $p\,L^kG$ depends on Taylor data only through total degree $b+k\max_im_i$; on the matrix side the word argument gives the same eventual vanishing. Hence the polynomial identity proves the identity at that exponent, and one repeats at each exponent.
\end{proof}

\begin{theorem}[All characters are evaluations]
\label{finite:thm:characters}
Over $\HH_\Gamma(D)$ and under \cref{finite:hyp:ordinary}, evaluation induces a bijection between the common zeros of $F\sh$ in $D\sh_K$ and the $K$-algebra characters $\Alg\to K$, the character of $\zeta$ being $\chi_\zeta([G])=G\sh(\zeta)$. Every such zero lies in $\mm_K^n$; there are at least one and at most $\mu$ of them; and every zero of $F\sh$ in the full surcomplex monad $\mm_{\SC}^{\,n}$ already belongs to $\mm_K^n$, so that no additional zeros appear on passing to the class field.
\end{theorem}

\begin{proof}
If $c+\varepsilon$ is a coherent zero then \cref{finite:prop:evaluation} gives $f_i(c)=0$ for every $i$, hence $c=0$ by \cref{finite:hyp:ordinary}; so all zeros lie in $\mm_K^n$ and evaluation there is an algebra homomorphism annihilating the ideal, that is, a character.

Conversely let $\chi$ be a character and $\zeta_j=\chi([z_j])$. A nonzero finite-dimensional commutative algebra over the algebraically closed $K$ has finitely many maximal ideals, all with residue field $K$. Cayley--Hamilton gives $P_j(\zeta_j)=0$, so $\zeta\in\mm_K^n$ by \cref{finite:lem:coordpoly}. Apply $\chi$ to \eqref{finite:eq:shifteddivision}: every term $(z_j-\zeta_j)Q_j$ dies and
\[
 \chi([G])=G\sh(\zeta).
\]
This is a \emph{finite} identity in the algebra; it assumes no continuity of $\chi$ and no commutation of $\chi$ with an infinite Hahn sum. For coordinate-power leading systems there is also the alternative \eqref{finite:eq:matrixcalculus}: apply its matrix identity to the coordinate vector of $[1]$, and then apply $\chi$ as a linear functional on the finite basis. Multiplication by finitely many fixed scalars and finite addition preserve strong summability, giving the same conclusion. Taking $G=F_j$ shows $\zeta$ is a zero; distinct zeros give distinct values on some coordinate, and the two assignments are inverse.

Finally, $P_j(z_j)\in(F)$ is an identity in $\HH_\Gamma(D)$, so it may be evaluated at any proposed zero in the full surcomplex monad after enlarging the workspace to contain that point's supports. Each coordinate is then a root of $P_j\in K[T]$, which already splits over the algebraically closed $K$; hence the coordinate lies in $K$.
\end{proof}

\subsection{Local lengths are the multiplicities at the displaced zeros}
\label{finite:sec:lengths}

For a zero $\zeta$ put
\begin{equation}\label{finite:eq:localtaylor}
 \widehat F_{j,\zeta}(u)
 =\sum_{\alpha\in\N^n}\frac{(\partial^\alpha F_j)\sh(\zeta)}{\alpha!}\,u^\alpha
 \in K[[u_1,\ldots,u_n]],
 \qquad
 B_\zeta=\frac{K[[u]]}{(\widehat F_{1,\zeta},\ldots,\widehat F_{n,\zeta})} .
\end{equation}
Finiteness of $B_\zeta$ is \emph{proved} below, not built into the definition. The expansion is taken at the actual displaced zero $\zeta$, not at its standard part $0$; the distinction is not cosmetic, as \cref{finite:ex:unit} shows.

\begin{theorem}[Local algebra comparison and exact conservation]
\label{finite:thm:length}
Let $\Alg_\zeta$ be the local Artinian factor of $\Alg$ belonging to $\chi_\zeta$. Formal Taylor expansion induces a canonical isomorphism
\begin{equation}\label{finite:eq:localiso}
 \Alg_\zeta\simeq B_\zeta,
 \qquad
 [z_j]-\zeta_j\longleftrightarrow u_j .
\end{equation}
In particular every $B_\zeta$ is finite dimensional, every zero is formally isolated, and
\begin{equation}\label{finite:eq:conservation}
 \sum_{F\sh(\zeta)=0}\mu_\zeta=\mu,
 \qquad \mu_\zeta:=\dim_KB_\zeta .
\end{equation}
A zero $\zeta$ is \emph{simple}, that is $\mu_\zeta=1$, exactly when $J_F\sh(\zeta)\neq0$; if \emph{every} zero is simple there are exactly $\mu$ distinct zeros. The same local dimensions are obtained after extending scalars from $K$ to any larger algebraically closed Hahn workspace, or to $\SC$ in the finite-dimensional computations.
\end{theorem}

\begin{proof}
Taylor expansion at $\zeta$ is an algebra homomorphism $\HH_\Gamma(D)\to K[[u]]$ by \cref{finite:lem:shifted}, sending $F_j$ to $\widehat F_{j,\zeta}$, hence inducing $\Alg\to B_\zeta$. Cayley--Hamilton gives $P_j(z_j)\in(F)$; at the zero, $P_j(\zeta_j+u_j)=u_j^{r_j}U_j(u_j)$ with $r_j\geq1$ and $U_j(0)\neq0$, so $U_j$ is a formal unit and $u_j^{r_j}=0$ in $B_\zeta$. Therefore $B_\zeta$ is finite dimensional, spanned by the monomials $u^\alpha$ with $\alpha_j<r_j$, and the map onto it is surjective; it is a nonzero local algebra with residue field $K$, so it factors through $\Alg_\zeta$.

For the inverse, put $n_j=[z_j]-\zeta_j$ in $\Alg_\zeta$. By \cref{finite:lem:shifted} the maximal ideal is generated by the $n_j$, and it is nilpotent because the factor is Artinian; choose $N$ with its $N$th power zero. Formal substitution $u_j\mapsto n_j$ is then a \emph{finite} operation on each series, and \eqref{finite:eq:finitetaylor} applied to $F_j$ shows $\widehat F_{j,\zeta}(n)=0$, since the coherent remainder of order $N$ vanishes in the factor. This induces $B_\zeta\to\Alg_\zeta$. The two maps fix the coordinate generators, which generate both algebras over $K$, so they are mutually inverse; this is \eqref{finite:eq:localiso}. The product decomposition $\Alg\simeq\prod_\zeta\Alg_\zeta$ of a finite-dimensional commutative algebra --- obtained, for instance, by the Chinese remainder theorem applied to sufficiently high powers of the finitely many maximal ideals --- and $\dim_K\Alg=\mu$ from \cref{finite:thm:finitefree} give \eqref{finite:eq:conservation}.

For the simple-zero criterion: if the Jacobian is invertible, the formal inverse-function theorem, proved by degree-by-degree recursion with invertible linear part, makes the ideal of the $\widehat F_{j,\zeta}$ equal to $(u)$, so $B_\zeta=K$. Conversely if $\dim_KB_\zeta=1$ then its ideal is the maximal ideal, so the $n$ linear parts must span the $n$-dimensional cotangent space and their determinant is nonzero. Finally, all isomorphisms are between finite-dimensional spaces presented by finitely many linear relations among monomials of bounded degree; extension of scalars preserves the rank of such a presentation, and over the algebraically closed $K$ no residue-field splitting remains.
\end{proof}

\begin{corollary}[Conservation of multiplicity; a monad form of Rouch\'e]
\label{finite:cor:rouche}
Every strictly positive Hahn perturbation of the same ordinary system has total zero multiplicity exactly $\mu$ in the monad of $0$, regardless of how the zeros split, at what scales they lie, or whether they collide.
\end{corollary}

The word Rouch\'e here refers to conservation under a positive Hahn perturbation of the leading system. It is not an unrestricted boundary-norm theorem for arbitrary functions on $\SC^n$, and its hypothesis is agreement of the leading tuple rather than a scalar boundary inequality imported into several variables without a definition.

\begin{corollary}[Finite mapping of an entire monad]
\label{finite:cor:monadmap}
Over $\HH_\Gamma(D)$, the map $F\sh:\mm_{\SC}^{\,n}\to\mm_{\SC}^{\,n}$ is surjective, each fibre is finite, and each fibre has total local multiplicity $\mu$: for every infinitesimal target $\eta$,
\[
 \sum_{F\sh(\zeta)=\eta}
 \dim_{K'}K'[[u]]\big/
 \bigl(\widehat F_{1,\zeta}-\eta_1,\ldots,\widehat F_{n,\zeta}-\eta_n\bigr)=\mu,
\]
where $K'$ is any divisible Hahn workspace containing the original coefficients and the target.
\end{corollary}

\begin{proof}
The leading functions vanish at $0$, so evaluation at an infinitesimal input gives an infinitesimal output. For a target $\eta$, the system $F-\eta$ has the same leading complete intersection and still has well-ordered positive perturbation support after adjoining the finitely many target supports; apply \cref{finite:thm:characters,finite:thm:length}. The resulting algebra has positive dimension $\mu$, hence at least one character and therefore at least one root.
\end{proof}

This is a finite-mapping assertion in the coherent algebraic sense. It does not claim that a real-parameter path in a fine topology can sweep out a monad, and no compactness or properness of the fine-topological monad has been used. ``Finite map'' here means exactly the stated finite-fibre and finite-algebra properties; it does not assert properness for a topology on a proper class. Similarly, ``finite free'' is the explicitly proved statement, and finite flatness is its consequence.

%======================================================================
\section{The formal-coefficient variant}
\label{finite:sec:formal}

This section is stated over the \emph{formal-coefficient ring} $\Af_{n,\Gamma}=\C[[z_1,\ldots,z_n]]((t^\Gamma))$ of \cref{finite:sec:rings}(2). Nothing in it is transferred to or from $\HH_\Gamma(D)$, except through the explicit Taylor homomorphism \eqref{finite:eq:taylorbridge}, which is used once, in \cref{finite:thm:monadcount}, and only in the direction stated there.

\subsection{The warning example that fixes the correct algebra}
\label{finite:sec:warning}

\begin{example}[{Why $K[[z]]$ is the wrong ring, and why one well-ordered support is indispensable}]
\label{finite:ex:unit}
In $K[[z]]$ the element $F(z)=z^2-t$ is a \emph{unit}: its constant coefficient $-t$ is nonzero in the field $K$. Consequently $K[[z]]/(z^2-t)=0$: the local ring at the ordinary origin sees no quotient and no displaced roots at all. Its inverse is
\begin{equation}\label{finite:eq:badinverse}
 (z^2-t)^{-1}=-\sum_{r\geq0}t^{-r-1}z^{2r},
\end{equation}
whose exponents $-1,-2,-3,\ldots$ are \emph{not} well ordered, so \eqref{finite:eq:badinverse} does not lie in $\Af_{1,\Gamma}$. In $\Af_{1,\Gamma}$ --- and equally in $\HH_\Gamma(D)$ --- the quotient by $(z^2-t)$ has dimension two, with maximal ideals corresponding to the two infinitesimal zeros $\pm t^{1/2}$.
\end{example}

Two consequences deserve to be stated plainly. First, the theorems of this report are \emph{not} base-change assertions about $K[[z]]$: the natural map $\Af_{n,\Gamma}\hookrightarrow K[[z]]$ is injective, but an arbitrary family of coefficients in $K[[z]]$ need not have a common well-ordered support, and uniform support is exactly what \eqref{finite:eq:badinverse} violates. Second, this is why one must not expand at the ordinary centre: only after locating an \emph{actual} zero does one pass to its formal local ring, which is the content of \eqref{finite:eq:localiso}.

\subsection{Evaluation, translation, and division over \texorpdfstring{$\Af_{n,\Gamma}$}{the formal coefficient ring}}
\label{finite:sec:formaleval}

For $H=\sum_\gamma h_\gamma(z)t^\gamma\in\Af_{n,\Gamma}$ with $h_\gamma=\sum_\alpha h_{\gamma,\alpha}z^\alpha$ and $a\in\mm_K^n$, define $H(a)=\sum_{\gamma,\alpha}h_{\gamma,\alpha}t^\gamma a^\alpha$. The complex formal series $h_\gamma$ need not converge on any ordinary neighbourhood; the evaluation is by strong Hahn summability, not by ordinary analytic convergence.

\begin{lemma}[Evaluation, translation, division by coordinates]
\label{finite:lem:translation}
For every $a\in\mm_K^n$, evaluation is a $K$-algebra homomorphism $\Af_{n,\Gamma}\to K$; translation $T_aH(z)=H(z+a)=\sum_\beta(a^\beta/\beta!)\partial^\beta H(z)$ is a well-defined $K$-algebra automorphism with inverse $T_{-a}$; both have support in $\supp H+E_a^*$, $E_a=\bigcup_j\supp a_j$. Moreover $H(z)-H(a)=\sum_j(z_j-a_j)H_j(z)$ for suitable $H_j\in\Af_{n,\Gamma}$, so that $\ker(\mathrm{ev}_a)=(z_1-a_1,\ldots,z_n-a_n)$. For every $q\geq1$, Taylor expansion at $a$ induces $\Af_{n,\Gamma}/\mathfrak n_a^q\simeq K[u]/(u)^q$, where $\mathfrak n_a=\ker(\mathrm{ev}_a)$, and the same after localizing at $\mathfrak n_a$.
\end{lemma}

\begin{proof}
All nonzero coordinate supports of $a$ are positive, so every displayed term lies in $\supp H+E_a^*$; at each exponent \cref{finite:lem:support} bounds the possible powers and the contributing input coefficients, giving strong summability, and the formal binomial identities give multiplicativity and $T_aT_b=T_{a+b}$. For division, a complex formal series with zero constant term may be divided by the coordinate ideal by assigning each nonconstant monomial to its first coordinate of positive exponent and removing one occurrence; these are complex-linear operators changing no Hahn exponents, and applying $T_{-a}$ transports the identity to $a$. For the jets, use $T_a$ to reduce to $a=0$; truncating each coefficient to total $z$-degree $<q$ gives the polynomial ring, and a datum in the kernel has only monomials of total degree $\geq q$, each of which may be written as a degree-$q$ monomial times an element of $\Af_{n,\Gamma}$, without enlarging support. An element outside $\mathfrak n_a$ has a nonzero constant Taylor coefficient and is invertible in the truncation, so localization changes nothing.
\end{proof}

\begin{theorem}[Uniform-support division and conservation of colength over $\Af_{n,\Gamma}$]
\label{finite:thm:formaldivision}
Let $f_1,\ldots,f_n\in\C[[z]]$ with $f_i(0)=0$ and $\mu=\dim_\C\C[[z]]/(f)<\infty$, let $F_i=f_i+E_i\in\Af_{n,\Gamma}$ with $\supp E_i\subseteq\Gamma_{>0}$, $E=\bigcup_i\supp E_i$, $M=E^*$, and let $b_1,\ldots,b_\mu$ be polynomial representatives of a basis of $\C[[z]]/(f)$. Then every $H\in\Af_{n,\Gamma}$ has
\begin{equation}\label{finite:eq:formaldivision}
 H=\sum_{\ell=1}^{\mu}c_\ell b_\ell+\sum_{j=1}^{n}F_jQ_j,
 \qquad c_\ell\in K,\quad Q_j\in\Af_{n,\Gamma},
\end{equation}
with the remainder unique, one choice of quotients given by $h_FH$, and every remainder coefficient and quotient supported in $\supp H+M$. Hence $\Af_{n,\Gamma}/(F)\simeq K^\mu$ and $\Af^+_{n,\Gamma}/(F)\simeq(\Kc)^\mu$ as modules, on the displayed basis. The Koszul complex of $F$ over $\Af_{n,\Gamma}$ is exact in positive degrees, with $Z=\partial h_FZ$ for every positive-degree cycle. Enlarging $\Gamma$ to $\Gamma'$ gives a natural algebra isomorphism $\bigl(\Af_{n,\Gamma}/(F)\bigr)\otimes_K K'\simeq\Af_{n,\Gamma'}/(F)$, so no colength is lost or gained on enlarging the workspace.
\end{theorem}

\begin{proof}
$\C[[z]]$ is a regular local ring of dimension $n$, hence Cohen--Macaulay; the finite-colength hypothesis makes $f_1,\ldots,f_n$ a system of parameters and therefore a regular sequence, so its Koszul complex is exact in positive degrees \cite{finite:StacksKoszul}. Some power of the maximal ideal lies in $(f)$, so polynomial representatives exist. Take a contraction by \cref{finite:constr:h} and extend coefficientwise; these maps are $K$-linear and do not enlarge support. Now \cref{finite:thm:hpl}, or equivalently the conjugation \eqref{finite:eq:conjugation}, applies verbatim, and the argument of \cref{finite:thm:finitefree} gives existence, uniqueness of the remainder, the support bound, and freeness over $\Kc$ in the nonnegative-support version. For the base change, the same coefficientwise contraction over the larger group has unchanged formulas on inputs from the smaller group, and both quotients have the same basis with the same structure constants.
\end{proof}

\begin{theorem}[Uniform valuation localization of the zeros]
\label{finite:thm:localization}
Assume $E\neq\varnothing$ and put $\lambda=\min E>0$. Let $r\geq1$ be an integer with $z_j^r=0$ in $\C[[z]]/(f)$ for every $j$; one may always take $r=\mu$. Then every common infinitesimal zero $a$ of $F$ satisfies
\begin{equation}\label{finite:eq:rootbound}
 v(a_j)\geq\lambda/r\qquad(1\leq j\leq n).
\end{equation}
If $E=\varnothing$ the only zero is the origin.
\end{theorem}

\begin{proof}
Every entry of $M_j^r$ has support in $M$ and zero coefficient at exponent zero, its reduction being zero; so every nonzero entry has valuation at least $\lambda$. The coordinate $a_j$ is an eigenvalue of $M_j$; choose an eigenvector and a component of minimal valuation. In that component of $M_j^rw=a_j^rw$ the left side has valuation at least $\lambda$ plus the minimal component valuation, while cancellation can only raise valuation; hence $r\,v(a_j)\geq\lambda$. A local algebra of dimension $\mu$ has $\mu$th power of its maximal ideal zero, by Nakayama and strict decrease of the successive nonzero powers, which permits $r=\mu$. When $E=\varnothing$ all $M_j$ are the original nilpotent complex matrices.
\end{proof}

The exponent $\lambda/r$ is a bound, not an assertion that the roots have Puiseux expansions in one chosen parameter: an arbitrary-rank group may contain additional independent and infinitely smaller scales.

\subsection{Finite fibres of the monad map}
\label{finite:sec:fibres}

\Cref{finite:thm:characters,finite:thm:length} hold verbatim over $\Af_{n,\Gamma}$, with \cref{finite:lem:translation} in place of \cref{finite:lem:shifted}: the maximal ideals of $\Af_{n,\Gamma}/(F)$ are in bijection with the common zeros of $F$ in $\mm_K^n$, the character at $a$ is $[H]\mapsto H(a)$, no zero in the full surcomplex monad lies outside $K^n$, the local factor at $a$ is $K[[u]]/(T_aF_1,\ldots,T_aF_n)$, and $\sum_a\mu_a(F)=\mu$. The next theorem is the statement about the \emph{map} that this makes possible, and it is not a statement about the zero set of one system.

\begin{theorem}[Finite infinitesimal fibres]
\label{finite:thm:fibres}
In a surcomplex realization, the evaluated map $F:\mm_{\SC}^{\,n}\to\mm_{\SC}^{\,n}$ is surjective, every fibre is a finite set, and
\begin{equation}\label{finite:eq:fibres}
 \sum_{a:F(a)=b}\mu_a(F-b)=\mu
 \qquad\text{for every }b\in\mm_{\SC}^{\,n}.
\end{equation}
The finite algebra of each fibre is obtained from \cref{finite:thm:formaldivision} after enlarging the workspace to contain $b$.
\end{theorem}

\begin{proof}
Every monomial of $f_i$ has positive valuation at an infinitesimal argument since $f_i(0)=0$, and every evaluated perturbation term also has positive valuation, so $F$ maps the monad into itself. For a target $b$, the tuple $F-b$ has the same leading system and still has positive perturbation support; apply conservation. The total is the positive integer $\mu$, so the fibre is nonempty as well as finite.
\end{proof}

\begin{corollary}[A bijection of the whole monad, with fine-analytic inverse]
\label{finite:cor:inverse}
If $\det Df(0)\neq0$ --- equivalently $\mu=1$ --- then $F$ is a bijection of the full infinitesimal monad onto itself, and its inverse is fine-analytic at every point of that monad.
\end{corollary}

\begin{proof}
By \cref{finite:thm:fibres} each target has exactly one preimage, of multiplicity one. To prove analyticity, we construct its inverse with a common Hahn support, rather than infer it from a formal inverse over an unrestricted $K[[w]]$.

Let $g\in\C[[w]]^n$ be the ordinary formal inverse of $f$, with $g(0)=0$. Seek
\[
 H(w)=g(w)+\sum_{\nu\in M\setminus\{0\}}H_\nu(w)t^\nu,
 \qquad H_\nu\in\C[[w]]^n,\qquad F(H(w))=w.
\]
Expand each coefficient of $F$ formally about $g(w)$. This substitution is legitimate because $g(0)=0$. At a positive exponent $\nu$, the sole term involving the unknown $H_\nu$ is $Df(g(w))H_\nu$; all other terms involve an input error or at least two positive corrections, and so depend only on lower exponents. The matrix $Df(g(w))$ has invertible constant part $Df(0)$ and is invertible over $\C[[w]]$. Multiplying the negative known term by its inverse recursively defines $H_\nu$. The support lemma gives finitely many contributing words at each exponent, so $H\in(\Af^+_{n,\Gamma})^n$ with support in $M$. In particular this is a uniform-support construction, even when the ordinary formal coefficients diverge.

Evaluation of the composition is compatible with this recursion: at every prescribed Hahn exponent, and then at each ordinary formal degree, only finitely many positive corrections and Taylor terms contribute. Hence $F(H(b))=b$ for every infinitesimal $b$, in a workspace containing it, and the established uniqueness identifies $H$ with the inverse map. For each such $b$, \cref{finite:lem:translation} gives the exact strongly summable Taylor identity for $H(b+u)$ at every infinitesimal displacement $u$. This is fine analyticity in the support-summation sense used here. It also gives the usual fine derivative: the terms of degree at least two split into finitely many $u_ju_kR_{jk}(u)$, whose common nonnegative Hahn support makes their values finite. Thus the remainder is bounded by $n^2\omega\|u\|_\infty^2$, and its difference quotient tends to zero in the fine epsilon--delta sense as $u$ tends to zero. No convergence of Taylor partial sums in the fine topology is used.
\end{proof}

The domain here is an \emph{entire} monad, not an unspecified sufficiently small ball, and arbitrarily divergent ordinary formal coefficient functions are permitted provided the uniform Hahn support condition holds. No fine-topological convergence of a sequence of partial sums is claimed anywhere in this argument, and the proposition does not turn arbitrary fine-local analytic functions into uniformly supported data on a monad.

\subsection{Conservation monad by monad over an ordinary base}
\label{finite:sec:monadcount}

\begin{theorem}[Coherent conservation of isolated intersections]
\label{finite:thm:monadcount}
Let $U\subseteq\C^n$ be ordinary open and let $F=f+E$ be an $n$-tuple in $\HH_\Gamma(U)$ with $f:U\to\C^n$ ordinary holomorphic and $E$ of strictly positive support. Let $c\in U$ be a common zero of $f$ whose Taylor ideal has finite colength $\mu_c$. Then, with $\mu(c)=c+\mm_{\SC}^{\,n}$,
\begin{equation}\label{finite:eq:monadcount}
 \sum_{a\in\mu(c),\ F\sh(a)=0}\mu_a(F)=\mu_c .
\end{equation}
If $f(c)\neq0$ there is no common zero in $\mu(c)$. Consequently, if $\Omega\Subset U$ is bounded and the common zero set of $f$ in $\overline\Omega$ consists of finitely many isolated zeros $c_1,\ldots,c_s$ in $\Omega$ with no boundary zeros, then $F$ has finitely many common zeros in $\Omega\sh$ and their total local multiplicity is $\sum_{k}\mu_{c_k}$.
\end{theorem}

\begin{proof}
Apply the Taylor bridge \eqref{finite:eq:taylorbridge} at $c$: it sends each coefficient function to an element of $\C[[u]]$ and does not enlarge the Hahn support, so the image tuple is exactly of the form required by \cref{finite:thm:formaldivision} over $\Af_{n,\Gamma}$, with leading colength $\mu_c$. Its evaluation at $u\in\mm_{\SC}^{\,n}$ is the original coherent evaluation at $c+u$, whence \eqref{finite:eq:monadcount}. If some $f_i(c)\neq0$ then that coordinate of $F\sh(c+u)$ has the same nonzero standard part and cannot vanish. For the last assertion, the standard part of any common zero of $F$ must be a common zero of $f$, so add the counts over the finitely many possible monads.
\end{proof}

The finiteness and separation assumptions in the last assertion concern the \emph{ordinary} base; they do not assert compactness of $\Omega\sh$ in the fine topology. Note also the direction of the bridge: a coherent system on $U$ is transported to the formal ring, and no theorem is carried back.

%======================================================================
\section{Spectral invariants without choosing root branches}
\label{finite:sec:spectral}

\begin{theorem}[Characteristic polynomials, traces and norms]
\label{finite:thm:spectral}
Over $\HH_\Gamma(D)$, for every $G$,
\begin{align}
 \det(TI-M_G)&=\prod_{F\sh(\zeta)=0}\bigl(T-G\sh(\zeta)\bigr)^{\mu_\zeta},
 \label{finite:eq:characteristic}\\
 \Tr(M_G)&=\sum_{F\sh(\zeta)=0}\mu_\zeta\,G\sh(\zeta),
 \label{finite:eq:trace}\\
 \det(M_G)&=\prod_{F\sh(\zeta)=0}G\sh(\zeta)^{\mu_\zeta}.
 \label{finite:eq:norm}
\end{align}
The entries of $M_G$ have support in $\supp G+M$, and lie in $\Kc$ when $\supp G\subseteq\Gamma_{\geq0}$. In particular, for $\xi\in K^n$, $\det\bigl(TI-\sum_j\xi_jM_j\bigr)=\prod_\zeta(T-\xi\cdot\zeta)^{\mu_\zeta}$. The identical statements hold over $\Af_{n,\Gamma}$ for the systems of \cref{finite:thm:formaldivision}.
\end{theorem}

\begin{proof}
The support assertion is \cref{finite:thm:finitefree} applied to each $Ge_a$. In the local factor $\Alg_\zeta$ the element $[G]-G\sh(\zeta)$ lies in the nilpotent maximal ideal, so its multiplication operator is nilpotent; hence multiplication by $G$ on that factor has the single eigenvalue $G\sh(\zeta)$ with algebraic multiplicity $\dim_K\Alg_\zeta=\mu_\zeta$ by \cref{finite:thm:length}. Take characteristic polynomial, trace and determinant of the block decomposition.
\end{proof}

Individual roots can acquire fractional exponents and can exchange under changes of auxiliary description. The matrices above have no such labelling problem: their entries are produced directly by strong-summation division, and their characteristic coefficients are invariant under any permutation of the roots. This is particularly useful over high-rank groups, where ordering roots by a single real-valued ``size'' would discard information. Thus the total moments of a cluster are computable by finite matrix arithmetic without expanding or even individually locating its roots.

\begin{proposition}[A bounded list of separating linear forms]
\label{finite:prop:separation}
There is an integer $0\leq N\leq(n-1)\binom{\mu}{2}$ such that $\ell_N(z)=z_1+Nz_2+\cdots+N^{n-1}z_n$ takes distinct values at all distinct zeros of $F$. For $n=1$ take $\ell_0(z)=z_1$.
\end{proposition}

\begin{proof}
For distinct roots $a,b$ the polynomial $(a_1-b_1)+X(a_2-b_2)+\cdots+X^{n-1}(a_n-b_n)$ is nonzero of degree $<n$, so it has at most $n-1$ roots in $K$ and hence at most $n-1$ ordinary integer roots. There are at most $\binom\mu2$ pairs, so among $(n-1)\binom\mu2+1$ integer candidates at least one separates every pair.
\end{proof}

When every zero is simple a separating form produces the idempotent projector $e_a=\prod_{b\neq a}(\ell_N-\ell_N(b))/(\ell_N(a)-\ell_N(b))$, with $a_j=\Tr(M_{z_j}M_{e_a})$; in nonreduced fibres one uses the Chinese remainder projectors for the powers $(T-\ell_N(a))^{\mu_a}$ instead. These are algebraic extraction formulas, not a claim that arbitrary infinite input coefficients have an effective finite encoding.

%======================================================================
\section{Coefficientwise residues and perfect duality}
\label{finite:sec:residues}

\subsection{What a contour is here, and what it is not}
\label{finite:sec:contour}

For the ordinary isolated complete intersection $f$, choose a sufficiently small admissible local residue cycle
\[
 C_f=\{z:|f_1(z)|=\rho_1,\ldots,|f_n(z)|=\rho_n\}
\]
surrounding the origin and lying compactly inside $D$, with regular positive radii, oriented by $\dd\arg f_1\wedge\cdots\wedge\dd\arg f_n>0$, and taken inside a neighbourhood on which the map germ $f$ is finite. The cycle can have several sheets and is not assumed to be a single embedded torus. It is attached to the \emph{leading} equations $f$ and is never moved with the perturbation. Its ordinary functional is
\begin{equation}\label{finite:eq:ordinaryresidue}
 \mathcal R_f(g)=\frac1{(2\pi i)^n}\int_{C_f}
 \frac{g\,\dd z_1\wedge\cdots\wedge\dd z_n}{f_1\cdots f_n}.
\end{equation}
We use the classical facts that $\mathcal R_f$ annihilates $(f)$, that $(g,g')\mapsto\mathcal R_f(gg')$ is a perfect pairing on $A_0$, that it is compatible with a change of generators, that at a simple zero it is evaluation divided by the Jacobian, and that it depends holomorphically on finitely many ordinary perturbation parameters when the whole local cluster is enclosed; also the ordinary trace identity
\begin{equation}\label{finite:eq:ordinarytrace}
 \mathcal R_f(g\,J_f)=\Tr_{A_0}(m_g).
\end{equation}
These are classical local-residue and complete-intersection duality facts, not new surcomplex assertions \cite{finite:GriffithsHarris,finite:Hartshorne,finite:Kunz,finite:CattaniDickenstein,finite:NayakSastry}. The multivariate residue literature contains many further global results, for example \cite{finite:CDS}, which are not being claimed here.

\begin{definition}[Hahn local-cluster residue]
\label{finite:def:residue}
For $F=f+E$ expand, near the \emph{fixed} ordinary cycle of the leading system,
\begin{equation}\label{finite:eq:residueexpansion}
 \frac{G}{F_1\cdots F_n}
 =\sum_{k\in\N^n}(-1)^{|k|}
 \frac{G\,E_1^{k_1}\cdots E_n^{k_n}}{f_1^{k_1+1}\cdots f_n^{k_n+1}},
\end{equation}
every ordinary denominator being nonzero on $C_f$, and define
\begin{equation}\label{finite:eq:Hresidue}
 \Res_F(G)=\frac1{(2\pi i)^n}\HInt_{C_f}
 \frac{G\,\dd z_1\wedge\cdots\wedge\dd z_n}{F_1\cdots F_n}
 :=\sum_{k\in\N^n}(-1)^{|k|}\,
 \mathcal R^{(k)}\bigl(G\,E_1^{k_1}\cdots E_n^{k_n}\bigr),
\end{equation}
where $\mathcal R^{(k)}$ is the ordinary local residue with denominator $f_1^{k_1+1}\cdots f_n^{k_n+1}$, applied to every Hahn coefficient separately.
\end{definition}

It must be said emphatically, once: \emph{this series is the definition}. The superscript $H$ records only which operation is being performed --- apply the ordinary complex contour integral to each ordinary coefficient of \eqref{finite:eq:residueexpansion} and reassemble the Hahn series. It does not denote a fine-topological Riemann integral over a surcomplex parametrized torus, and it does not integrate along a fine-continuous path in $\SC$; all contours here are ordinary compact complex cycles. Supplying a genuine contour that represents this functional is a separate question, treated in the companion contours report. The sign $(-1)^{|k|}$ and the factor $(2\pi i)^{-n}$ are the normalization, chosen so that $\Res_F(J_F)=\mu$ exactly and $\Res_F(G)=\sum_\zeta G\sh(\zeta)/J_F\sh(\zeta)$ at simple zeros; the normalization is fixed by three data, namely the \emph{ordered} tuple $F_1,\ldots,F_n$, the volume form $\dd z_1\wedge\cdots\wedge\dd z_n$, and the orientation convention. No numerical smallness condition is used in \eqref{finite:eq:residueexpansion}: its correction has strictly positive Hahn support.

\begin{proposition}[The residue descends to the finite algebra]
\label{finite:prop:resdescend}
Over $\HH_\Gamma(D)$, $\Res_F$ is $K$-linear, satisfies $\supp\Res_F(G)\subseteq\supp G+M$ and $v\Res_F(G)\geq v(G)$, is independent of the admissible small leading cycle, and annihilates $(F_1,\ldots,F_n)$. It therefore induces a $K$-linear functional on $\Alg$ and a $\Kc$-linear functional $\Alg^+\to\Kc$, whose reduction is the ordinary residue $\mathcal R_f$.
\end{proposition}

\begin{proof}
The product expansions have support in $\supp G+M$, and by \cref{finite:lem:support} a fixed exponent leaves only finitely many words in the perturbation coefficients, hence finitely many ordinary meromorphic integrands, whose integrals are complex scalars; this gives strong summability, the support bound, linearity, and the reduction statement. Independence of the cycle holds coefficientwise, being ordinary residue invariance for those integrands.

For ideal annihilation, two arguments are available. Directly: for $G=F_jH$, cancel $F_j$ against its strong inverse before expanding; each remaining coefficient is an ordinary holomorphic numerator over powers of the other $f_i$, whose local residue on $C_f$ vanishes. Alternatively, and without presuming any rule for moving surcomplex contours: fix one output exponent of $\Res_F(F_jG)$; by \cref{finite:lem:witness} only finitely many perturbation coefficient functions and words contribute; replace them by independent ordinary parameters; for sufficiently small complex values of those parameters the original cycle encloses the deformed local cluster with no denominator zero on it, and the ordinary cluster residue of a numerator divisible by the $j$th defining equation is zero; its Taylor identity in the parameters is therefore zero, and assigning the original positive Hahn monomials proves that the specified coefficient vanishes.
\end{proof}

\begin{corollary}[Explicit coefficient extraction for coordinate-power leading systems]
\label{finite:cor:coefficientresidue}
If $F_i=z_i^{m_i}-A_i(z)$ with $m_i\geq1$ and $\supp A_i\subseteq\Gamma_{>0}$, so that the ordinary basis is $z^\alpha$ with $0\leq\alpha_i<m_i$ and $\mu=\prod_im_i$, then for every coherent numerator $G$
\begin{equation}\label{finite:eq:coefficientresidue}
 \Res_F(G)=\sum_{r\in\N^n}
 \bigl[z_1^{m_1(r_1+1)-1}\cdots z_n^{m_n(r_n+1)-1}\bigr]
 \Bigl(G(z)\prod_{i=1}^nA_i(z)^{r_i}\Bigr),
\end{equation}
coefficient extraction on the right being ordinary Taylor coefficient extraction applied to each Hahn coefficient. The whole sum is strongly summable and supported in $\supp G+(\bigcup_i\supp A_i)^*$. Equivalently, writing $F_i=z_i^{m_i}+E_i$, one has $\Res_F(G)=\sum_k(-1)^{|k|}[z_1^{m_1(k_1+1)-1}\cdots z_n^{m_n(k_n+1)-1}](G\prod_iE_i^{k_i})$.
\end{corollary}

\begin{proof}
Take an ordinary product of small circles as $C_f$ and expand
\[
 (z_i^{m_i}-A_i)^{-1}=\sum_{r_i\geq0}A_i^{r_i}z_i^{-m_i(r_i+1)}.
\]
Apply the ordinary one-variable Cauchy coefficient formula in each variable. The support proof is that of \cref{finite:prop:resdescend}; the resulting functional is independent of the radii, and, since the integrands at each exponent are finite sums of ordinary functions continuous on a product of circles, of the order of iterated integration.
\end{proof}

This extraction is often far more explicit than constructing a full Koszul contraction, and together with \cref{finite:cor:residuereconstruction} it computes the entire finite algebra. Let it be said once, however, that the residue is generally \emph{not} merely the top coefficient of the normal form. That simplification holds for certain polynomial systems, including \cref{finite:sec:collision} below, but not for an arbitrary holomorphic error; the theorems use the full, explicitly invertible Gram matrix.

\subsection{Perfect duality over the valuation ring}
\label{finite:sec:duality}

\begin{theorem}[Perfect integral residue duality]
\label{finite:thm:duality}
Over $\HH_\Gamma(D)$ and under \cref{finite:hyp:ordinary}, the bilinear form
\begin{equation}\label{finite:eq:pairing}
 \Alg^+\times\Alg^+\longrightarrow\Kc,
 \qquad (a,b)\longmapsto\Res_F(ab)
\end{equation}
is perfect: it induces an isomorphism $\Alg^+\xrightarrow{\ \sim\ }\Hom_{\Kc}(\Alg^+,\Kc)$ of $\Alg^+$-modules, so that $\Alg^+$ is a finite free Frobenius $\Kc$-algebra. The analogous pairing over $K$ is perfect as well, \emph{including} when the zero set is nonreduced. Relative to the fixed basis, the Gram matrix $\Gram_F=(\Res_F(e_ae_b))_{a,b=1}^\mu$ and its inverse both have entries in $\Kc$ with support in $M$, and $\det\Gram_F\in(\Kc)^\times$.
\end{theorem}

\begin{proof}
All entries lie in $\Kc$ by \cref{finite:prop:resdescend}, and their exponent-zero coefficients are exactly $\mathcal R_f(e_ae_b)$, because every occurrence of an $E_i$ has strictly positive support. The ordinary residue pairing on $A_0$ is perfect, so the complex matrix $\Gram_0$ is invertible and $\Gram_F=\Gram_0(1+\mathcal U)$ with $\mathcal U$ of positive support inside $M$. The matrix Neumann series of \cref{finite:lem:operator} gives $\Gram_F^{-1}=(1+\mathcal U)^{-1}\Gram_0^{-1}$, still supported in $M$ because $M+M=M$. In the finite free basis of \cref{finite:thm:finitefree}, invertibility of $\Gram_F$ is exactly perfection of \eqref{finite:eq:pairing}, and $\Alg^+$-linearity is $\Res_F((ca)b)=\Res_F(a(cb))$.
\end{proof}

Thus degenerating or colliding individual roots do not make the cluster pairing singular. The discriminant of the zero set may vanish; the residue Gram determinant nevertheless remains a unit, as long as the leading complete intersection and the chosen basis are fixed. Most sharply: $\det\Gram_F$ has \emph{nonzero constant leading coefficient}, so no inverse discriminant is needed anywhere.

\begin{corollary}[Contour recovery of the finite algebra]
\label{finite:cor:residuereconstruction}
Let $e^1,\ldots,e^\mu$ be the residue-dual basis, $\Res_F(e_ae^b)=\delta_a^b$, whose coordinates are the entries of $\Gram_F^{-1}$, supported in $M$, with representatives taken as $\Kc$-linear combinations of the $e_a$. Then for every $G\in\HH_\Gamma(D)$,
\begin{equation}\label{finite:eq:residuenormalform}
 \NF_F(G)_a=\Res_F(G\,e^a),
 \qquad
 (M_G)_{ab}=\Res_F(G\,e_b\,e^a);
\end{equation}
equivalently, if $m_a=\Res_F(e_aG)$ then the normal-form coordinate vector is $\Gram_F^{-1}(m_1,\ldots,m_\mu)^{\mathsf T}$, with all resulting coefficients supported in $\supp G+M$. Thus the residue functional together with a fixed ordinary basis recovers every normal form and every multiplication matrix, with no explicit choice of Koszul splitting, and this recovery is valid without simple zeros and without a reduced quotient.
\end{corollary}

\begin{proof}
Write the unique remainder of $G$ as $\sum_bc_be_b$, pair with $e^a$ and use ideal annihilation; apply the same to $Ge_b$ for the matrix formula. For the support statement use $M+M=M$.
\end{proof}

\begin{corollary}[Relative residue duality over a parameter domain]
\label{finite:cor:relative}
For the families of \cref{finite:thm:family}, coefficientwise integration on a fixed product residue cycle in the $z$-variables --- independent of $y$, because the leading equations are $z_j^{m_j}$ --- defines
\[
 \Res_{F/B}:\HH_\Gamma(B\times D)\longrightarrow\HH_\Gamma(B),
\]
which descends to the finite free family algebra. Its pairing is perfect over $\Hp_\Gamma(B)$ on the positive family algebra, and over $\HH_\Gamma(B)$ with arbitrary supports. In the rectangular monomial basis the leading Gram matrix is the \emph{constant permutation matrix}
\begin{equation}\label{finite:eq:gramleading}
 (\Gram_0)_{\alpha\beta}=
 \begin{cases}1,&\alpha_j+\beta_j=m_j-1\ \text{for every }j,\\
 0,&\text{otherwise},\end{cases}
\end{equation}
so that $\Gram=\Gram_0+\mathcal C$ with $\mathcal C$ of positive support in $M$, and $\Gram^{-1}=\sum_{k\geq0}(-\Gram_0\mathcal C)^k\Gram_0$ has support in $M$. Moreover $\Res_{F/B}(GJ_F)=\Tr(M_G)$ as an equality of Hahn-coherent functions on $B$, and every construction commutes with specialization at points of $B\sh$.
\end{corollary}

\begin{proof}
Each ordinary coefficient integral is holomorphic on $B$ by differentiation under the compact cycle, and the geometric expansion has support in $\supp G+M$ uniformly on $B$. At exponent zero every error disappears and the functional is extraction of the coefficient of $z_1^{m_1-1}\cdots z_n^{m_n-1}$, which gives \eqref{finite:eq:gramleading}; the map $\alpha\mapsto(m_1-1-\alpha_1,\ldots,m_n-1-\alpha_n)$ is an involution of $\bx$, so $\Gram_0$ is a permutation matrix with $\Gram_0^2=I$, and its leading coefficient is the nowhere-zero constant $\pm1$, which is what makes the geometric inverse valid globally on the parameter domain. At every ordinary $b\in B$, specialization gives the scalar constructions already proved, so ideal annihilation and the trace identity hold at every such $b$ and hence coefficientwise as holomorphic identities on $B$. Finally, parameter differentiation commutes with the fixed integral and with the coordinate division operators, and strong Taylor evaluation in a parameter displacement commutes with the resulting Hahn expressions.
\end{proof}

We describe this as a finite free Frobenius structure, obtained from one explicitly defined functional, rather than silently invoking a Noetherian version of relative Gorenstein duality; the proof shows directly what remains valid over the present coefficient ring.

%======================================================================
\section{Transfer, the trace--Jacobian identity, and discriminants}
\label{finite:sec:trace}

\subsection{The finite-witness transfer principle}
\label{finite:sec:transfer}

Transferring an ordinary analytic identity to infinitely many Hahn coefficients is the one place where a convergence assumption could hide. One cannot pretend that an arbitrary Hahn family is a holomorphic family in one ordinary parameter. The correct comparison is coefficientwise and uses finitely many independent parameters at a time.

\begin{lemma}[Finite-witness transfer]
\label{finite:lem:witness}
Consider expressions obtained from finitely many Hahn inputs by finitely many additions, multiplications, fixed coefficientwise complex-linear maps, derivatives, divided differences, normal forms, finite matrices and determinants, the residue expansion \eqref{finite:eq:residueexpansion}, and inverses of operators or units whose error from an invertible exponent-zero part has strictly positive support, all positive shifts lying in a common well-ordered $T\subseteq\Gamma_{>0}$. Numerators may have arbitrary well-ordered support. Then each output coefficient depends on only finitely many input coefficient functions and finitely many words in $T$.

Consequently an identity between such expressions transfers: if it holds for every sufficiently small finite-parameter ordinary holomorphic deformation of the leading system, it holds for arbitrary well-ordered positive-support Hahn errors. Concretely, replace each selected positive input coefficient by its own independent ordinary parameter times the same ordinary coefficient, prove the identity formally in those parameters, and then assign each parameter its original positive Hahn monomial.
\end{lemma}

\begin{proof}
Expand every inverse as in \cref{finite:lem:operator}, after factoring out its invertible exponent-zero part. At a prescribed output exponent, \cref{finite:lem:support} leaves only finitely many words and finitely many external input exponents; repeated finite products, finite matrices and coefficientwise maps do not change this. So the coefficient is a finite algebraic expression in exactly this finite collection of data. Keep the labels of the selected coefficient functions separate even when their exponents coincide, assign one independent variable to each, and record the words by ordinary multiplication of those variables. The corresponding coefficient of the formal identity in the independent parameters is zero by hypothesis; assigning them positive weights is strongly summable, again by \cref{finite:lem:support}, and gives the required zero coefficient. Doing this at every exponent proves the Hahn identity.
\end{proof}

\begin{remark}[Two supplements to \cref{finite:lem:witness}]
\label{finite:rem:transfer}
Two further points make the transfer usable rather than merely formal.
\begin{enumerate}
\item \emph{Finite-jet reduction.} For normal forms \emph{over a coordinate-power leading system}, which is the setting of \cref{finite:lem:matrix}, the total $z$-adic order estimate in the proof of that lemma reduces the computation at a fixed exponent to finitely many Taylor \emph{jets} of the finitely many contributing coefficient functions. A ``finite witness'' means finitely many ordinary coefficient \emph{functions}; the jet estimate is what additionally bounds the Taylor data needed from each.
\item \emph{Where ordinary convergence is licensed.} The ordinary operator formulas really are the Taylor formulas of the finite ordinary deformations. Work on a smaller \emph{closed} polydisk strictly inside $D$, in the Banach space of functions holomorphic inside and continuous on the closure. Each $R_i$ is bounded by Cauchy estimates and finite truncation, and the removable quotient $D_i$ is bounded by the maximum principle on the $z_i$ boundary, a bound $(m_i+1)\rho_i^{-m_i}$ sufficing. Hence $\delta h$ has norm less than one for sufficiently small ordinary parameters and its Neumann series converges in that ordinary Banach space, while the contour denominators are nonzero with normally convergent geometric expansions. This licenses the ordinary Taylor identities used \emph{before} substitution. \emph{No corresponding Banach convergence is claimed for the Hahn substitution}, and the smaller closed polydisk is a device internal to this proof: the actual Hahn operators and their output coefficient functions remain defined on the original open $D$, by \cref{finite:cor:restriction}.
\end{enumerate}
\end{remark}

\begin{lemma}[The ordinary finite-parameter comparison]
\label{finite:lem:finitedependence}
Choose finitely many ordinary holomorphic perturbation vectors $a_1,\ldots,a_r$ on $D$ and set $f_s(z)=f(z)+\sum_{\ell=1}^r s_\ell a_\ell(z)$. For sufficiently small complex parameters the local cluster near $0$ has a finite flat ordinary analytic algebra of rank $\mu$, the chosen $e_a$ remain a basis, and the local analytic algebra $\C\{z,s\}/(f_s)$ is finite over $\C\{s\}$ near $(0,0)$ and free of rank $\mu$ on that basis. The Taylor series in $s$ of its multiplication matrices and of its residues on the fixed leading cycle agree with the matrices and residues obtained from the formal contraction formulas over $\C[[s_1,\ldots,s_r]]$.
\end{lemma}

\begin{proof}
Finiteness follows from the analytic finite-mapping theorem, the special fibre being zero-dimensional. The total local ring before quotienting is regular; the $n$ equations have height $n$, so their quotient is Cohen--Macaulay of dimension $r$, the parameters form a system of parameters and hence a regular sequence, and the resulting finite module over the regular parameter ring has full depth and is free, of rank the dimension $\mu$ of the special fibre. Nakayama's lemma lifts the chosen special-fibre basis to a small parameter neighbourhood. These are ordinary finite-map and complete-intersection facts \cite{finite:GunningRossi,finite:Demailly,finite:Kunz}. Consequently the analytic normal-form coefficients and multiplication matrices have Taylor series in $s$; the formal Koszul contraction gives a free quotient in the same basis over $\C[[s]]$, its inverses now justified by the ordinary total-degree filtration, and uniqueness of representatives in a basis identifies the two multiplication laws. If the analytic comparison was made on a smaller $z$-neighbourhood, \cref{finite:cor:restriction} identifies it with the formal matrices constructed on $D$. Residue integrands on the fixed leading cycle have uniformly convergent geometric expansions for sufficiently small ordinary $s$, and differentiation under the ordinary integral identifies their Taylor coefficients with \eqref{finite:eq:residueexpansion}. This is an ordinary finite-parameter comparison only; no convergence assertion is made after replacing $s$ by Hahn monomials.
\end{proof}

\subsection{Trace equals residue against the Jacobian}
\label{finite:sec:tracejac}

\begin{theorem}[Trace--Jacobian identity and simple-zero residues]
\label{finite:thm:trace}
Over $\HH_\Gamma(D)$ and under \cref{finite:hyp:ordinary}, for every $G\in\HH_\Gamma(D)$,
\begin{equation}\label{finite:eq:tracejac}
 \Res_F(G\,J_F)=\Tr(M_G)=\sum_{F\sh(\zeta)=0}\mu_\zeta\,G\sh(\zeta).
\end{equation}
In particular $\Res_F(J_F)=\mu$ \emph{exactly}, with no positive-support correction. If all zeros are simple then
\begin{equation}\label{finite:eq:simpleresidue}
 \Res_F(G)=\sum_{F\sh(\zeta)=0}\frac{G\sh(\zeta)}{J_F\sh(\zeta)} .
\end{equation}
The identity holds over $\HH_\Gamma(B)$ in the parameter version of \cref{finite:cor:relative}.
\end{theorem}

\begin{proof}
Fix an exponent in the proposed equality $\Res_F(GJ_F)=\Tr(M_G)$. The residue expansion, the differentiation and determinant defining $J_F$, and the normal-form expansion defining $M_G$ all satisfy the hypotheses of \cref{finite:lem:witness}. Select the finitely many ordinary coefficient functions contributing to that exponent and replace their positive shifts by independent parameters. By \cref{finite:lem:finitedependence} the formal coefficients on both sides agree with Taylor coefficients in an ordinary finite flat family, and the classical identity \eqref{finite:eq:ordinarytrace} holds throughout that family; hence those coefficients agree. Assigning the original positive monomials proves the Hahn coefficient equality, and linearity absorbs the finitely many coefficients of $G$ contributing to the chosen output exponent. As the exponent was arbitrary, the first equality follows.

The classical input itself survives collisions: translate the target of the finite local map $f_s$ by a small regular value; on the resulting simple fibres the residue of $gJ_{f_s}$ at a point is $g$ at that point and the trace of multiplication is the sum of the same values; the finite local algebra varies flatly, and both sides vary holomorphically as the target returns to zero.

The second equality in \eqref{finite:eq:tracejac} is \eqref{finite:eq:trace}. For \eqref{finite:eq:simpleresidue}, a zero is simple exactly when its formal Jacobian determinant is nonzero (\cref{finite:thm:length}); if all zeros are simple then $[J_F]$ is a unit of the product algebra $\Alg$, and substituting a representative of $[G][J_F]^{-1}$ into the first equality, with all $\mu_\zeta=1$, gives the formula. Taking $G=1$ gives $\Res_F(J_F)=\Tr(I)=\mu$.
\end{proof}

\begin{remark}[Multiple zeros: what remains true and what must not be written]
\label{finite:rem:multiple}
The perfect pairing \cref{finite:thm:duality} and the trace--Jacobian identity \eqref{finite:eq:tracejac} do \emph{not} require simple zeros. The formula with individual terms $G\sh(\zeta)/J_F\sh(\zeta)$ does. At a multiple zero its denominator is zero, so that expression must not be used there. The functional $\Res_F$ still exists and retains the nilpotent information that the trace alone cannot recover, and \cref{finite:thm:localres} supplies its localized summands at singular intersections without ever writing a divergent quotient by the Jacobian. No unproved identification of separately localized multiple-zero residues is needed for any result in this report.
\end{remark}

\subsection{The B\'ezoutian kernel, and a divergence among the sources}
\label{finite:sec:bezout}

For the coordinate-power case there is an independent route to \eqref{finite:eq:tracejac} that also produces a reproducing kernel. Introduce a second tuple $w$ and the telescoping divided differences
\begin{equation}\label{finite:eq:bezentries}
 B_{ij}(z,w)=
 \frac{F_i(w_1,\ldots,w_{j-1},z_j,z_{j+1},\ldots,z_n)
      -F_i(w_1,\ldots,w_j,z_{j+1},\ldots,z_n)}{z_j-w_j},
\end{equation}
whose singularity on $z_j=w_j$ is removable, and put $\mathcal D_F(z,w)=\det(B_{ij}(z,w))$. Then
\begin{equation}\label{finite:eq:bezrelation}
 F_i(z)-F_i(w)=\sum_jB_{ij}(z,w)(z_j-w_j),
 \qquad \mathcal D_F(z,z)=J_F(z).
\end{equation}
The coefficients are holomorphic on $D\times D$ and their support lies in $\{0\}\cup E$ before the finite determinant is taken, so applying box normal form in both tuples gives a well-defined class in $\Alg\otimes_K\Alg$.

\begin{theorem}[Residue reproducing kernel]
\label{finite:thm:bezout}
For coordinate-power leading systems over $\HH_\Gamma(D)$, and for every $G\in\HH_\Gamma(D)$,
\begin{equation}\label{finite:eq:reproduce}
 \NF_{F,z}\Bigl(\Res_{F,w}\bigl(G(w)\,\mathcal D_F(z,w)\bigr)\Bigr)=\NF_F(G)(z);
\end{equation}
equivalently, if $b_1,\ldots,b_\mu$ is the box basis and $b^1,\ldots,b^\mu$ its residue-dual basis, then
\begin{equation}\label{finite:eq:casimir}
 [\mathcal D_F(z,w)]=\sum_{j=1}^{\mu}b_j(z)\,b^j(w)
 \qquad\text{in }\Alg\otimes_K\Alg .
\end{equation}
Both identities hold over the parameter algebra as well. Consequently the trace element of \eqref{finite:eq:Eulerelement} satisfies $\Eul_F=\sum_jb_jb^j=[J_F]$, which is \eqref{finite:eq:tracejac} again.
\end{theorem}

\begin{proof}
Consider first a sufficiently small finite ordinary deformation with simple zeros. If $a\neq b$ are two zeros, \eqref{finite:eq:bezrelation} puts the nonzero vector $a-b$ in the kernel of $B(a,b)$, so $\mathcal D_F(a,b)=0$; on the diagonal it is $J_F(a)$. The ordinary residue theorem therefore gives $\Res_{F,w}(G(w)\mathcal D_F(a,w))=G(a)$, and for a simple cluster evaluation at its points identifies the finite quotient with the product of its residue fields, so the normal-form identity holds. Adding independent small constant parameters reaches simple clusters, and holomorphic dependence extends the identity to all small parameters including zero. \Cref{finite:lem:witness} now proves \eqref{finite:eq:reproduce} for the Hahn system. In a finite free module with a perfect pairing the tensor representing the identity endomorphism is uniquely $\sum_jb_j\otimes b^j$, and \eqref{finite:eq:reproduce} says exactly that $[\mathcal D_F]$ represents that endomorphism; this is \eqref{finite:eq:casimir}. Applying the multiplication map $\Alg\otimes\Alg\to\Alg$ and using $\mathcal D_F(z,z)=J_F(z)$ gives $\Eul_F=[J_F]$.
\end{proof}

This proof isolates the genuinely new summability issue from the classical B\'ezoutian identity, which is not claimed as a new theorem of ordinary complex or polynomial analysis \cite{finite:CDS,finite:Cox}.

\begin{remark}[Where the sources diverge, and how it is resolved]
\label{finite:rem:eulerconflict}
Define, for any of the finite free algebras above with its residue-dual basis,
\begin{equation}\label{finite:eq:Eulerelement}
 \Eul_F=\sum_{a=1}^{\mu}e_a\,e^a\in\Alg,
 \qquad\text{so that}\qquad
 \Tr(M_G)=\Res_F(G\,\Eul_F)
\end{equation}
for every $G$, by expanding $Ge_a$ and summing the diagonal coefficients $\Res_F(e^aGe_a)$; the perfect pairing determines $\Eul_F$ uniquely from the trace functional, so $\Eul_F$ is basis independent, its normal-form coefficients have support in $M$, and in the coordinate-power case its exponent-zero part is $\mu\,z_1^{m_1-1}\cdots z_n^{m_n-1}$.

Two of the five sources merged here state \eqref{finite:eq:Eulerelement} and then \emph{explicitly decline} to assert $\Eul_F=[J_F]$ in this coefficient category, verifying it only in their worked examples: one writes that it does not need, and does not assert without further proof, a general identification of $\Eul_F$ with the Jacobian determinant for every analytic perturbation; the other writes that a general identity $\Eul_F=\det(\partial F_i/\partial w_j)$ is not needed to prove duality or trace conservation, and that extending the identification in complete generality is a compatibility question beyond the theorems it asserts. One of the remaining three proves it outright; the other two prove
the trace-Jacobian identity $\Res_F(GJ_F)=\Tr(M_G)$ without identifying the
trace element, and the step from the one to the other is supplied here. We print the proofs rather than averaging: \cref{finite:thm:trace} proves $\Res_F(GJ_F)=\Tr(M_G)$ for an \emph{arbitrary} isolated complete intersection over $\HH_\Gamma(D)$, by finite-witness transfer, and perfectness of the pairing then forces $\Eul_F=[J_F]$; \cref{finite:thm:bezout} proves the same for coordinate-power leading systems by an independent route which additionally yields the reproducing kernel \eqref{finite:eq:casimir}. The two declining manuscripts are not contradicted --- their theorems are true as stated --- but their caution is superseded here, and the identification is used below only in the ring and generality in which it has been proved.
\end{remark}

\subsection{Localized residues at actual surcomplex points}
\label{finite:sec:localres}

Let $e_\zeta\in\Alg$ be the idempotent equal to $1$ in the local factor $\Alg_\zeta$ and $0$ in the others, and define
\begin{equation}\label{finite:eq:localized}
 \Res_{\zeta,F}(G)=\Res_F(e_\zeta G),
\end{equation}
using any representative of $e_\zeta$; the value is independent of the choice by ideal annihilation. This definition works at singular intersections without writing a divergent quotient by the Jacobian.

\begin{theorem}[Moving-intersection residue theorem]
\label{finite:thm:localres}
There are finitely many localized residues and
\begin{equation}\label{finite:eq:localres}
 \Res_F(G)=\sum_{F\sh(\zeta)=0}\Res_{\zeta,F}(G),
 \qquad
 \Res_{\zeta,F}(G\,J_F)=\mu_\zeta\,G\sh(\zeta).
\end{equation}
Each localized pairing on $\Alg_\zeta$ is nondegenerate. If $\zeta$ is simple then $\Res_{\zeta,F}(G)=G\sh(\zeta)/J_F\sh(\zeta)$. Hence, when all roots are simple, the fixed-cycle functional equals the sum of the actual simple-intersection residues at the displaced points, and not merely a residue attached to their common standard part.
\end{theorem}

\begin{proof}
The idempotents sum to $1$, giving the first formula. Applying \eqref{finite:eq:tracejac} to $e_\zeta G$ restricts the trace to the local factor, where its value is $\mu_\zeta G\sh(\zeta)$. Distinct factors are orthogonal for the pairing since their product is zero, so perfection of the global pairing implies perfection on each factor: an element annihilating its own factor pairs to zero with every other factor and hence is zero. If $\zeta$ is simple the factor is $K$ and $J_F$ acts by the nonzero scalar $J_F\sh(\zeta)$.
\end{proof}

Via \eqref{finite:eq:localiso} this is the residue-dual functional of the ordered formal complete intersection at the displaced point, and \eqref{finite:eq:casimir} is the usual algebraic characterization of that normalization, agreeing with the formal Grothendieck convention of \cite{finite:CDS}. It should be stated that $\Res_F$ is specified by its defining formula, including the order of the defining equations and the orientation convention, and that we do not assume without proof a universal identification with every Grothendieck-residue construction over non-Noetherian Hahn analytic rings. Neither an infinitesimal ordinary-radius contour nor a choice of path between fine components is needed.

\subsection{Change of generators and of coordinates}
\label{finite:sec:transform}

\begin{corollary}[Change of generators]
\label{finite:cor:transformation}
Let $A(z)\in\Mat_n(\Hp_\Gamma(D))$ have exponent-zero coefficient matrix holomorphically invertible on $D$, and put $\widetilde F=AF$. Then $\Res_{\widetilde F}(G\det A)=\Res_F(G)$, both residues using the same volume form and the admissible leading cycle appropriate to their own leading equations.
\end{corollary}

\begin{proof}
The leading transformed equations are an equivalent isolated complete intersection. Every coefficient of the claimed equality has a finite witness, including the coefficients of the positive part of $A$; the ordinary transformation formula \cite{finite:GriffithsHarris} in the resulting finite-parameter family gives the coefficient identity, and \cref{finite:lem:witness} transfers it.
\end{proof}

\begin{theorem}[Coordinate change by an infinitesimal displacement]
\label{finite:thm:transform}
Let the leading system be coordinate powers and let $\Phi(z)=z+\eta(z)$ with every coordinate of $\eta$ of strictly positive support. Coherent Taylor composition defines $F\circ\Phi$ on the same $D$, with support controlled by the original support plus $(\bigcup_j\supp\eta_j)^*$ and coefficients still holomorphic on $D$, and
\begin{equation}\label{finite:eq:coordchange}
 \Res_{F\circ\Phi}\bigl((G\circ\Phi)\det D\Phi\bigr)=\Res_F(G).
\end{equation}
The induced map $\Phi\sh:\mm_K^n\to\mm_K^n$ is a bijection.
\end{theorem}

\begin{proof}
For sufficiently small finite ordinary parameters on a smaller polydisk, $\Phi$ is a coordinate change near the cluster, and the ordinary residue substitution rule --- or the simple-root formula with the chain rule, followed by parameter continuation --- gives \eqref{finite:eq:coordchange}; \cref{finite:lem:witness} transfers it. For bijectivity, the equations $z_i+\eta_i(z)-w_i=0$ have all degrees $m_i=1$, so by \cref{finite:cor:monadmap} each $w\in\mm_K^n$ has exactly one solution, of multiplicity one. This is bijectivity on the whole infinitesimal polydisk, not on an unspecified smaller fine neighbourhood.
\end{proof}

These laws show that the residue construction is not an artifact of a particular infinitesimal presentation of the same finite intersection. The ordinary leading coordinates remain part of the hypotheses.

\subsection{Two discriminants}
\label{finite:sec:discriminant}

Recall the notation fixed in \eqref{finite:eq:twodiscriminants}: $\Delta_F=\det M_{J_F}$ is basis independent, $\disc_F=\det\TGram_F$ with $\TGram_F=(\Tr M_{e_ae_b})_{a,b=1}^\mu$ is basis relative --- a basis change multiplies it by the square of the change-of-basis determinant.

\begin{theorem}[Discriminant--Jacobian identity and the reducedness criterion]
\label{finite:thm:discriminant}
Over $\HH_\Gamma(D)$ and in the fixed basis,
\begin{equation}\label{finite:eq:gramfactorization}
 \TGram_F=\Gram_F\,M_{J_F},
 \qquad
 \disc_F=(\det\Gram_F)\,\Delta_F,
 \qquad
 v(\disc_F)=v(\Delta_F),
\end{equation}
including the convention $v(0)=+\infty$. Moreover
\begin{equation}\label{finite:eq:deltaproduct}
 \Delta_F=\prod_{F\sh(\zeta)=0}J_F\sh(\zeta)^{\mu_\zeta},
\end{equation}
and the following are equivalent: $\Delta_F\neq0$; $\disc_F\neq0$; $\Alg$ is reduced; all zeros of $F$ in the origin monad are simple. In the reduced case $v(\Delta_F)=\sum_\zeta v(J_F\sh(\zeta))$. For the families of \cref{finite:thm:family}, $\Delta_F\in\Hp_\Gamma(B)$ and a specialized fibre is simple exactly when $\Delta_F\sh(y)\neq0$.
\end{theorem}

\begin{proof}
The $(a,b)$ entry of $\Gram_FM_{J_F}$ is $\Res_F(e_aJ_Fe_b)$, which equals $\Tr M_{e_ae_b}$ by \cref{finite:thm:trace}; taking determinants and using $\det\Gram_F\in(\Kc)^\times$ (\cref{finite:thm:duality}) gives \eqref{finite:eq:gramfactorization}. Formula \eqref{finite:eq:deltaproduct} is \eqref{finite:eq:norm} applied to $J_F$: on a local Artinian factor, multiplication by an element has determinant its residue value raised to the dimension, the nilpotent part having all eigenvalues zero. Over an algebraically closed field of characteristic zero the trace pairing of a finite commutative algebra is nondegenerate exactly when the algebra is a product of copies of the field: it is nondegenerate on such a product, and otherwise a nonzero nilpotent ideal lies in its radical, since multiplication by a nilpotent element times any other element has trace zero. So reducedness is equivalent to $\disc_F\neq0$, hence to $\Delta_F\neq0$, and by \cref{finite:thm:length} to simplicity of all zeros. In the reduced case \eqref{finite:eq:deltaproduct} gives the valuation sum. In a family, normal forms, differentiation in $z$ and determinants all commute with specialization.
\end{proof}

The criterion uses the full Hahn value, not merely its standard part: a cluster can split into distinct simple roots while $\Delta_F$ is a nonzero \emph{infinitesimal} and the leading ordinary fibre is multiple. And the contrast between the two matrices is the whole point of the word \emph{duality} here: $\Gram_F$ remains unimodular over $\Kc$ even when $\disc_F$ is infinitesimal or zero, so the residue basis is preferable to the trace basis near a multiple ordinary zero. \Cref{finite:sec:cubic} exhibits the contrast explicitly.

%======================================================================
\section{One-variable collision invariants and lifted monodromy}
\label{finite:sec:onevar}

Two complementary consequences of parameter division are recorded here. The one-variable residue identity itself is classical in the coherent setting; the point is its common-domain, collision-stable form, and the monodromy statement that follows from it.

\begin{proposition}[One-variable trace and residue formulas, stable across collisions]
\label{finite:prop:univtrace}
Let $P\in\Hp_\Gamma(U)[w]$ be monic of degree $m$, with leading ordinary polynomial $P_0$ satisfying the root-disk condition of \cref{finite:lem:ordinarydivision} and $P-P_0$ of strictly positive support. For $G\in\HH_\Gamma(U\times D_r)$ let $R_G$ be its unique remainder of degree $<m$ and put
\[
 \lambda_P(G)=\frac1{2\pi i}\Hoint_{|w|=\rho}\frac{G}{P}\,\dd w,
 \qquad r_0<\rho<r,
\]
the inverse of $P$ being obtained by geometric expansion about $P_0$. Then, for $\Alg_P=\Hp_\Gamma(U\times D_r)/(P)$ of rank $m$,
\begin{equation}\label{finite:eq:univresidue}
 \lambda_P(G)=[w^{m-1}]R_G,
 \qquad
 \Tr(M_G)=\lambda_P(P_wG)
 =\frac1{2\pi i}\Hoint_{|w|=\rho}\frac{G\,P_w}{P}\,\dd w .
\end{equation}
The pairing $(G,G')\mapsto\lambda_P(GG')$ is perfect over the parameter ring, with Gram determinant $(-1)^{m(m-1)/2}$ independently of the coefficients of $P$. All these quantities are coherent \emph{across} the ordinary discriminant locus, and specialize to the corresponding formulas for the actual surcomplex roots with their multiplicities.
\end{proposition}

\begin{proof}
$G/P-R_G/P$ is coherent holomorphic on the disk, so its contour integral vanishes; at each Hahn exponent the expansion of $R_G/P$ is an ordinary rational function with poles only at roots of $P_0$, all inside the circle, and its contour integral is its coefficient of $w^{-1}$ at infinity, which for monic $P$ and $\deg R_G<m$ is $[w^{m-1}]R_G$. For monomials $w^i,w^j$ the Gram entry is zero when $i+j<m-1$ and one when $i+j=m-1$; reversing one basis index makes the matrix triangular with unit diagonal, giving the stated determinant and perfection even when $P$ has repeated roots. For the trace, specialize an ordinary parameter and factor $P=\prod_\xi(w-\xi)^{d_\xi}$; all roots are finite with standard parts among the roots of $P_0$, polynomial division identifies the quotient with $\prod_\xi K[u]/(u^{d_\xi})$, and multiplication by $G$ has trace $\sum_\xi d_\xi G(\xi)$. On the other side $P_w/P=\sum_\xi d_\xi/(w-\xi)$; for a finite root $\xi=c+\varepsilon$ expand $(w-\xi)^{-1}=\sum_{k\geq0}\varepsilon^k(w-c)^{-k-1}$ on the ordinary circle, with support the numerator support plus $\supp(\varepsilon)^*$, so the computation is strongly summable, and ordinary Cauchy derivative formulas applied to each Hahn coefficient give the value $G(\xi)$. Both sides are coherent in the parameter, so coefficient uniqueness extends the identity to $U$ and to every finite surcomplex parameter. Division, integration and finite matrix traces all have holomorphic parameter coefficients, and none uses an inverse discriminant.
\end{proof}

\begin{theorem}[Support-preserving lift of the ordinary root cover]
\label{finite:thm:monodromy}
Let $P=P_0+E_0$ be as above and suppose $P_0(x,\cdot)$ has only simple roots for every $x\in U$, so that its ordinary root space $X_0=\{(x,c)\in U\times D_r:P_0(x,c)=0\}$ is an unramified finite analytic covering of $U$. Then there is a \emph{unique} coherent function on $X_0$ of the form
\[
 r(x,c)=c+\eta(x,c),
 \qquad \supp\eta\subseteq\supp(E_0)^*\setminus\{0\},
\]
satisfying $P(x,r(x,c))=0$. On every root chart its coefficients are holomorphic on that chart's \emph{entire} ordinary domain. Analytic continuation of the lifted roots along an ordinary loop permutes their labels in exactly the same way as continuation of the ordinary roots of $P_0$: the monodromy is unchanged.
\end{theorem}

\begin{proof}
On $X_0$ the function $P_{0,w}(x,c)$ is ordinary holomorphic and nowhere zero. Substituting $c+\eta$ into $P$, the coefficient of $\eta_\nu$ at a positive exponent $\nu$ is $P_{0,w}(x,c)\eta_\nu$, while all other terms at that exponent involve the input and previously determined coefficients: the nonlinear terms contain at least two positive corrections, and a correction from $E_0$ is multiplied by a positive input exponent. Recursively divide the negative of their sum by $P_{0,w}(x,c)$. The support-word argument of \cref{finite:lem:support} constructs $\eta$ on $\supp(E_0)^*$ with finite contributions at each exponent. Every operation uses ordinary holomorphic functions on $X_0$ and division by the globally zero-free derivative, so the coefficients glue on the root space itself and no globally labelled choice of roots over $U$ is needed. Uniqueness follows from the least positive exponent at which two candidate corrections could differ. Continuing a chart expression along a loop gives the expression on the continued ordinary sheet, and uniqueness excludes any additional permutation.
\end{proof}

The loops here live in the \emph{ordinary} parameter base and continuation concerns ordinary holomorphic coefficients; the theorem does not postulate nonconstant fine-continuous real-parameter loops in the surcomplex class. New collisions whose parameter standard part lies on the ordinary discriminant are not excluded either: the theorem deliberately works over the complement of that locus, while the finite algebra and the residue pairing remain meaningful across such collisions on a larger parameter chart to which preparation applies. \Cref{finite:sec:cubic} contains a worked instance.

%======================================================================
\section{Stability: two certificates, and which is sharper for what}
\label{finite:sec:stability}

Two different stability questions must be kept apart. Normal-form coefficients and cluster residues can be insensitive to a high-valuation perturbation, while individual roots are sensitive to an almost singular Jacobian. Both are proved here; and here the five sources genuinely diverge in what they are willing to certify, so both certificates are printed with a statement of which is sharper for which purpose.

All inequalities in this section concern the \emph{coefficient valuation}, never an ordinary analytic norm.

\subsection{The aggregate algebra does not lose precision}
\label{finite:sec:precision}

\begin{theorem}[Precision of normal forms, residues and spectral invariants]
\label{finite:thm:precision}
Fix $D$, $f$ and the representative basis, and let $F=f+E$ and $\widetilde F=f+\widetilde E$ be two strictly positive Hahn perturbations of the \emph{same} isolated ordinary complete intersection over $\HH_\Gamma(D)$, with
\[
 \eta=\min_i v\bigl(E_i-\widetilde E_i\bigr)>0
\]
($\eta=+\infty$ when the systems agree). Then for every $G\in\HH_\Gamma(D)$,
\begin{align}
 v\bigl(\NF_F(G)-\NF_{\widetilde F}(G)\bigr)&\geq v(G)+\eta,
 \label{finite:eq:precisionnormal}\\
 v\bigl(\Res_F(G)-\Res_{\widetilde F}(G)\bigr)&\geq v(G)+\eta,
 \label{finite:eq:precisionresidue}\\
 v\bigl(M_G^F-M_G^{\widetilde F}\bigr)&\geq v(G)+\eta.
 \label{finite:eq:precisionmatrix}
\end{align}
Witnesses chosen with the same $h$ satisfy \eqref{finite:eq:precisionnormal} as well. For $G$ of nonnegative support, every corresponding characteristic-polynomial coefficient differs by valuation at least $\eta$; so do the residue Gram matrices, their inverses, the trace Gram matrices, and the three determinants $\det\Gram$, $\Delta$, $\disc$. If the numerator is changed too, the normal forms, residues, multiplication matrices and witnesses for $(F,G)$ and $(\widetilde F,\widetilde G)$ differ by valuation at least
\begin{equation}\label{finite:eq:inputprecision}
 \min\bigl\{v(G-\widetilde G),\ v(G)+\eta\bigr\}.
\end{equation}
These estimates remain valid at higher-rank exponents where an entire initial segment of the support is infinite.
\end{theorem}

\begin{proof}
Use the same ordinary contraction for both systems. Their resolvents satisfy the exact identity
\begin{equation}\label{finite:eq:resolventdifference}
 S_F-S_{\widetilde F}=S_F(\widetilde\delta-\delta)h\,S_{\widetilde F}.
\end{equation}
Every $S$ preserves the lower support bound and $p,h$ do not lower it, while $\widetilde\delta-\delta$ raises it by at least $\eta$; applying \eqref{finite:eq:resolventdifference} to $G$ gives \eqref{finite:eq:precisionnormal} and the witness bound, and applying it to $Ge_a$ gives \eqref{finite:eq:precisionmatrix}. For residues, on the fixed leading cycle $F_i^{-1}-\widetilde F_i^{-1}=(\widetilde F_i-F_i)/(F_i\widetilde F_i)$, the inverses on the right having nonnegative Hahn support as meromorphic coefficient data near the cycle; a finite telescoping product of these identities bounds $v\bigl(G/(F_1\cdots F_n)-G/(\widetilde F_1\cdots\widetilde F_n)\bigr)$ below by $v(G)+\eta$, and coefficientwise integration cannot introduce a smaller exponent. For a nonnegative-support $G$, both multiplication matrices have entries of nonnegative valuation, and any difference between characteristic coefficients is a sum of products with at least one changed entry; the same argument covers the Gram matrices and the determinants, while for the inverse Gram matrices use $\Gram_F^{-1}-\Gram_{\widetilde F}^{-1}=\Gram_F^{-1}(\Gram_{\widetilde F}-\Gram_F)\Gram_{\widetilde F}^{-1}$ with both inverses having entries in $\Kc$. For $\Delta_F$, also note that differentiation does not lower coefficient valuation: expansion of the two Jacobian determinants gives $v(J_F-J_{\widetilde F})\geq\eta$, since all their derivative entries have nonnegative support. Combining this with the matrix estimate and its changed-input version gives $v(M^F_{J_F}-M^{\widetilde F}_{J_{\widetilde F}})\geq\eta$, and then the determinant bound. Estimate \eqref{finite:eq:inputprecision} for the stated linear constructions follows by adding and subtracting the result for $(\widetilde F,G)$.
\end{proof}

\begin{corollary}[A discriminant threshold that preserves simple zeros]
\label{finite:cor:discstable}
Suppose $F$ has only simple zeros and $\eta>v(\Delta_F)$. Then $\widetilde F$ also has exactly $\mu$ distinct simple zeros, and $\Delta_{\widetilde F}$ has the \emph{same leading Hahn term} as $\Delta_F$.
\end{corollary}

\begin{proof}
By \cref{finite:thm:precision}, $v(\Delta_{\widetilde F}-\Delta_F)\geq\eta$. The strict inequality prevents cancellation of the leading term of $\Delta_F$, so $\Delta_{\widetilde F}\neq0$; apply \cref{finite:thm:discriminant,finite:thm:length}.
\end{proof}

\begin{remark}[What this certificate deliberately does not assert]
\label{finite:rem:rootinstability}
\Cref{finite:thm:precision} does \emph{not} assert that a labelling of individual roots is Lipschitz. Roots of a multiple leading system generally involve fractional exponents, and inverse Jacobians can have negative valuation. Normal forms and the total residue nevertheless preserve the input precision, because their construction uses only positive-support inverses about ordinary invertible data. \Cref{finite:cor:discstable} therefore provides one certified geometric consequence \emph{without choosing root branches}, and it is available even when no root labelling exists. The next subsection supplies the complementary certificate, which does match roots but demands a hypothesis that \cref{finite:cor:discstable} does not.
\end{remark}

\subsection{Sharp conditioned root stability}
\label{finite:sec:rootstability}

\begin{theorem}[Sharp conditioned root stability]
\label{finite:thm:rootstability}
Let the leading system be coordinate powers, $F_i=z_i^{m_i}+E_i$ over $\HH_\Gamma(D)$, and let $a\in\mm_K^n$ be a \emph{simple} zero. Put
\begin{equation}\label{finite:eq:kappa}
 A=DF\sh(a),
 \qquad
 \kappa=\max\{0,-v(A^{-1})\},
\end{equation}
the inverse-Jacobian loss. Let $\Delta$ be a coherent vector on $D$ and let $\sigma\in\Gamma$ satisfy
\begin{equation}\label{finite:eq:threshold}
 v(\Delta)\geq\sigma>2\kappa .
\end{equation}
Then $\widetilde F=F+\Delta$ has exactly one zero $a+h$ in the ball $\mathcal B_a=\{a+u:v(u)>\kappa\}$; it is simple; and
\begin{align}
 v(h)&\geq\sigma-\kappa,\label{finite:eq:shiftbound}\\
 v\bigl(h+A^{-1}\Delta\sh(a)\bigr)&\geq2\sigma-3\kappa.\label{finite:eq:errorbound}
\end{align}
All statements hold in arbitrary divisible Hahn workspaces, with no rank-one and no sequential-completeness assumption.
\end{theorem}

\begin{proof}
Every coefficient of the Taylor expansion of $F\sh(a+u)$ has nonnegative valuation, the coherent coefficients of $F$ having nonnegative support and $a$ being infinitesimal. Write
\[
 F\sh(a+u)=Au+Q(u),\quad \ord_uQ\geq2,
 \qquad
 \Delta\sh(a+u)=\Delta\sh(a)+Bu+Q_\Delta(u),\quad \ord_uQ_\Delta\geq2,
\]
with $v(B)\geq\sigma$ and all coefficients of $Q_\Delta$ of valuation at least $\sigma$.

Set $r=\sigma-\kappa$, so $r>\kappa\geq0$, and substitute $u=t^ry$. The normalized equation is
\begin{equation}\label{finite:eq:normalizedhensel}
 t^{-r}A^{-1}\widetilde F\sh(a+t^ry)=y+b+P(y)=0,
 \qquad b=t^{-r}A^{-1}\Delta\sh(a),
\end{equation}
where $v(b)\geq0$ and every nonzero coefficient of $P$ has \emph{strictly positive} valuation: the extra linear term has valuation at least $\sigma-\kappa>0$, and a degree-$k\geq2$ term from $Q$ has valuation at least $(k-1)r-\kappa\geq r-\kappa=\sigma-2\kappa>0$, terms from $Q_\Delta$ gaining more. Equation \eqref{finite:eq:normalizedhensel} is a genuinely \emph{coherent} equation, not merely a formal one: the Taylor coefficients at $a$ have a common nonnegative support inside the original support plus the positive monoid generated by $\supp a$, multiplication by the fixed matrix $A^{-1}$ and by the powers $t^{(k-1)r}$ has well-ordered support with finite contribution, and at a fixed exponent only finitely many degrees $k$ contribute, so each coefficient function of $y$ is a polynomial.

Translate $y=-\st(b)+\nu$. The new leading system is $\nu_i$, all remaining coefficients having positive support, so the degree-one case of \cref{finite:thm:characters,finite:thm:length} gives a unique infinitesimal $\nu$, hence a finite $y$ solving \eqref{finite:eq:normalizedhensel}. This constructs a zero with $v(h)\geq r$, which is \eqref{finite:eq:shiftbound}, and places it in $\mathcal B_a$.

For uniqueness on the \emph{whole} ball, not merely on the rescaled monad used in the construction, take $u,w$ with $v(u),v(w)>\kappa$. Taylor divided differences give
\begin{equation}\label{finite:eq:injectiveball}
 \widetilde F\sh(a+u)-\widetilde F\sh(a+w)=\bigl(A+C(u,w)\bigr)(u-w),
 \qquad
 v\bigl(C(u,w)\bigr)\geq\min\{v(u),v(w),\sigma\}>\kappa,
\end{equation}
every nonlinear divided-difference term containing at least one coordinate of $u$ or $w$, and the perturbation derivative having valuation at least $\sigma$; these expansions are strongly summable. Since $v(A^{-1}C)>0$, the matrix $I+A^{-1}C$ is invertible by its geometric series, so \eqref{finite:eq:injectiveball} proves injectivity on $\mathcal B_a$ and hence uniqueness; the same estimate for the derivative at $a+h$ shows the new zero is simple. Finally the root equation gives
\[
 h+A^{-1}\Delta\sh(a)=-A^{-1}\bigl(Q(h)+Bh+Q_\Delta(h)\bigr),
\]
and $v(h)\geq r$ yields the lower bound $\min\{2r-\kappa,\ \sigma+r-\kappa,\ \sigma+2r-\kappa\}=2\sigma-3\kappa$, since $\kappa\geq0$. This is \eqref{finite:eq:errorbound}.
\end{proof}

\begin{corollary}[Stable matching of an entire simple cluster]
\label{finite:cor:matching}
Suppose all $\mu$ zeros of $F$ are simple, with $\kappa_a$ as in \eqref{finite:eq:kappa}, and suppose $v(\Delta)\geq\sigma>2\max_a\kappa_a$. Then the zeros of $F$ and of $F+\Delta$ have a unique matching, characterized by $v(\widetilde a-a)>\kappa_a$; all perturbed zeros are simple; and each matched pair satisfies \eqref{finite:eq:shiftbound}--\eqref{finite:eq:errorbound} with its own $\kappa_a$.
\end{corollary}

\begin{proof}
The balls $\mathcal B_a$ are disjoint: if two intersected, the ultrametric inequality would place the centre of the ball with the larger threshold inside the other, and applying \eqref{finite:eq:injectiveball} with $\Delta=0$ would put two distinct zeros of $F$ inside a ball on which $F$ is injective. \Cref{finite:thm:rootstability} constructs one simple perturbed root in each of the $\mu$ balls, and the total multiplicity of the perturbed system is $\mu$ by \cref{finite:thm:length}, so there are no others.
\end{proof}

\begin{proposition}[Both bounds attained, and the threshold strict]
\label{finite:prop:sharpness}
The strict inequality $\sigma>2\kappa$ cannot be weakened to $\sigma\geq2\kappa$ in any uniform simple-root matching theorem, and the two valuation bounds \eqref{finite:eq:shiftbound} and \eqref{finite:eq:errorbound} are both attained.
\end{proposition}

\begin{proof}
Let $\alpha>0$, $\tau=t^\alpha$ and $F(x)=x^2-\tau^2$, with roots $\pm\tau$. At $a=\tau$ the inverse derivative is $(2\tau)^{-1}$, so $\kappa=\alpha$. The perturbation $\Delta=\tau^2$, of valuation exactly $2\kappa$, changes $F$ to $x^2$: two simple roots merge into one double root. This proves strictness.

For $\Delta=t^\sigma$ with $\sigma>2\alpha$, the root near $\tau$ is
\[
 \tau+h=\tau\Bigl(1-\frac{t^\sigma}{\tau^2}\Bigr)^{1/2}
 =\tau-\frac{t^\sigma}{2\tau}-\frac{t^{2\sigma}}{8\tau^3}-\cdots,
\]
the binomial series being strongly summable because $\sigma-2\alpha>0$. Its displayed nonzero leading coefficients give
\[
 v(h)=\sigma-\alpha,
 \qquad
 v\Bigl(h+\frac{t^\sigma}{2\tau}\Bigr)=2\sigma-3\alpha,
\]
so both lower bounds are attained. The examples embed in any number of variables by adjoining equations $z_i=0$.
\end{proof}

\subsection{Which certificate is sharper for which purpose}
\label{finite:sec:whichcertificate}

The two subsections above certify \emph{different} things, and averaging them into a single ``discriminant or Jacobian threshold'' would lose real content. The following is the resolution used throughout this report.

\begin{itemize}
\item \textbf{Use \cref{finite:thm:rootstability,finite:cor:matching} when individual roots are to be matched}, and when the two attained bounds \eqref{finite:eq:shiftbound}--\eqref{finite:eq:errorbound} are wanted. This requires a simple zero and the conditioning number $\kappa$, and it delivers a labelled bijection of root sets together with a first-order formula $h\approx-A^{-1}\Delta\sh(a)$ with a quantified second-order error. It is proved here only for coordinate-power leading systems.
\item \textbf{Use \cref{finite:cor:discstable} when the only statement needed is that all zeros remain simple}, and when no root labelling is available. It applies to an \emph{arbitrary} isolated complete intersection, it needs no conditioning number, and its threshold $\eta>v(\Delta_F)$ is stated in an invariant of the whole cluster rather than of one root. It certifies nothing about where individual roots go, and this is deliberate: one of the sources merged here declines root matching outright on the ground that roots of a multiple leading system generally involve fractional exponents and inverse Jacobians can have negative valuation, and certifies stability instead through the discriminant, calling it one certified geometric consequence obtained without choosing root branches. That refusal is not an error; it is a different and strictly more general certificate, and it is printed above as such.
\item \textbf{The determinant-only bound is provably lossy.} One always has $0\leq\kappa_a\leq v(\det A)$, so a threshold stated in terms of $v(\det A)$ --- equivalently obtained from the adjugate --- is weaker than the inverse-Jacobian threshold of \eqref{finite:eq:threshold}, and there are worked cases in which $\kappa=\max(\alpha,\beta)$ beats a determinant-only threshold $\sigma>2(\alpha+\beta)$. This remark is recorded here, where both certificates are stated; the companion analytic-geometry report cites it rather than restating it.
\item \textbf{Aggregate versus individual.} \Cref{finite:thm:precision} and \cref{finite:thm:rootstability} are not in competition at all: the first says that the \emph{cluster} data --- normal forms, matrices, residues, Gram matrices, discriminants --- lose no valuation whatever, while the second says that \emph{individual roots} lose exactly $\kappa$. \Cref{finite:sec:collision} exhibits both numbers in one example.
\end{itemize}

%======================================================================
\section{Worked examples}
\label{finite:sec:examples}

The examples are chosen to exhibit, in concrete equations, each of the phenomena the theorems allow: collision with an unimpaired pairing; roots separated by genuinely incommensurable scales; a triple collision resolved at a scale smaller than every power of the first infinitesimal; the difference between residue duality and trace degeneration; a support with infinitely many exponents below a bound; and nonpolynomial coefficient functions.

\subsection{A coupled rank-four cluster: collision, discriminant, infinite residues}
\label{finite:sec:collision}

Let $\tau=t^\alpha$ with $\alpha>0$, let $\sigma\in\mm_K$, and over $\HH_\Gamma(D)$ take
\begin{equation}\label{finite:eq:exampleF}
 F_1=x^2-\tau y-\sigma,
 \qquad
 F_2=y^2-\tau x .
\end{equation}
The leading ordinary equations are $(x^2,y^2)$, so $\mu=4$ with basis $\bx=(1,x,y,xy)$; the coupling prevents this from being a pair of independent quadratics. In the quotient,
\[
 x^2=\tau y+\sigma,\quad y^2=\tau x,\quad
 x^2y=\tau^2x+\sigma y,\quad xy^2=\tau^2y+\tau\sigma,\quad
 x^2y^2=\tau^2xy+\tau\sigma x,
\]
so that, columns being the images of the basis elements,
\begin{equation}\label{finite:eq:examplematrices}
 M_x=\begin{pmatrix}0&\sigma&0&0\\1&0&0&\tau^2\\0&\tau&0&\sigma\\0&0&1&0\end{pmatrix},
 \qquad
 M_y=\begin{pmatrix}0&0&0&\tau\sigma\\0&0&\tau&0\\1&0&0&\tau^2\\0&1&0&0\end{pmatrix},
\end{equation}
with $M_xM_y=M_yM_x$, $M_x^2=\tau M_y+\sigma I$ and $M_y^2=\tau M_x$. Since $\tau\neq0$ in $K$, eliminating $x=y^2/\tau$ gives
\begin{equation}\label{finite:eq:quarticP}
 P(y)=y^4-\tau^3y-\sigma\tau^2,
\end{equation}
so $\Alg\simeq K[y]/(P)$; the integral basis $\bx$ is nevertheless preferable for residue duality because it avoids dividing by the infinitesimal $\tau$. The Jacobian and the basis-independent discriminant are
\begin{equation}\label{finite:eq:examplejacdisc}
 J_F=4xy-\tau^2,
 \qquad
 \Delta_F=\det\bigl(4M_xM_y-\tau^2I\bigr)=-\tau^2\bigl(27\tau^6+256\sigma^3\bigr),
\end{equation}
the quartic discriminant being $\tau^4\Delta_F$. All four points are simple unless $27\tau^6+256\sigma^3=0$, and a collision is exhibited exactly by $\sigma=-3\cdot4^{-4/3}\tau^2$: then $y_0=\tau/4^{1/3}$ is a double root of $P$, since $P(y_0)=P'(y_0)=0$ and $P''(y_0)\neq0$, while the other two roots are simple, so the local lengths are $2,1,1$ with total $4$, as \eqref{finite:eq:conservation} requires.

\paragraph{The pairing does not notice the collision.}
On the basis,
\begin{equation}\label{finite:eq:exampleresvalues}
 \Res_F(1)=\Res_F(x)=\Res_F(y)=0,\qquad \Res_F(xy)=1 .
\end{equation}
These follow from \cref{finite:cor:coefficientresidue} without assuming simple roots: expanding
\[
 \frac{x^ry^s}{F_1F_2}
 =\sum_{p,q\geq0}\frac{x^ry^s(\tau y+\sigma)^p(\tau x)^q}{x^{2p+2}y^{2q+2}},
\]
a term whose $y$-power $c$ is chosen from $(\tau y+\sigma)^p$ contributes to the coefficient of $x^{-1}y^{-1}$ only if $q=2p+1-r$, $c=4p+3-2r-s$ and $0\leq c\leq p$; for $r,s\in\{0,1\}$ the unique solution is $p=q=c=0$, $r=s=1$. Combining with the multiplication rules,
\begin{equation}\label{finite:eq:examplegram}
 \Gram_F=\begin{pmatrix}0&0&0&1\\0&0&1&0\\0&1&0&0\\1&0&0&\tau^2\end{pmatrix},
 \qquad\det\Gram_F=1,
\end{equation}
which does not involve $\sigma$ at all and remains a unit at the double-root collision where \eqref{finite:eq:examplejacdisc} vanishes --- \cref{finite:thm:duality} in a single display. Also $\Res_F(J_F)=4$, $\Tr M_x=\Tr M_y=0$ and $\Tr(M_xM_y)=3\tau^2$, consistent with \eqref{finite:eq:tracejac}.

\paragraph{A finite total assembled from infinite summands.}
At $\sigma=0$ the roots are $(0,0)$ and $(\tau\zeta^2,\tau\zeta)$ for $\zeta^3=1$; the Jacobian is $-\tau^2$ at the first and $3\tau^2$ at each of the others, so \eqref{finite:eq:simpleresidue} reads
\begin{equation}\label{finite:eq:infiniterescancel}
 \Res_F(1)=-\tau^{-2}+3\cdot\frac1{3\tau^2}=0,
 \qquad
 \Res_F(xy)=3\cdot\frac{\tau^2}{3\tau^2}=1 .
\end{equation}
The individual residues in the first sum are \emph{infinite} surcomplex numbers of valuation $-2\alpha$. Their cancellation is an exact equality in $K$, not a limiting regularization, and it is why the residue functional cannot be replaced by a procedure that discards infinite individual terms.

\paragraph{Both stability numbers in one place.}
At all four roots the least valuation of the inverse-Jacobian entries is $-\alpha$, so $\kappa=\alpha$ in \eqref{finite:eq:kappa}. Any coherent perturbation with $\sigma'>2\alpha$ therefore matches all four roots uniquely by \cref{finite:cor:matching}, each displacement having valuation at least $\sigma'-\alpha$: the individual roots lose exactly $\alpha$ of accuracy. At the same time \cref{finite:thm:precision} says that the cluster residue of a numerator of valuation zero changes only at valuation $\sigma'$ or higher: the cluster data lose nothing. The two certificates of \cref{finite:sec:whichcertificate} are visible side by side.

\subsection{A root separated by a higher Archimedean scale}
\label{finite:sec:higherrank}

Keep \eqref{finite:eq:exampleF} and now choose
\[
 \Gamma=\Q+\Q\omega,\qquad \tau=t,\qquad\sigma=t^\omega .
\]
The support $\{1,\omega\}$ is positive and well ordered, but $v(\sigma)>Nv(\tau)$ for \emph{every} ordinary positive integer $N$: the rank-four theorem requires no common Archimedean scale for the two positive valuations. The quartic \eqref{finite:eq:quarticP} then has one root much smaller than $\tau$. Setting
\[
 q=\frac{\sigma^3}{\tau^6}=t^{3\omega-6},
 \qquad
 y_*=\frac\sigma\tau\,u(q),
\]
equation \eqref{finite:eq:quarticP} becomes $u+1=qu^4$, whose unique formal branch with $u(0)=-1$ is strongly summable at this positive infinitesimal $q$ and begins
\begin{equation}\label{finite:eq:usmall}
 u(q)=-1+q-4q^2+22q^3-140q^4+969q^5+O(q^6),
\end{equation}
so that
\begin{align}
 y_*&=-\frac\sigma\tau+\frac{\sigma^4}{\tau^7}-4\frac{\sigma^7}{\tau^{13}}
 +22\frac{\sigma^{10}}{\tau^{19}}+\cdots,
 \label{finite:eq:ysmall}\\
 x_*&=\frac{\sigma^2}{\tau^3}\bigl(1-2q+9q^2-52q^3+340q^4+O(q^5)\bigr).
 \label{finite:eq:xsmall}
\end{align}
Here $O(q^N)$ means a formal power series with no terms below degree $N$, evaluated by strong summation; it is not a sequential asymptotic statement. The supports are
\[
 (\omega-1)+\N(3\omega-6)
 \quad\text{and}\quad
 (2\omega-3)+\N(3\omega-6),
\]
explicit well-ordered positive sets. For each cube root of unity $\zeta$, writing $y=\tau v$ and $e=\sigma/\tau^2$ gives $v^4-v-e=0$ with ordinary leading derivative $3$ at $v=\zeta$, so
\[
 y_\zeta=\tau\Bigl(\zeta+\frac e3-\frac{2\zeta^2e^2}9+\frac{20\zeta e^3}{81}+O(e^4)\Bigr),
 \qquad x_\zeta=\frac{y_\zeta^2}\tau .
\]
The nonzero discriminant and the rank-four count show these four branches exhaust the cluster. The leading exponents of the small branch involve $\omega$, so they cannot be encoded by a Puiseux series whose exponents are rational multiples of $v(\tau)=1$: this is precisely why the arbitrary-rank Hahn statement is more than a one-parameter power-series calculation.

\subsection{A rank-three example compared lexicographically in \texorpdfstring{$1,\omega,\omega^2$}{1, omega, omega squared}}
\label{finite:sec:quartic}

Let $a,b,c\in\mm_K$ and take, over $\HH_\Gamma(D)$,
\begin{equation}\label{finite:eq:quarticsystem}
 F_1=x^2-y-a,\qquad F_2=y^2-bx-c .
\end{equation}
The leading system $(x^2-y,\,y^2)$ has sole zero $(0,0)$ with quotient $\C[x]/(x^4)$, so $\mu=4$. It is triangular after ordering the equations as $y^2,x^2-y$, but it is not a coordinate-power leading system in the displayed coordinates; therefore \cref{finite:thm:finitefree} applies directly, while \cref{finite:thm:family} does not. Eliminating $y$,
\begin{equation}\label{finite:eq:quartic}
 P(x)=(x^2-a)^2-bx-c=x^4-2ax^2-bx+(a^2-c),
\end{equation}
so $\Alg\simeq K[x]/(P)$ with $y=x^2-a$, basis $1,x,x^2,x^3$, companion matrix $M_x$ and $M_y=M_x^2-aI$, and
\[
 J_F=4xy-b\equiv4x^3-4ax-b=P'(x).
\]
In this basis the residue is the coefficient of $x^3$ in the remainder modulo $P$:
\[
 \Res_F(g)=[x^3]\operatorname{rem}_P\bigl(g(x,x^2-a)\bigr).
\]
For coherent $g$ the right side means its coherent normal form, not polynomial long division applied to an unevaluated infinite expression. Writing $S_k=\Res_F(x^k)$,
\begin{equation}\label{finite:eq:moments}
 S_0=S_1=S_2=0,\quad S_3=1,
 \qquad
 S_{k+4}=2aS_{k+2}+bS_{k+1}-(a^2-c)S_k\quad(k\geq0),
\end{equation}
whence $S_4=0$, $S_5=2a$, $S_6=b$, $S_7=3a^2+c$, $S_8=4ab$, $S_9=4a^3+4ac+b^2$, $S_{10}=10a^2b+2bc$, and
\begin{equation}\label{finite:eq:quarticgram}
 \Gram_F=\begin{pmatrix}0&0&0&1\\0&0&1&0\\0&1&0&2a\\1&0&2a&b\end{pmatrix},
 \qquad\det\Gram_F=1,
\end{equation}
invertible at every degeneration of the root configuration. The generating function $\sum_kS_kw^k=w^3/(1-2aw^2-bw^3+(a^2-c)w^4)$ gives the closed form
\begin{equation}\label{finite:eq:momentmultinomial}
 S_k=\sum_{\substack{r,s,u\geq0\\ 3+2r+3s+4u=k}}
 \frac{(r+s+u)!}{r!\,s!\,u!}(2a)^rb^s(c-a^2)^u,
\end{equation}
and, for the coherent entire numerator $e^x$, the exact strong identity
\begin{equation}\label{finite:eq:quarticexp}
 \Res_F(e^x)=\sum_{r,s,u\geq0}
 \frac{(r+s+u)!}{r!\,s!\,u!}\,
 \frac{(2a)^rb^s(c-a^2)^u}{(3+2r+3s+4u)!},
 \qquad \st\Res_F(e^x)=\tfrac16 .
\end{equation}
Grouping by numerator degree gives $\tfrac1{3!}+\tfrac{2a}{5!}+\tfrac b{6!}+\tfrac{3a^2+c}{7!}+\tfrac{4ab}{8!}+\cdots$; this display is ordered by numerator degree, \emph{not} by increasing Hahn valuation, and \eqref{finite:eq:quarticexp} is the precise statement.

Now choose the three genuinely separated scales
\begin{equation}\label{finite:eq:threescales}
 a=t,\qquad b=t^\omega,\qquad c=t^{\omega^2},
\end{equation}
in the divisible group generated by $1,\omega,\omega^2$. Since $\disc(P)=-256a^3b^2+256a^2c^2+288ab^2c-27b^4-256c^3$ and $\det\Gram_F=1$, \cref{finite:thm:discriminant} identifies this with both $\disc_F$ and $\Delta_F$. Comparing the five exponents
\[
 3+2\omega,\quad 2+2\omega^2,\quad 1+2\omega+\omega^2,\quad 4\omega,\quad 3\omega^2
\]
lexicographically by their $\omega^2$, $\omega$ and finite parts --- \emph{not} by a single-real-parameter Puiseux approximation --- the leading term is $-256a^3b^2$, of valuation $3+2\omega$. Hence $\Delta_F\neq0$ and all four zeros are simple. Writing $s=\sqrt a$ and $r_\epsilon^2=b\epsilon s+c$ for $\epsilon=\pm1$, the substitution $x=\epsilon s+(r_\epsilon/2\epsilon s)W$ turns \eqref{finite:eq:quartic} into
\[
 \Bigl(W+\frac{r_\epsilon}{4a}W^2\Bigr)^2-1-\frac b{2\epsilon s\,r_\epsilon}W=0,
\]
whose two displayed correction parameters are infinitesimal under \eqref{finite:eq:threescales}; at their ordinary value $0$ the equation is $W^2-1=0$ with two simple roots, and the ordinary formal implicit-function expansions at $W=\pm1$ evaluate strongly, giving
\begin{equation}\label{finite:eq:quarticrootleading}
 x_{\epsilon,\delta}=\epsilon s+\frac{\delta r_\epsilon}{2\epsilon s}\bigl(1+o_v(1)\bigr),
 \qquad
 y_{\epsilon,\delta}=\delta r_\epsilon\bigl(1+o_v(1)\bigr),
\end{equation}
where $o_v(1)$ denotes an element of positive valuation. The four sign choices exhaust the cluster by \eqref{finite:eq:conservation}. At each root
\[
 \frac{e^{x_{\epsilon,\delta}}}{J_F(x_{\epsilon,\delta},y_{\epsilon,\delta})}
 =\frac1{4\epsilon\delta\,s\,r_\epsilon}\bigl(1+o_v(1)\bigr),
 \qquad
 v\Bigl(\frac1{4s\,r_\epsilon}\Bigr)=-\frac\omega2-\frac34<0,
\]
so every individual residue is infinitely large, while their total has standard part $\tfrac16$ by \eqref{finite:eq:quarticexp}: the fixed coefficientwise cycle performs the cancellation across all three scales automatically.

\subsection{A triple collision resolved at an infinitely smaller scale}
\label{finite:sec:triple}

Choose a positive \emph{infinite} surreal exponent $\Omega$, work in $\Gamma=\Q+\Q\Omega$, and set
\[
 t=\omega^{-1},\qquad s=t^\Omega,\qquad a_0=\tfrac34t^2+s,
\]
so that $s\prec t^N$ for every ordinary $N$. Consider
\begin{equation}\label{finite:eq:tripleF}
 F_1(x,y)=x^2-ty-a_0,
 \qquad
 F_2(x,y)=y^2-tx-a_0 ,
\end{equation}
with leading system $(x^2,y^2)$, positive support inside $\{1,2,\Omega\}$ and $\mu=4$ on the basis $(1,x,y,xy)$. Subtracting, $F_1-F_2=(x-y)(x+y+t)$, so every root lies on one of two affine lines:
\begin{equation}\label{finite:eq:tripleroots}
 (r_\pm,r_\pm)\ \text{with}\ r_\pm=\frac t2\pm\sqrt{t^2+s},
 \qquad
 u_\pm=\Bigl(-\frac t2\pm\sqrt s,\ -\frac t2\mp\sqrt s\Bigr).
\end{equation}
The binomial series, strongly evaluated at $s/t^2$, gives $r_+=\tfrac{3t}2+\tfrac s{2t}-\tfrac{s^2}{8t^3}+\tfrac{s^3}{16t^5}-\cdots$ and $r_-=-\tfrac t2-\tfrac s{2t}+\tfrac{s^2}{8t^3}-\cdots$; the supports are well ordered because successive exponents differ by $\Omega-2>0$.

The geometry uses two infinitesimal levels. One point sits near $(3t/2,3t/2)$ and three near $(-t/2,-t/2)$: the valuation of every difference between the distant point and the others is $1$, the valuation of every pairwise difference among the three nearby points is $\Omega/2$, and the diagonal member's displacement from their common centre has the still larger valuation $\Omega-1$. With $J=4xy-t^2$, the Jacobian is exactly $-4s$ at the off-diagonal points, $2s+\tfrac{s^2}{2t^2}-\tfrac{s^3}{4t^4}+\cdots$ at the near diagonal point, and has leading term $8t^2$ at the distant one; all four are nonzero, so each point is simple.

Setting $s=0$ instead collapses this to a simple point at $(3t/2,3t/2)$ and a point of \emph{multiplicity three} at $(-t/2,-t/2)$. This is visible in the algebra itself. In the basis $(1,x,y,xy)$,
\begin{equation}\label{finite:eq:triplematrices}
 M_x=\begin{pmatrix}0&a_0&0&ta_0\\1&0&0&t^2\\0&t&0&a_0\\0&0&1&0\end{pmatrix},
 \qquad
 M_y=\begin{pmatrix}0&0&a_0&ta_0\\0&0&t&a_0\\1&0&0&t^2\\0&1&0&0\end{pmatrix},
\end{equation}
which commute and satisfy $M_x^2=tM_y+a_0I$, $M_y^2=tM_x+a_0I$, with
\begin{equation}\label{finite:eq:triplechar}
 \chi_x(X)=X^4-2a_0X^2-t^3X+a_0^2-a_0t^2
 =(X^2-tX-a_0)(X^2+tX+t^2-a_0),
\end{equation}
and at $s=0$ this factors as $(X-3t/2)(X+t/2)^3$, proving the local lengths $1$ and $3$ without any picture of the roots. Exact cluster moments come for free from \eqref{finite:eq:trace}:
\[
 \sum_\zeta x(\zeta)=\sum_\zeta y(\zeta)=0,
 \qquad
 \sum_\zeta x(\zeta)^2=4a_0=3t^2+4s,
 \qquad
 \sum_\zeta x(\zeta)y(\zeta)=3t^2 ,
\]
the last being completely insensitive to $s$, although $s$ is what separates three roots at a scale invisible in any fixed finite power expansion in $t$.

The residue data are equally explicit: $\Res_F(1)=\Res_F(x)=\Res_F(y)=0$ and $\Res_F(xy)=1$ by the same coefficient count as in \cref{finite:sec:collision}; hence
\[
 \Gram_F=\begin{pmatrix}0&0&0&1\\0&0&1&0\\0&1&0&0\\1&0&0&t^2\end{pmatrix},
 \quad\det\Gram_F=1,
 \quad
 \text{dual basis }(xy-t^2,\ y,\ x,\ 1),
\]
so that $\Eul_F=(xy-t^2)+xy+yx+xy=4xy-t^2=J_F$, an exact instance of \cref{finite:thm:bezout}. For $s\neq0$ every zero is simple and $J_F$ is a unit, so \eqref{finite:eq:simpleresidue} gives the complete moving-residue formula
\begin{equation}\label{finite:eq:triplemoving}
 \Res_F(G)=\sum_{\zeta\in\{(r_+,r_+),(r_-,r_-),u_+,u_-\}}\frac{G\sh(\zeta)}{J_F\sh(\zeta)}
\end{equation}
for \emph{every} coherent numerator on $D\sh$, not only for polynomials. At $G=1$ it reads
\begin{equation}\label{finite:eq:triplecancel}
 \frac1{J(r_+,r_+)}+\frac1{J(r_-,r_-)}-\frac1{2s}=0,
\end{equation}
the near-diagonal contribution having leading term $1/(2s)$ and the two off-diagonal contributions being each $-1/(4s)$: their largest terms cancel exactly, the remaining cancellation involving the distant point. Replacing the displaced roots by their common standard part would destroy this information entirely.

\subsection{Residue duality versus trace degeneration, and lifted monodromy}
\label{finite:sec:cubic}

Let $a,b\in\Ip_\Gamma(U)$ be positive-support coherent parameters and $P(w)=w^3-aw-b$, with basis $(1,w,w^2)$ and $w^3=aw+b$. By \cref{finite:prop:univtrace} the \emph{residue} Gram matrix and the \emph{trace} Gram matrix are
\begin{equation}\label{finite:eq:cubicgrams}
 \Gram_{\mathrm{res}}=\begin{pmatrix}0&0&1\\0&1&0\\1&0&a\end{pmatrix},
 \quad\det=-1,
 \qquad
 \TGram=\begin{pmatrix}3&0&2a\\0&2a&3b\\2a&3b&2a^2\end{pmatrix},
 \quad\det=4a^3-27b^2 .
\end{equation}
The residue-dual basis is $(w^2-a,\ w,\ 1)$, so $\Eul_P=3w^2-a=P_w$. At a repeated root the trace pairing degenerates while the residue pairing does not; this distinction is exactly what justifies the word \emph{duality} in \cref{finite:thm:duality} and the two-symbol convention \eqref{finite:eq:twodiscriminants}.

For an explicit collision take $a=3t^2$, $b=2t^3$, so that $P=(w-2t)(w+t)^2$ with local multiplicities $1$ at $2t$ and $2$ at $-t$. Partial fractions,
\[
 \frac1P=\frac1{9t^2(w-2t)}-\frac1{9t^2(w+t)}-\frac1{3t(w+t)^2},
\]
together with the one-variable Cauchy derivative formula give
\begin{equation}\label{finite:eq:doublepolelambda}
 \lambda_P(G)=\frac{G(2t)-G(-t)}{9t^2}-\frac{G'(-t)}{3t},
 \qquad\text{while}\qquad
 \Tr(M_G)=G(2t)+2G(-t).
\end{equation}
The nondegenerate functional remembers a first derivative at the double point --- information wholly invisible to the trace pairing on nilpotent directions.

For \cref{finite:thm:monodromy}, take $P(z,w)=w^3-z-tw$ on an ordinary annulus $0<r<|z|<R$ with a sufficiently large ordinary fibre disk. The leading polynomial $w^3-z$ has three simple roots with the ordinary cyclic monodromy around zero, and \cref{finite:thm:monodromy} lifts these sheets coherently, preserving the three-cycle; on a local choice $c^3=z$ the first correction is
\[
 w=c+\frac t{3c}+\text{terms of higher positive support}.
\]
The exact polynomial discriminant is $4t^3-27z^2$, so collisions occur at $z=\pm\tfrac2{3\sqrt3}t^{3/2}$, both of standard part zero and therefore outside the thickening of the ordinary annulus. The monodromy theorem neither misses nor forbids them: its domain intentionally excludes the whole standard-zero parameter monad, while the finite algebra and the residue pairing remain meaningful across such collisions on a larger parameter chart to which \cref{finite:thm:parameterprep} applies.

\subsection{Infinitely many exponents below a bound}
\label{finite:sec:infsupport}

The theorems cover inputs that no finite list of ordinary deformation parameters describes.

\begin{example}[A well-ordered support with infinitely many exponents below $1$]
\label{finite:ex:infsupport}
In a workspace containing $\Q$, let $\lambda_m=1-\tfrac1{m+1}$ for $m\geq1$ and, on an ordinary bidisk,
\[
 F_1(x,y)=x^3+\sum_{m\geq1}t^{\lambda_m}e^{m!x}(1+y^m),
 \qquad
 F_2(x,y)=y^2+\sum_{m\geq1}t^{\lambda_m}e^{m!y}(x+x^my).
\]
The support $\{\lambda_m\}$ is strictly increasing and well ordered although infinitely many of its exponents lie below $1$, and the coefficient functions are entire with no uniform growth bound in $m$. Nevertheless $\Alg$ has basis $1,x,x^2,y,xy,x^2y$ and rank six, all six intersections counted with local multiplicity lie in the infinitesimal bidisk, the residue pairing is perfect, and the multiplication and residue coefficients are computed by the support-controlled formulas.
\end{example}

\begin{example}[A support of order type $\omega+1$, with no ordinary numerical convergence]
\label{finite:ex:nondiscrete}
Let $\gamma_k=1-\tfrac1{k+2}$ and
\[
 F_1(x,y)=x^2+\sum_{k\geq0}t^{\gamma_k}e^{k!x+y},
 \qquad
 F_2(x,y)=y^3+t\,e^{x-y},
\]
so that the common error support is $\{\gamma_k\}\cup\{1\}$, of order type $\omega+1$ with supremum attained. Here $m_1=2$, $m_2=3$, the basis is $1,y,y^2,x,xy,xy^2$, every actual fibre has total intersection multiplicity six, and all normal-form and residue coefficients are controlled by $M=(\{\gamma_k\}\cup\{1\})^*$. At $x=y=0$, replacing $t$ by a fixed ordinary real $q\in(0,1)$ would make the summands $q^{\gamma_k}$ tend to $q$ rather than to zero, so ordinary numerical convergence is unavailable even at that point; and the perturbation cannot be an ordinary integer-power series in one positive monomial, since infinitely many distinct exponents lie in a bounded interval. The strong-support proof needs neither property.
\end{example}

An algorithm asked for ``all coefficients of valuation less than $1$'' would have infinite output in these examples. That does not contradict finite contribution at each \emph{specified} exponent; the two notions of truncation must not be conflated, as \cref{finite:rem:truncation} says.

\subsection{Nonpolynomial coefficient functions}
\label{finite:sec:nonpoly}

\begin{example}[A nonpolynomial complete intersection of multiplicity five]
\label{finite:ex:five}
The general theorem is not restricted to a rectangular monomial basis. Take $f_1=x^2-y^3$, $f_2=xy$, whose only common zero is the origin, with relations $x^2=y^3$, $xy=0$, $y^4=0$ and basis $1,x,y,y^2,y^3$, so $\mu=5$; a lexicographic Gr\"obner basis is $\{x^2-y^3,xy,y^4\}$. Its embedding dimension is two, both defining equations having vanishing linear part, so it cannot be carried by an analytic coordinate change to a two-variable rectangular monomial complete intersection of length five, such a rectangle having side lengths $1$ and $5$ and embedding dimension one. With $t=\omega^{-1}$ and $\Gamma=\Q+\Q\omega$ put
\[
 F_1=x^2-y^3-t\,e^{x+y},
 \qquad
 F_2=xy-t^\omega\cos x,
\]
genuinely nonpolynomial coherent equations whose coefficient functions are entire, hence elements of $\HH_\Gamma(D)$ for every fixed polydisk $D$, with error support inside $\{1,\omega\}$. \Cref{finite:thm:finitefree,finite:thm:length} give a rank-five integral algebra in the unchanged basis, every coherent input having a unique remainder there and witnesses supported in $S+\{m+n\omega:m,n\in\N\}$. There are exactly five roots counted with formal local multiplicity in the origin monad, all in $\C((t^{\Q+\Q\omega}))$, and for every infinitesimal target $(\eta_1,\eta_2)$ the equations $F_1=\eta_1$, $F_2=\eta_2$ have the same total multiplicity five; this neither assumes nor needs that any individual root is simple. Since $J_f=2x^2+3y^3\equiv5y^3\pmod{(f)}$, \cref{finite:thm:trace} gives the exact statement $\Res_F(J_F)=5$.
\end{example}

\begin{example}[A transcendental coupled system with four simple roots]
\label{finite:ex:transcendental}
Let $a,b\in\mm_K\setminus\{0\}$ and $F_1=x^2-ae^y$, $F_2=y^2-be^x$, the exponentials being coherent lifts of ordinary entire functions. The leading system $(x^2,y^2)$ has $\mu=4$. At a zero, $x$ and $y$ are nonzero, and since $x^2y^2=abe^{x+y}$,
\[
 J_F=4xy-abe^{x+y}=xy(4-xy),
\]
with $xy$ infinitesimal, so $4-xy$ is a unit and every zero is simple: there are exactly four. Writing $s=\sqrt a$, $r=\sqrt b$, $x=sU$, $y=rV$ turns the system into $U^2=e^{rV}$, $V^2=e^{sU}$, whose ordinary specialization at $(s,r)=(0,0)$ has the four simple solutions $(U,V)=(\epsilon,\delta)$, $\epsilon,\delta\in\{\pm1\}$; the ordinary implicit-function expansions in the two independent parameters evaluate strongly, giving $x_{\epsilon,\delta}=\epsilon\sqrt a(1+o_v(1))$ and $y_{\epsilon,\delta}=\delta\sqrt b(1+o_v(1))$. This works equally for $a=t$, $b=t^\omega$ and for positive supports with infinitely many exponents. \Cref{finite:cor:coefficientresidue} with numerator $e^{x+y}$ gives the exact two-scale formula
\begin{equation}\label{finite:eq:transresidue}
 \Res_F(e^{x+y})=\sum_{r,s\geq0}a^rb^s\,
 \frac{(s+1)^{2r+1}(r+1)^{2s+1}}{(2r+1)!\,(2s+1)!}
 =1+\frac{a+b}3+\frac{a^2+b^2}{40}+\frac{16ab}9+\cdots,
\end{equation}
strongly summable for every pair of nonzero infinitesimals, while \cref{finite:thm:trace} equates it to the finite sum $\sum_{\epsilon,\delta}e^{x+y}/\bigl(xy(4-xy)\bigr)$ over the four roots, each summand having leading value $\epsilon\delta/(4\sqrt{ab})$, which is infinitely large. If instead $a=0$ and $b\neq0$, the zeros are $(0,\pm\sqrt b)$, each of multiplicity two: the second equation solves formally for $y$, leaving $x^2=0$. The total is still four, \eqref{finite:eq:transresidue} still defines the residue after setting $a=0$, and the simple-root expression must not be used there, its Jacobians vanishing --- \cref{finite:rem:multiple}.
\end{example}

%======================================================================
\section{Computation, scope, and what is not proved}
\label{finite:sec:scope}

\subsection{What can be computed coefficient by coefficient}
\label{finite:sec:computation}

When an explicit contraction and effective finite dependency lists are supplied, \eqref{finite:eq:NFformula} and the recursion \eqref{finite:eq:recursion} give a direct coefficient procedure: to obtain one normal-form coefficient at exponent $\nu$, enumerate the finitely many words in the positive perturbation supports and the input exponents summing to $\nu$, apply the corresponding finite compositions of ordinary operators and add. In the coordinate-power case, \cref{finite:lem:matrix} further reduces those ordinary operations to finitely many Taylor jets. Repeating on $Ge_a$ produces multiplication matrices; characteristic polynomials then bound every coordinate of every zero; traces, $\Delta_F$ and $\disc_F$ require only finite matrix operations. Cluster residues can be computed by \eqref{finite:eq:coefficientresidue} without solving for any roots, or recovered from finitely many residue moments by \cref{finite:cor:residuereconstruction}; when all roots are simple the formula $G\sh(\zeta)/J_F\sh(\zeta)$ is equivalent but can expose large cancelling quantities, as \cref{finite:sec:quartic} shows.

Three qualifications are part of the statement, not a disclaimer appended to it.

\begin{enumerate}
\item The finiteness is \emph{mathematical}. An arbitrary ordered group, or an arbitrary well-ordered input support, need not come with an effective representation or an algorithm for solving the required exponent equations. The mathematical recursion is well-founded; it may not be an algorithm on finite input. To obtain a terminating symbolic procedure one must supply computable coefficient operations and a procedure returning the complete finite list of input exponents and positive-support words contributing to a requested output exponent. A finite positive generating set by itself, or a membership test without a completeness bound, does not provide this procedure in an arbitrary effectively ordered group. In finite-rank rational exponent models with explicit finite decomposition routines, exact sparse exponent arithmetic supplies suitable input data.
\item An initial segment below a valuation bound can be infinite, by \cref{finite:rem:truncation} and \cref{finite:ex:infsupport}. Software should therefore compute \emph{specified} coefficients, finite algebraic identities, or finitely generated support expressions, and should not promise a finite enumeration of every term below a chosen valuation.
\item An unspecified complex-linear contraction can be nonconstructive; a finite witness means finitely many ordinary coefficient \emph{functions}, not finitely many Taylor coefficients of them, and the jet bound is a separate statement. The parameter preparation operators \eqref{finite:eq:ordinaryremainder}--\eqref{finite:eq:ordinaryquotient} and the coordinate operators \eqref{finite:eq:Rj} are, by contrast, explicit in terms of ordinary differentiation, division by a coordinate monomial, and a fixed contour.
\end{enumerate}

\subsection{Why the hypotheses cannot simply be dropped}
\label{finite:sec:hypotheses}

\paragraph{The ordinary zero must be isolated.} The pair $(xy,0)$ has no finite-dimensional leading zero algebra and its zero set contains whole coordinate monads; neither a finite basis nor a finite degree follows from the number of displayed equations. Similarly $(x,0)$ has infinite colength, and its perturbations $(x,ty)$ and $(x,0)$ have completely different zero sets.

\paragraph{The system must be a complete intersection.} Arbitrary redundant generators cannot be fed into the same positive-degree Koszul exactness argument; an isolated ideal with more than $n$ generators has syzygies not encoded by the Koszul complex.

\paragraph{The perturbation must preserve the leading fibre.} Starting from $(x^2,y^2)$ and changing the first equation to $x^2-x$ adds a term at exponent zero; the origin-monad multiplicity becomes two rather than four, the other ordinary root lying in a different monad. Strict positivity is exactly what keeps the exponent-zero resolution and its rank unchanged.

\paragraph{Integral normalization matters.} Over $K$, multiplying an equation by a nonzero Hahn scalar does not change its ideal. Over $\Kc$, multiplying by a positive infinitesimal need not preserve the \emph{integral} zero algebra: the quotient by $\tau x$ has a nonzero class of $x$ annihilated by $\tau$ and is not the finite free normalized quotient by $x$. Equations with different leading monomial factors must first be normalized, and the integral finite-flat assertion is about that normalized model.

\paragraph{There must be a common ordinary domain.} A Hahn family of convergent coefficient germs admitting no common neighbourhood belongs to the radius-free ring (3) of \cref{finite:sec:rings}, but neither to any fixed-domain ring (1) nor to the common-domain germ ring (4). An algebra of all coefficientwise germs is a different object, and none of the arguments above quietly enlarges the function class at an infinite stage.

\paragraph{Uniform support is indispensable.} \Cref{finite:ex:unit} is the decisive instance: in $K[[z]]$ the basic quotient for $z^2-t$ disappears entirely, because the non-well-ordered inverse \eqref{finite:eq:badinverse} is allowed there. It is not legitimate to run the proof in that ring and then declare that all displaced roots have been counted.

\paragraph{Strong summability is not a topological iteration theorem.} \Cref{finite:rem:notlimit}.

\subsection{What is explicitly not proved}
\label{finite:sec:notproved}

The following non-claims are carried over verbatim in substance from the five merged sources and are part of the statement of the results. They are listed separately rather than pooled, because each attaches to a specific theorem.

\paragraph{On the meaning of finiteness and flatness.}
``Finite flat'' follows here from the strictly stronger, explicitly proved ``finite free''; no Noetherian deformation theorem is invoked for $\Kc$, which can have high-rank value group and need not be Noetherian, and ordinary Noetherian local algebra is used only to construct the exponent-zero resolution. In the usual algebraic sense $\Spec\Alg^+\to\Spec\Kc$ is finite flat of rank $\mu$; \emph{that statement is not a claim about proper maps in an undeclared topology on the class $\SC^n$}. Correspondingly, ``finite map'' means the stated finite-fibre and finite-algebra properties and \emph{does not assert properness for a topology on a proper class}. The observation that the positive ideal of $\Hp_\Gamma(\mathrm{pt})$ is not finitely generated \emph{does not say that the field $K$ is non-Noetherian}.

\paragraph{On summation, truncation, and iteration.}
$E^*\cap(-\infty,\eta)$ need not be finite; $1+t+t^2+\cdots$ is a valid Hahn sum although its ordinary finite partial sums do not approximate it to valuation $>\omega$. Infinite formulas here are justified by strong summation, \emph{never} by silently treating $k\min E$ as cofinal in $\Gamma$. There is no assertion that $T^kX$ tends to zero in the full fine topology, and ``the valuation increases'' is never substituted for the finite-word support argument. The coefficient recursion \eqref{finite:eq:recursion} is an exact algorithm when ordinary operations and complete finite dependency lists are computable; decidable support membership alone is insufficient; it is \emph{not} a blanket claim that every valuation truncation can be listed in finite time, and an arbitrary set-sized well-ordered support is not automatically an effective input format. All index sets are sets, not proper classes.

\paragraph{On the contraction and the perturbation lemma.}
The existence proof for $h$ uses ordinary choices of complex-vector-space complements and \emph{is not, by itself, an effective algorithm for an arbitrary black-box holomorphic system}. \Cref{finite:thm:hpl} is not a new abstract perturbation lemma: invertibility of $1+\delta h$ is already the appropriate algebraic smallness condition in the classical statement, and what has been checked is that the inverse exists on the full Hahn modules, respects one ordinary coefficient domain, and carries the exact support estimate. The general perturbation lemma allows homology in several degrees with an induced differential on it; here that differential vanishes only because the target is concentrated in degree zero. The family criterion of \cref{finite:rem:familycriterion} is \emph{a precise additional hypothesis, not an assertion that every family automatically has such a splitting on any prescribed domain}, and the abstract criterion is an algebraic extension, not a claim that an arbitrary vector-space splitting is analytically satisfactory.

\paragraph{On contours and residues.}
The integrals in the coefficients are ordinary complex integrals: \emph{the superscript $H$ does not denote a fine-topological Riemann integral over a surcomplex parametrized torus}, and no ordinary circle is claimed to be a fine-continuous path in $\SC$. The ordinary residue inputs are \emph{classical local-residue and complete-intersection duality facts, not new surcomplex assertions}, and the multivariate residue literature contains many further global results which are not being claimed here. $\Res_F$ is specified by its defining formula including the order of the equations and the orientation convention; \emph{we do not assume without proof a universal identification with every Grothendieck-residue construction} over non-Noetherian Hahn analytic rings. The residue is generally \emph{not} the top coefficient of the normal form. Two of the merged sources decline to assert $\Eul_F=[J_F]$ in general; see \cref{finite:rem:eulerconflict} for exactly how far it is proved here and by which argument.

\paragraph{On multiplicities, roots, and stability.}
The perfect pairing and the trace--Jacobian identity do not require simple zeros; the formula with individual terms $G(\zeta)/J_F(\zeta)$ does, and at a multiple zero that expression must not be used. \emph{No unproved identification of separately localized multiple-zero residues is needed for any result here.} The theorems do \emph{not} assert that a labelling of individual roots is Lipschitz: roots of a multiple leading system generally involve fractional exponents and inverse Jacobians can have negative valuation. All inequalities in \cref{finite:sec:stability} concern the coefficient valuation, not an ordinary analytic norm. The valuation localization \eqref{finite:eq:rootbound} is a bound, not an assertion that roots have Puiseux series in one chosen parameter. The support bounds for multiplication matrices are \emph{not} asserted for individual roots: root extraction can require divisible exponents, $z^2-t$ already having roots $\pm t^{1/2}$, and the root theorem keeps all roots in the divisible workspace $K$, not necessarily in the monoid generated by the original error support. The word Rouch\'e refers to conservation under a positive Hahn perturbation of the leading system and is \emph{not} an unrestricted boundary-norm theorem for arbitrary functions on $\SC^n$. The monodromy theorem does not postulate fine-continuous loops in $\SC$ and deliberately excludes the standard-zero parameter monad.

\paragraph{On the coefficient rings and on evaluation.}
The formal local statements are formal and unambiguous; the coherent Taylor series also represents the function at sufficiently small surcomplex displacements, but this \emph{does not assume a general theorem identifying arbitrary fine-analytic sheaves with Noetherian convergent-series rings}, and \cref{finite:cor:inverse} does not turn arbitrary fine-local analytic functions into uniformly supported data on a monad. A possible pitfall is to assume that an arbitrary algebra homomorphism commutes with an infinite Hahn sum; that assumption is avoided throughout by using finite identities in the coherent ring. The finiteness and separation hypotheses of \cref{finite:thm:monadcount} concern the \emph{ordinary} base and do not assert compactness of $\Omega\sh$ in the fine topology. ``Global support control'' means uniform support on a stated ordinary domain and \emph{not} the all-scale coherence condition on every surreal affine chart, which is a different and much more restrictive condition. Finally, the terminology of Hahn holomorphic and Hahn meromorphic functions in the literature refers to convergent generalized-power expansions near zero on a logarithmic covering of the ordinary complex plane \cite{finite:MullerStrohmaier}; \emph{it should not be confused with the uniformly supported holomorphic coefficient families used here}, and no comparison of all such constructions is claimed. Nor is any extension of surreal integration and resummation programs to arbitrary complex sectors asserted \cite{finite:CostinEhrlich}.

\subsection{What remains outside the theorems}
\label{finite:sec:outside}

The results settle a delimited target: fixed-domain division and finite-mapping theorems for positive Hahn perturbations of isolated ordinary complete intersections, with holomorphic-parameter versions, multidimensional residue duality, stability, and an automatic-coherence theorem for the function class. They do not provide a complete analytic-set theory for every ideal of $\HH_\Gamma(U)$, nor a sheaf-coherence theorem for arbitrary covers whose supports change from chart to chart, nor a global theory of essential singularities, nor a Noetherian analytic geometry over the whole surcomplex class.

Four concrete continuations are well posed. \emph{Positive-dimensional intersections}: when the ordinary quotient has infinite dimension the perturbation lemma still has formal meaning if a contraction exists, but the finite spectral and residue arguments fail; a useful target is a coherent module over an ordinary parameter domain with a locally free finite-rank direct image and transition maps of globally compatible support --- applying the present theorem at each ordinary point would not prove that compatibility. \emph{General finite analytic ideals}: a finite free resolution with a perturbation of its full differential satisfying the square-zero equations is a plausible replacement for the Koszul complex, the correct compatibility hypothesis being part of the problem. \emph{Explicit contractions for general leading singularities}: \cref{finite:rem:familycriterion} states exactly what additional input would extend the parameter results; ordinary local analytic division suggests a route, but its domain and support consequences must be proved rather than assumed. \emph{Globalization}: gluing finite coherent families across ordinary charts, comparing these coefficientwise residues with sectorial or resurgent procedures, and characterizing which parameter-dependent finite free algebras admit coherent root branches after a controlled ramification. The present results count fibres and supply coherent multiplication matrices; they do not promise globally single-valued root branches around an ordinary discriminant locus, and the all-scale rigidity phenomenon shows why a naive atlas condition can exclude transcendental functions altogether.

\subsection{Consolidated status}
\label{finite:sec:status}

All five merged sources carry a status disclaimer of the same shape; it is stated here once, and it governs the whole report.

A targeted literature search was carried out through 21 September 2026, covering the supplied manuscript, surreal analytic and exponential structure \cite{finite:Alling,finite:vdDE,finite:BM}, Hahn-field algebra and closure \cite{finite:Neumann,finite:Poonen}, analytic structures on valued fields \cite{finite:CluckersLipshitz}, homological perturbation \cite{finite:Crainic}, ordinary local residue and finite-algebra theory \cite{finite:Demailly,finite:Kunz,finite:CattaniDickenstein,finite:CDS,finite:Cox,finite:NayakSastry}, Hahn-meromorphic function theory \cite{finite:MullerStrohmaier}, surreal integration \cite{finite:CostinEhrlich}, and higher-rank Hahn geometry \cite{finite:JoswigSmith}. The search was targeted, not exhaustive. It did not locate the exact combined statements proved here. \emph{That is a statement about the search, not a proof of priority}: there may be equivalent results under different terminology, or stronger general theorems from which some assertions can be derived, and a specialist might identify a theorem from formal, rigid or non-Archimedean analytic geometry that subsumes some assertions after a careful comparison of categories. Such an identification would change the priority assessment, not the proofs.

No named published conjecture is claimed solved. The open target came from an unpublished supplied manuscript, and no published conjecture whose established name and exact hypotheses coincide with any theorem here has been identified. The article has not been independently refereed and has not been formally verified. The symbolic checks recorded in the source archives --- exact SymPy computations of the multiplication matrices, characteristic polynomials, discriminants, residue Gram matrices, dual bases, trace--Jacobian identities, moment recursions, series expansions and finite instances of the Koszul contraction, for the examples of \cref{finite:sec:examples} --- validate \emph{displayed examples}; they are not machine verification of the general analytic, homological or arbitrary-support theorems, which rest on the mathematical arguments and the identified classical inputs above. Independent review would be most valuable for the two analytic-to-algebraic bridges, \cref{finite:lem:matrix,finite:lem:witness}, since finite dimension by itself does not identify abstract characters with analytic evaluation, and ordinary continuity in a real parameter does not automatically justify a Hahn substitution.

%======================================================================
\section{Conclusion}
\label{finite:sec:conclusion}

A multiple ordinary zero does not merely leave behind a list of displaced surcomplex points. Under a Hahn perturbation of strictly positive well-ordered support it leaves behind a finite free algebra of unchanged rank, and that algebra is the correct carrier of multiplicity, residues and spectral information. The decisive construction is to extend one ordinary Koszul contraction coefficientwise and to sum only along a positive well-ordered support: this keeps every division coefficient on the original ordinary domain and avoids every unjustified convergence assumption, because the inverse is legitimated by finite contribution at each exponent rather than by an iteration in a topology.

The finite algebra is not by itself a theorem about zeros. Shifted division at an infinitesimally displaced centre provides the bridge; for coordinate-power leading systems, a summable matrix functional calculus supplies an independent alternative. These arguments identify the characters with actual surcomplex evaluations and the local Artinian factors with the formal local algebras \emph{at the displaced points}. Only then does conservation of multiplicity become a statement of surcomplex analysis, and only then does the monad map acquire finite fibres of constant total multiplicity over every infinitesimal target. A third bridge, finite-witness transfer, carries the ordinary residue and trace identities coefficientwise onto the Hahn functional, whose pairing is perfect even over the full valuation ring and whose trace formula survives collisions and individually infinite residues.

Around that core, four independent blocks complete the picture. Automatic Hahn--Hartogs coherence shows that the function class itself is better behaved than its definition suggests: separate coherence forces joint coherence, and the global well-ordered support is a conclusion. Matrix division is what controls coupling when equations are eliminated in succession. The precision theorem and the sharp conditioned stability theorem separate robust cluster data --- which lose no valuation at all --- from sensitive individual roots, which lose precisely the inverse-Jacobian loss $\kappa$, with the threshold $\sigma>2\kappa$ and both bounds attained in an explicit quadratic model. And the monodromy lift shows that the ordinary covering structure above the simple-root locus survives the perturbation unchanged.

Throughout, every theorem names its coefficient ring, and no theorem has been moved between rings.

%======================================================================
\appendix
\section{Support audit}
\label{finite:app:support}

Every index set in this report is a set. Let $S$ be an input support, $E$ the common strictly positive perturbation support with $M=E^*$, and $A$ the union of the supports of an infinitesimal centre, target or parameter displacement. Each entry below is a well-ordered set by \cref{finite:lem:support}, and each bound is the one actually proved.

\begin{center}\small
\begin{longtable}{@{}L{0.47\textwidth}L{0.45\textwidth}@{}}
\toprule
Construction & Controlling support\\
\midrule
\endhead
Product of two Hahn data & Minkowski sum of the two supports\\[0.3em]
Coefficientwise complex-linear map & A subset of the input support $S$\\[0.3em]
Coherent evaluation at $c+\varepsilon$; translation & $S+A^*$\\[0.3em]
Shifted division and finite Taylor remainder & $S+A^*$\\[0.3em]
Preparation corrections $A$, $C$ & $M\setminus\{0\}$\\[0.3em]
Division by a prepared polynomial; matrix division & $S+M$\\[0.3em]
Resolvent $S=(1+\delta h)^{-1}$ & $S+M$\\[0.3em]
Perturbed contraction, normal forms, witnesses & $S+M$\\[0.3em]
Syzygy lifting $h_FZ$ & $S+M$\\[0.3em]
Structure constants; multiplication matrices of ordinary basis inputs & $M$\\[0.3em]
Characteristic-polynomial coefficients for ordinary-coefficient inputs & $M$\\[0.3em]
Hahn residue $\Res_F(G)$ & $S+M$\\[0.3em]
Residue Gram matrix $\Gram_F$ and its inverse & $M$\\[0.3em]
Residue-dual basis; trace element $\Eul_F$ & $M$\\[0.3em]
Remainder recovered from residue moments & $S+M$\\[0.3em]
Coherent parameter pullback along $\phi_0+\eta$ & $S+(\bigcup_j\supp\eta_j)^*$\\[0.3em]
Specialized family constructions & The corresponding bound plus $A^*$\\[0.3em]
Monodromy lift $\eta$ of a simple root & $\supp(E_0)^*\setminus\{0\}$\\[0.3em]
A prescribed output coefficient & Finitely many contributing words; \emph{not} an entire initial segment\\
\bottomrule
\end{longtable}
\end{center}

For every infinite operator formula, finite word multiplicity supplies both well-ordering and finite contribution at each exponent; finite-dimensional matrix operations then need only finite sums and products, except for unit inverses, whose positive-support geometric series are covered by the same lemma. None of these bounds requires $\Gamma$ to have rank one. Finally, the bounds above are \emph{not} asserted for individual roots: root extraction can require divisible exponents, and the root theorems place all roots in $K$, not necessarily in $M$.

\section{Proof dependencies and provenance}
\label{finite:app:audit}

\subsection{Established inputs and the extension steps}

\begin{center}\small
\begin{longtable}{@{}L{0.20\textwidth}L{0.33\textwidth}L{0.40\textwidth}@{}}
\toprule
Stage & Established input & Extension step proved here\\
\midrule
\endhead
Coefficient calculus & Positive-support Neumann lemma \cite{finite:Neumann,finite:Poonen} &
Strong evaluation; global divided differences on the unshrunk domain; support-preserving linear maps (\cref{finite:prop:evaluation,finite:lem:shifted,finite:lem:operator})\\[0.35em]
Function class & Classical Hartogs separate holomorphy and extension \cite{finite:BK,finite:Lebl} &
Automatic common support and separate $\Rightarrow$ joint coherence, with the global support a conclusion (\cref{finite:lem:autosupport,finite:thm:hartogs})\\[0.35em]
Ordinary finite algebra & Regular local rings, Koszul exactness, Cartan's Theorem~B \cite{finite:Cartan,finite:GunningRossi,finite:Demailly,finite:StacksKoszul} &
One global algebraic contraction on $\OO(D)$, fixed-domain (\cref{finite:prop:ordinarykoszul,finite:constr:h})\\[0.35em]
Ordinary division & Weierstrass division and preparation \cite{finite:GunningRossi,finite:Lebl} &
Fixed-domain contour operators with exactly one shrink; Hahn preparation and matrix division with the $S+M$ bound (\cref{finite:lem:ordinarydivision,finite:thm:matrixdivision,finite:thm:parameterprep})\\[0.35em]
Perturbed division & Homological perturbation identities \cite{finite:Crainic} &
Strong resolvent on the full Hahn modules; fixed-domain witnesses; integral finite freeness over a possibly non-Noetherian $\Kc$ (\cref{finite:thm:hpl,finite:thm:finitefree})\\[0.35em]
Geometry & Hahn algebraic closedness \cite[Cor.~4]{finite:Poonen}; finite Artinian algebras &
Characters are actual evaluations; every local factor is the formal germ algebra \emph{at the displaced zero}; conservation of $\mu$; finite fibres over every infinitesimal target (\cref{finite:thm:characters,finite:thm:length,finite:cor:monadmap,finite:thm:fibres})\\[0.35em]
Residues & Ordinary Grothendieck duality, local cycles, B\'ezoutians \cite{finite:GriffithsHarris,finite:Hartshorne,finite:Kunz,finite:CattaniDickenstein,finite:CDS,finite:NayakSastry} &
Perfect pairing over $\Kc$ surviving collisions; finite-witness trace transfer; reproducing kernel; localized residues at actual points; discriminant identities and valuation estimates (\cref{finite:thm:duality,finite:thm:trace,finite:thm:bezout,finite:thm:localres,finite:thm:discriminant,finite:thm:precision})\\[0.35em]
Stability & Hensel-type conditioning &
Sharp threshold $\sigma>2\kappa$ with both valuation bounds attained, over arbitrary-rank value groups (\cref{finite:thm:rootstability,finite:cor:matching,finite:prop:sharpness})\\
\bottomrule
\end{longtable}
\end{center}

The assumptions may also be audited by the step at which each enters. Isolatedness gives a finite $A_0$ with nilpotent coordinate classes. The complete-intersection condition gives positive-degree Koszul exactness. The common ordinary domain permits one global $h$ and forbids an infinite sequence of shrinkings. Strict positivity of the error supports gives the geometric inverse and prohibits support loss. Divisibility of $\Gamma$ and algebraic closedness of $\C$ ensure that every character of the finite algebra has coordinates in the chosen $K$. Ordinary residue duality gives the nonzero leading determinant of $\Gram_F$. None of these steps uses a discrete valuation, an Archimedean value group, a Noetherian Hahn valuation ring, eventual convergence of finite partial sums in the fine topology, or a choice of individual root branches.

\subsection{Provenance of this report}

This is a merged report. Its five sources are the research archives listed in \cref{finite:sec:provenance}: \emph{Finite Hahn Deformations in Surcomplex Analysis} (05), \emph{Finite Infinitesimal Fibres in Surcomplex Analysis} (06), \emph{Finite Maps and Residue Duality in Hahn-Coherent Surcomplex Analysis} (08, the base), \emph{Finite Intersections in Surcomplex Analysis} (11), and \emph{Hartogs Coherence, Finite Maps, and Residue Duality over the Surcomplex Numbers} (15), all dated 21 September 2026. All five extend one supplied manuscript \cite{finite:Supplied}, taking from it, and not reproving: the definition of Hahn-coherent functions as well-ordered families of ordinary holomorphic coefficient functions on one common domain; strong summability and its distinction from fine-topological convergence; standard-part thickenings and infinitesimal evaluation; the coefficientwise contour interpretation; one-variable coherent Weierstrass preparation, exact root lifting, and moving-pole residues; and the Hahn--Neumann support lemma. The one-variable results of that manuscript are not claimed as new here, and its closing research directions --- several-variable coherent division, analytic-set theory, and finite-mapping theorems with global support control --- are the target addressed.

The shared core of the five --- division with the $S+M$ certificate, conservation of $\mu$, characters as actual infinitesimal evaluations, the collision-stable residue pairing, and the trace--Jacobian identity --- amounted to roughly sixty to seventy per cent of each source and appears here once, in the base's formulation except where another source proved a strictly stronger statement. Each numbered result proved by only one source has been retained. Where the sources conflicted, they have been resolved rather than averaged: see \cref{finite:rem:eulerconflict} for the trace element, \cref{finite:sec:whichcertificate} for the two stability certificates, \eqref{finite:eq:twodiscriminants} for the two discriminants, and \cref{finite:sec:rings} for the four coefficient rings.

\section{Notation reference}
\label{finite:app:notation}

\begin{center}\small
\begin{longtable}{@{}L{0.26\textwidth}L{0.66\textwidth}@{}}
\toprule
Notation & Meaning\\
\midrule
\endhead
$\No$, $\SC=\No[i]$ & Surreal and surcomplex classes\\[0.25em]
$\Gamma$, $t^\gamma=\omega^{-\gamma}$ & Set-sized divisible exponent group; Conway normal-form monomial\\[0.25em]
$K=\C((t^\Gamma))$, $\Kc$, $\mm_K$ & Hahn workspace; its valuation ring (finite elements); its infinitesimals\\[0.25em]
$v$, $\st$ & Valuation ($v(0)=+\infty$); standard part, equivalently the exponent-zero coefficient\\[0.25em]
$\HH_\Gamma(U)$, $\Hp_\Gamma(U)$, $\Ip_\Gamma(U)$ & Fixed-polydisk ring $\OO(U)((t^\Gamma))$; its nonnegative-support subring; its positive-support ideal\\[0.25em]
$\Af_{n,\Gamma}=\C[[z]]((t^\Gamma))$ & Formal-coefficient ring (\cref{finite:sec:formal} only)\\[0.25em]
$U\sh$, $G\sh$ & Infinitesimal thickening; coherent evaluation\\[0.25em]
$f$, $F=f+E$, $\mu$ & Leading ordinary complete intersection; its Hahn perturbation; the leading colength\\[0.25em]
$E$, $M=E^*$ & Union of the positive error supports; the monoid it generates, $M+M=M$\\[0.25em]
$e_1=1,\ldots,e_\mu$; $\varepsilon_1,\ldots,\varepsilon_n$ & Fixed polynomial basis of $A_0$; Koszul exterior generators\\[0.25em]
$\bx$ & Box index set $\{0\leq\alpha_i<m_i\}$ in the coordinate-power case\\[0.25em]
$d$, $\delta$, $\partial=d+\delta$ & Unperturbed, perturbing and perturbed Koszul differentials\\[0.25em]
$(i,p,h)$; $S=(1+\delta h)^{-1}$; $p_F$, $h_F$ & Contraction; strong resolvent; perturbed projection and homotopy\\[0.25em]
$\NF_F$, $c_a(G)$, $M_G$ & Normal form; its coordinates; multiplication matrix in the fixed basis\\[0.25em]
$\Alg^+$, $\Alg$, $\Alg_\zeta$ & Integral and full zero algebras; the local Artinian factor at $\zeta$\\[0.25em]
$\mu_\zeta$, $B_\zeta$ & Local multiplicity and formal local algebra at the \emph{actual} zero $\zeta$\\[0.25em]
$J_F$, $\Res_F$, $\Res_{\zeta,F}$ & Jacobian determinant; cluster residue; localized residue\\[0.25em]
$\Gram_F$, $e^1,\ldots,e^\mu$, $\Eul_F$ & Residue Gram matrix; residue-dual basis; trace element $\sum_ae_ae^a$\\[0.25em]
$\TGram_F$, $\disc_F$, $\Delta_F$ & Trace Gram matrix; $\det\TGram_F$ (basis relative); $\det M_{J_F}$ (basis independent)\\[0.25em]
$A=DF\sh(a)$, $\kappa$ & Jacobian matrix at a simple zero; inverse-Jacobian loss $\max\{0,-v(A^{-1})\}$\\
\bottomrule
\end{longtable}
\end{center}

The dependency chain is: support calculus $\to$ operator inversion $\to$ contraction and conjugacy $\to$ finite free division $\to$ characters as evaluations $\to$ finite-jet local comparison $\to$ conservation of multiplicity and finite fibres $\to$ residue functional $\to$ perfect pairing $\to$ finite-witness transfer $\to$ trace--Jacobian, discriminants and stability. The parameter branch replaces the abstract contraction by explicit coordinate operators or by matrix elimination; the Hartogs branch is independent of all of it.

%======================================================================
\begingroup
\small
\setlength{\parskip}{0pt}
\begin{thebibliography}{99}
\setlength{\itemsep}{0.4em}

\bibitem{finite:Supplied}
\emph{Surcomplex Analysis: Infinitesimal Calculus, Hahn-Coherent Holomorphy, and Contour Theory}.
Unpublished manuscript supplied with the research request, dated September 21, 2026
(distributed as \texttt{surcomplex\_analysis(2).tex} and \texttt{surcomplex\_analysis(3).tex}).
In particular its sections on Hahn coherence, strong summability, Weierstrass preparation,
moving-pole residues, and the subsection ``Further mathematical targets suggested by the results.''

\bibitem{finite:Neumann}
B. H. Neumann, \emph{On ordered division rings},
Transactions of the American Mathematical Society \textbf{66} (1949), 202--252.
\url{https://doi.org/10.1090/S0002-9947-1949-0032593-5}.

\bibitem{finite:Poonen}
B. Poonen, \emph{Maximally complete fields},
L'Enseignement Math\'ematique (2) \textbf{39} (1993), 87--106.
Author's manuscript: \url{https://math.mit.edu/~poonen/papers/amsval.pdf}.

\bibitem{finite:Alling}
N. L. Alling, \emph{Foundations of Analysis over Surreal Number Fields},
North-Holland Mathematics Studies, vol.~141, North-Holland, Amsterdam, 1987.

\bibitem{finite:vdDE}
L. van den Dries and P. Ehrlich, \emph{Fields of surreal numbers and exponentiation},
Fundamenta Mathematicae \textbf{167} (2001), 173--188;
erratum, \textbf{168} (2001), 295--297.
\url{https://doi.org/10.4064/fm167-2-3}.

\bibitem{finite:BM}
A. Berarducci and V. Mantova,
\emph{Transseries as germs of surreal functions},
Transactions of the American Mathematical Society \textbf{371} (2019), 3549--3592.
\url{https://arxiv.org/abs/1703.01995}.

\bibitem{finite:CostinEhrlich}
O. Costin and P. Ehrlich, \emph{Integration on the surreals},
Advances in Mathematics \textbf{452} (2024), article 109823.
\url{https://arxiv.org/abs/2208.14331}.

\bibitem{finite:MullerStrohmaier}
J. M\"uller and A. Strohmaier,
\emph{The theory of Hahn meromorphic functions, a holomorphic Fredholm theorem, and its applications},
Analysis \& PDE \textbf{7} (2014), 745--770.
\url{https://doi.org/10.2140/apde.2014.7.745}.

\bibitem{finite:CluckersLipshitz}
R. Cluckers and L. Lipshitz, \emph{Fields with analytic structure},
Journal of the European Mathematical Society \textbf{13} (2011), 1147--1223.
\url{https://doi.org/10.4171/JEMS/278}.

\bibitem{finite:JoswigSmith}
Michael Joswig and Ben Smith,
``Convergent Hahn Series and Tropical Geometry of Higher Rank,''
\href{https://arxiv.org/abs/1809.01457}{arXiv:1809.01457}, version 4,
January 16, 2023.
Cited for the distinct existing higher-rank Hahn-geometric setting.

\bibitem{finite:Crainic}
M. Crainic, \emph{On the perturbation lemma, and deformations},
preprint (2004), arXiv:math/0403266.
\url{https://arxiv.org/abs/math/0403266}.

\bibitem{finite:Cartan}
H. Cartan, \emph{Faisceaux analytiques sur les vari\'et\'es de Stein:
d\'emonstration des th\'eor\`emes fondamentaux},
S\'eminaire Henri Cartan \textbf{4} (1951--1952), 1--15.
\url{https://eudml.org/doc/112267}.

\bibitem{finite:GunningRossi}
R. C. Gunning and H. Rossi,
\emph{Analytic Functions of Several Complex Variables},
Prentice-Hall, Englewood Cliffs, NJ, 1965;
reprinted, AMS Chelsea Publishing, vol.~368, 2009.

\bibitem{finite:Demailly}
J.-P. Demailly, \emph{Complex Analytic and Differential Geometry},
Universit\'e de Grenoble~I, open-content book, 2012.
Chapters II and IX.

\bibitem{finite:Lebl}
J. Lebl, \emph{Tasty Bits of Several Complex Variables},
open-content textbook, version 4.x.
Theorems 4.3.1 and 6.2.5.
\url{https://www.jirka.org/scv/scv.pdf}.

\bibitem{finite:BK}
Jacek Bochnak and Wojciech Kucharz,
``Global variants of Hartogs' theorem,''
\emph{Archiv der Mathematik} \textbf{113} (2019), 281--290.
\href{https://doi.org/10.1007/s00013-019-01340-7}{doi:10.1007/s00013-019-01340-7}.
The classical separate-holomorphy theorem recalled there is the background
input in \cref{finite:thm:hartogs}.

\bibitem{finite:GriffithsHarris}
P. Griffiths and J. Harris, \emph{Residues and zero-cycles on algebraic varieties},
Annals of Mathematics (2) \textbf{108} (1978), 461--505, especially \S I(a).
\url{https://publications.ias.edu/node/219}.

\bibitem{finite:Hartshorne}
R. Hartshorne, \emph{Residues and Duality},
Lecture Notes in Mathematics, vol.~20, Springer, Berlin, 1966.
\url{https://doi.org/10.1007/BFb0080482}.

\bibitem{finite:Kunz}
E. Kunz, \emph{Residues and Duality for Projective Algebraic Varieties},
University Lecture Series, vol.~47, American Mathematical Society, Providence, RI, 2008.
\url{https://doi.org/10.1090/ulect/047}.

\bibitem{finite:CattaniDickenstein}
E. Cattani and A. Dickenstein, \emph{Introduction to residues and resultants},
in: Solving Polynomial Equations, Algorithms and Computation in Mathematics, vol.~14,
Springer, Berlin, 2005, pp.~1--61.
\url{https://doi.org/10.1007/3-540-27357-3_1}.

\bibitem{finite:CDS}
E. Cattani, A. Dickenstein and B. Sturmfels,
\emph{Computing multidimensional residues},
in: Algorithms in Algebraic Geometry and Applications, Progress in Mathematics, vol.~143,
Birkh\"auser, Basel, 1996, pp.~135--164.
\url{https://arxiv.org/abs/alg-geom/9404011}.

\bibitem{finite:Cox}
David A. Cox,
\emph{Solving Equations via Algebras},
CIMPA Graduate School lecture notes, Buenos Aires, 2003.
\href{https://mate.dm.uba.ar/~visita16/cimpa/notes/Cox.pdf}{Author's course notes}.
An expanded chapter appears in \emph{Solving Polynomial Equations}, Springer, 2005.

\bibitem{finite:NayakSastry}
S. Nayak and P. Sastry,
\emph{Grothendieck duality and transitivity II: traces and residues via Verdier's isomorphism},
preprint (2019), arXiv:1903.01783.
\url{https://arxiv.org/abs/1903.01783}.

\bibitem{finite:StacksKoszul}
The Stacks Project Authors, \emph{The Stacks Project}.
``Koszul regular sequences,'' Section 15.31, especially Lemma 15.31.2 and the
distinctions preceding it,
\href{https://stacks.math.columbia.edu/tag/062D}{Tag 062D};
``Cohen--Macaulay rings,'' Section 10.104, especially Lemma 10.104.2,
\href{https://stacks.math.columbia.edu/tag/00N7}{Tag 00N7};
``Regular local rings,'' Section 10.106, especially Lemma 10.106.3,
\href{https://stacks.math.columbia.edu/tag/00NN}{Tag 00NN}.
Accessed September 21, 2026.

\end{thebibliography}
\endgroup
\end{document}
